% Les villes du Cameroun — Géographie Tle E — Cours signé OnBuch+
% Programme officiel MINESEC. Compiler avec : tectonic N03-villes-cameroun.tex
\newcommand{\DOCMATIERE}{Géographie}
\newcommand{\DOCNIVEAU}{Tle E}
\newcommand{\DOCMODULE}{Module 1 : Le Cameroun}
\newcommand{\DOCLECON}{N03}
\newcommand{\DOCTITRE}{Les villes du Cameroun}
% OnBuch+ — préambule « cours signés » ENRICHI (compilé avec tectonic / XeLaTeX)
% Système visuel « Pop » d'OnBuch+ : fond crème (jamais blanc), bordures encre
% épaisses, ombres « dures » décalées, titres Archivo Black en capitales.
% Version MATHÉMATIQUES : graphes (pgfplots/tikz), boîtes pédagogiques
% (définition, propriété, méthode, attention, le savais-tu, exercice résolu,
% exercice, TP, bilan), page de garde et signature OnBuch+ ancrée en bas.
%
% Macros attendues dans main.tex AVANT \input{preamble} :
%   \DOCMATIERE  \DOCNIVEAU  \DOCMODULE  \DOCLECON (numéro)  \DOCTITRE
\documentclass[11pt]{article}

\usepackage{fontspec}
\usepackage[french]{babel}
\usepackage[a4paper,margin=2cm,top=3.4cm,bottom=2.8cm,headheight=1.4cm,footskip=1.3cm]{geometry}
\usepackage{amsmath,amssymb,amsfonts}
\usepackage{mathtools} % \xmapsto, \coloneqq... souvent utilisés par les modèles
\usepackage{cancel} % simplifications barrées (bilans de forces)
% \abs{z} (module, valeur absolue) et \norm{u} (norme), utilisés par les modèles
\providecommand{\abs}{}\let\abs\relax\DeclarePairedDelimiter{\abs}{\lvert}{\rvert}
\providecommand{\norm}{}\let\norm\relax\DeclarePairedDelimiter{\norm}{\lVert}{\rVert}
\usepackage[table]{xcolor} % [table] : fournit aussi \rowcolors (alternance de lignes)
\usepackage{pagecolor}
\usepackage{tikz}
\usetikzlibrary{arrows.meta,positioning,calc,shapes.geometric,shapes.misc,shapes.arrows,decorations.pathmorphing,decorations.pathreplacing,decorations.markings,patterns,shadows,mindmap,trees,fit,backgrounds,matrix,intersections,angles,quotes}
\usepackage{pgfplots}
\usepackage[version=4]{mhchem} % équations nucléaires \ce{^{235}_{92}U}
\usepackage[european,siunitx]{circuitikz} % schémas électriques
\pgfplotsset{compat=1.18}
\usepgfplotslibrary{fillbetween,groupplots} % aires entre courbes (calcul intégral)
\usepackage{enumitem}
\usepackage{fancyhdr}
% [most] charge toutes les bibliothèques tcolorbox (listings, minted, raster,
% documentation...) et gonfle inutilement la mémoire TeX sur un document long
% (~300 boîtes) : on ne charge que ce qui est réellement utilisé.
\usepackage[skins,breakable]{tcolorbox}
\usepackage{graphicx}
\graphicspath{{/home/user/t-t/geo-onbuch/images/}}
\usepackage{colortbl}
\usepackage{booktabs}
\usepackage{tabularx}
\usepackage{multirow}
\usepackage{diagbox} % case d'en-tête diagonale des tableaux à double entrée
% Colonnes extensibles centrée / gauche / droite (C, L, R), souvent utilisées par les modèles
\newcolumntype{C}{>{\centering\arraybackslash}X}
\newcolumntype{L}{>{\raggedright\arraybackslash}X}
\newcolumntype{R}{>{\raggedleft\arraybackslash}X}
\usepackage{array}
\usepackage{multicol}
\usepackage{siunitx}
% parse-numbers=false : accepte aussi \SI{2,5 \times 10^{-3}}{...}
\sisetup{locale=FR,output-decimal-marker={,},group-minimum-digits=4,parse-numbers=false}
\DeclareSIUnit\ohm{\ensuremath{\symup{\Omega}}} % U+2126 absent de la police
\DeclareSIUnit\division{div} % réglages d'oscilloscope (ms/div, V/div)
\DeclareSIUnit\tr{tr} % tours (vitesse de rotation, ex. \SI{1500}{\tr\per\minute})
\DeclareSIUnit\molar{M} % concentration molaire (ex. \SI{3}{\molar}, \SI{1}{\micro\molar})
\DeclareSIUnit\an{an} % année (le modèle confond parfois avec \year, un compteur TeX) — ex. \SI{5}{\kilo\meter\per\an}
\DeclareSIUnit\FCFA{FCFA} % franc CFA (ex. \SI{500}{\FCFA\per\kilo\gram})
\DeclareSIUnit\degreeBrix{\textdegree Brix} % teneur en sucre d'un sirop/jus (réfractométrie)
\DeclareSIUnit\month{mois} % le modèle utilise parfois \month, un compteur TeX, comme unité de durée
\DeclareSIUnit\microliter{\micro\liter} % le modèle écrit parfois \microliter en un seul token
\DeclareSIUnit\million{millions} % dénombrement (ex. \SI{15}{\million\per\mL} spermatozoïdes)
\usepackage{qrcode}
% \degree : commande fréquemment utilisée par les modèles pour un angle ou une
% température en dehors de \SI{}{} (non fournie par les paquets chargés).
\providecommand{\degree}{^{\circ}}
% \qed (fin de démonstration), utilisé par les modèles sans amsthm
\providecommand{\qed}{\hfill$\square$}
% \barre : double barre des valeurs interdites dans un tableau de variations (tkz-tab)
\providecommand{\barre}{\ensuremath{\|}}
% environnement proof (amsthm non chargé) utilisé par les modèles
\ifdefined\proof\else\newenvironment{proof}[1][Démonstration]{\par\noindent\textit{#1.}\ }{\qed\par}\fi
\usepackage{lastpage}
\usepackage{hyperref}
\hypersetup{hidelinks}

% ============================================================
% Palette « Pop » OnBuch+ — mêmes hex que la classe P (plus_ui.dart)
% ============================================================
\definecolor{poporange}{HTML}{F59321}
\definecolor{popdark}{HTML}{E07A0C}
\definecolor{popcream}{HTML}{FFF6E8}
\definecolor{popink}{HTML}{1C1714}
\definecolor{poppurple}{HTML}{7A5AE0}
\definecolor{popgreen}{HTML}{2E9E6A}
\definecolor{poppink}{HTML}{E0457B}
\definecolor{popblue}{HTML}{2F7DE1}
\definecolor{popgold}{HTML}{C98A12}
\definecolor{popmuted}{HTML}{8A7F76}
\definecolor{popline}{HTML}{EADFD2}
\definecolor{popgray}{HTML}{8A7F76}
\colorlet{popgrey}{popgray}
% Alias tolérants (le modèle nomme parfois les couleurs différemment)
\colorlet{popwhite}{white}
\colorlet{popblack}{popink}
\colorlet{popred}{poppink}
\colorlet{popyellow}{popgold}
% Teintes claires (fonds de tableaux/encadrés) — le modèle nomme parfois
% les couleurs avec un suffixe L (ex. popblueL) sans les avoir définies.
\colorlet{poporangeL}{poporange!20}
\colorlet{popdarkL}{popdark!20}
\colorlet{poppurpleL}{poppurple!20}
\colorlet{popgreenL}{popgreen!20}
\colorlet{poppinkL}{poppink!20}
\colorlet{popblueL}{popblue!20}
\colorlet{popgoldL}{popgold!20}
\colorlet{popredL}{popred!20}
\colorlet{popgrayL}{popgray!20}
% Teintes claires pour fonds de tableaux / zones
\colorlet{popblueL}{popblue!10!white}
\colorlet{poporangeL}{poporange!14!white}
\colorlet{popgreenL}{popgreen!12!white}
\colorlet{poppurpleL}{poppurple!12!white}
\colorlet{poppinkL}{poppink!12!white}
\colorlet{popgoldL}{popgold!14!white}

\pagecolor{popcream}

% ============================================================
% Typographie
% ============================================================
\defaultfontfeatures{Ligatures=TeX}
\setmainfont{PlusJakartaSans-Medium.ttf}[Path=../../../fonts/,BoldFont=PlusJakartaSans-Bold.ttf]
% Maths en sans-serif (Fira Math) avec chiffres et lettres droites de Plus
% Jakarta, pour que les formules se fondent dans le texte.
\usepackage[math-style=ISO,bold-style=ISO]{unicode-math}
\setmathfont{FiraMath-Regular.otf}[Path=../../../fonts/]
\setmathfont{PlusJakartaSans-Medium.ttf}[Path=../../../fonts/,range={up/num,up/{Latin,latin},\mathup}]
\usepackage{icomma} % virgule décimale sans espace en mode math
% Caractères mathématiques Unicode absents des polices de texte (Plus Jakarta,
% Archivo Black) : rendus par leur équivalent mathématique, en texte comme en
% formule — sinon ils s'affichent en carré vide (ex. « Arithmétique dans ℤ »).
\usepackage{newunicodechar}
\newunicodechar{ℝ}{\ensuremath{\mathbb{R}}}
\newunicodechar{ℤ}{\ensuremath{\mathbb{Z}}}
\newunicodechar{ℂ}{\ensuremath{\mathbb{C}}}
\newunicodechar{ℕ}{\ensuremath{\mathbb{N}}}
\newunicodechar{ℚ}{\ensuremath{\mathbb{Q}}}
\newunicodechar{𝔻}{\ensuremath{\mathbb{D}}}
\newunicodechar{α}{\ensuremath{\alpha}}
\newunicodechar{β}{\ensuremath{\beta}}
\newunicodechar{γ}{\ensuremath{\gamma}}
\newunicodechar{δ}{\ensuremath{\delta}}
\newunicodechar{ε}{\ensuremath{\varepsilon}}
\newunicodechar{θ}{\ensuremath{\theta}}
\newunicodechar{λ}{\ensuremath{\lambda}}
\newunicodechar{μ}{\ensuremath{\mu}}
\newunicodechar{σ}{\ensuremath{\sigma}}
\newunicodechar{τ}{\ensuremath{\tau}}
\newunicodechar{φ}{\ensuremath{\varphi}}
\newunicodechar{ψ}{\ensuremath{\psi}}
\newunicodechar{ω}{\ensuremath{\omega}}
\newunicodechar{Σ}{\ensuremath{\Sigma}}
% majuscules produites par \MakeUppercase dans les titres (α-aminés -> Α)
\newunicodechar{Α}{\ensuremath{\alpha}}
\newunicodechar{Β}{\ensuremath{\beta}}
\newunicodechar{Γ}{\ensuremath{\Gamma}}
\newunicodechar{Δ}{\ensuremath{\Delta}}
\newunicodechar{Ε}{\ensuremath{\varepsilon}}
\newunicodechar{Θ}{\ensuremath{\Theta}}
\newunicodechar{Λ}{\ensuremath{\Lambda}}
\newunicodechar{Μ}{\ensuremath{\mu}}
\newunicodechar{Τ}{\ensuremath{\tau}}
\newunicodechar{Φ}{\ensuremath{\Phi}}
\newunicodechar{Ψ}{\ensuremath{\Psi}}
\newunicodechar{Ω}{\ensuremath{\Omega}}
\newunicodechar{↦}{\ensuremath{\mapsto}}
\newunicodechar{∈}{\ensuremath{\in}}
\newunicodechar{∉}{\ensuremath{\notin}}
\newunicodechar{∧}{\ensuremath{\wedge}}
\newunicodechar{⇔}{\ensuremath{\Leftrightarrow}}
\newunicodechar{⟺}{\ensuremath{\Longleftrightarrow}}
\newunicodechar{⇒}{\ensuremath{\Rightarrow}}
\newunicodechar{⟹}{\ensuremath{\Longrightarrow}}
\newunicodechar{∘}{\ensuremath{\circ}}
\newunicodechar{∪}{\ensuremath{\cup}}
\newunicodechar{∩}{\ensuremath{\cap}}
\newunicodechar{≡}{\ensuremath{\equiv}}
\newunicodechar{⊂}{\ensuremath{\subset}}
\newunicodechar{⊥}{\ensuremath{\perp}}
\newunicodechar{∃}{\ensuremath{\exists}}
\newunicodechar{∀}{\ensuremath{\forall}}
\newunicodechar{∛}{\ensuremath{\sqrt[3]{\,}}}
\newunicodechar{∜}{\ensuremath{\sqrt[4]{\,}}}
\newunicodechar{⃗}{\ensuremath{{}^{\to}}}
\newunicodechar{ⁿ}{\ifmmode^{n}\else\textsuperscript{\ensuremath{n}}\fi}
\newunicodechar{ˣ}{\ifmmode^{x}\else\textsuperscript{\ensuremath{x}}\fi}
\newunicodechar{ᵇ}{\ifmmode^{b}\else\textsuperscript{\ensuremath{b}}\fi}
\newunicodechar{ʳ}{\ifmmode^{r}\else\textsuperscript{\ensuremath{r}}\fi}
\newunicodechar{ᵘ}{\ifmmode^{u}\else\textsuperscript{\ensuremath{u}}\fi}
\newunicodechar{ʸ}{\ifmmode^{y}\else\textsuperscript{\ensuremath{y}}\fi}
\newunicodechar{ᵃ}{\ifmmode^{a}\else\textsuperscript{\ensuremath{a}}\fi}
\newunicodechar{ᵉ}{\ifmmode^{e}\else\textsuperscript{\ensuremath{e}}\fi}
\newunicodechar{ᵏ}{\ifmmode^{k}\else\textsuperscript{\ensuremath{k}}\fi}
\newunicodechar{ᵖ}{\ifmmode^{p}\else\textsuperscript{\ensuremath{p}}\fi}
\newunicodechar{ᴺ}{\ifmmode^{N}\else\textsuperscript{\ensuremath{N}}\fi}
\newunicodechar{ᶜ}{\ifmmode^{c}\else\textsuperscript{\ensuremath{c}}\fi}
\newunicodechar{ʰ}{\ifmmode^{h}\else\textsuperscript{\ensuremath{h}}\fi}
\newunicodechar{ᵠ}{\ifmmode^{q}\else\textsuperscript{\ensuremath{q}}\fi}
\newunicodechar{ₙ}{\ifmmode_{n}\else\textsubscript{\ensuremath{n}}\fi}
\newunicodechar{ₐ}{\ifmmode_{a}\else\textsubscript{\ensuremath{a}}\fi}
\newunicodechar{ᵢ}{\ifmmode_{i}\else\textsubscript{\ensuremath{i}}\fi}
\newunicodechar{ᵣ}{\ifmmode_{r}\else\textsubscript{\ensuremath{r}}\fi}
\newunicodechar{ₖ}{\ifmmode_{k}\else\textsubscript{\ensuremath{k}}\fi}
\newunicodechar{ᵦ}{\ifmmode_{\beta}\else\textsubscript{\ensuremath{\beta}}\fi}
\newunicodechar{ₓ}{\ifmmode_{x}\else\textsubscript{\ensuremath{x}}\fi}
\newfontfamily\popdisplay{ArchivoBlack-Regular.ttf}[Path=../../../fonts/]
\newfontfamily\poplogo{SpaceGrotesk-Bold.ttf}[Path=../../../fonts/]
\newfontfamily\popsemi{PlusJakartaSans-SemiBold.ttf}[Path=../../../fonts/]
\newfontfamily\popxbold{PlusJakartaSans-ExtraBold.ttf}[Path=../../../fonts/]
\newfontfamily\popxboldls{PlusJakartaSans-ExtraBold.ttf}[Path=../../../fonts/,LetterSpace=3.0]

\setlist[enumerate]{leftmargin=1.7em,itemsep=5pt,topsep=4pt}
\setlist[itemize]{leftmargin=1.7em,itemsep=4pt,topsep=4pt,label={\color{poporange}\textbullet}}
\setlength{\parindent}{0pt}
\setlength{\parskip}{5pt}
\renewcommand{\arraystretch}{1.25}

% pgfplots : style Pop par défaut
\pgfplotsset{
  every axis/.append style={axis line style={popink,line width=1pt},tick style={popink},
    grid style={popline,line width=0.5pt},label style={font=\small},tick label style={font=\footnotesize}},
  popcurve/.style={poporange,line width=1.8pt,no markers,smooth},
}
% Même style disponible dans un \draw TikZ simple (hors pgfplots)
\tikzset{popcurve/.style={poporange,line width=1.8pt}}
\tikzset{popgrid/.style={popline,very thin}}
\colorlet{popcurve}{poporange} % le nom du style est parfois utilisé comme couleur

% ============================================================
% Pill / squiggle
% ============================================================
\newcommand{\poppill}[3]{%
  \tikz[baseline=-0.55ex]\node[draw=popink,line width=1.2pt,rounded corners=2.6pt,%
    fill=#2,inner xsep=6.5pt,inner ysep=3pt]{%
    \popxboldls\fontsize{7}{7}\selectfont\color{#3}\MakeUppercase{#1}};%
}
\newcommand{\popsquiggle}[1]{%
  \tikz[baseline=-0.4ex]{
    \draw[#1,line width=2.6pt,line cap=round]
      plot [smooth] coordinates {(0,0.09) (0.34,0.01) (0.68,0.21) (1.02,0.01) (1.36,0.10)};
  }%
}
% Mot-clé surligné (marqueur orange sous le texte)
\newcommand{\cle}[1]{{\popxbold\color{popdark}#1}}

% ============================================================
% En-tête / pied
% ============================================================
\pagestyle{fancy}
\fancyhf{}
\renewcommand{\headrulewidth}{0pt}
\renewcommand{\footrulewidth}{0pt}
\lhead{\raisebox{-0.15ex}{\poplogo\fontsize{13}{13}\selectfont\color{popink}OnBuch\color{poporange}+}}
\rhead{\poppill{\DOCMATIERE\ \textbullet\ \DOCNIVEAU\ \textbullet\ Leçon \DOCLECON}{white}{popink}}
\lfoot{\popxbold\fontsize{7.6}{7.6}\selectfont\color{popmuted}\MakeUppercase{Cours signé OnBuch+}\ \textbullet\ {\popsemi\DOCTITRE}}
\rfoot{\tikz[baseline=-0.6ex]\node[rounded corners=8.5pt,fill=popink,minimum height=17pt,minimum width=17pt,inner xsep=5pt,inner sep=0]{\popxbold\fontsize{8}{8}\selectfont\color{popcream}\thepage\,/\,\pageref*{LastPage}};}

% ============================================================
% Titres
% ============================================================
\newcommand{\chaptitre}[2]{%
  \begin{tcolorbox}[enhanced,colback=poporange,colframe=popink,boxrule=2.6pt,arc=11pt,
      drop shadow=popink,left=15pt,right=15pt,top=12pt,bottom=11pt,before skip=3pt,after skip=18pt]
    \poppill{#2}{popink}{popcream}\par\vspace{9pt}
    {\raggedright\hyphenpenalty=10000\exhyphenpenalty=10000\popdisplay\fontsize{22}{24}\selectfont\color{popcream}\MakeUppercase{#1}\par}
  \end{tcolorbox}
}
% Section numérotée : pastille orange avec le numéro + titre Archivo
\newcounter{popsec}
\newcommand{\coursec}[1]{%
  \stepcounter{popsec}\setcounter{popsub}{0}%
  \par\vspace{16pt}\noindent
  \begin{minipage}{\linewidth}
  \tikz[baseline=-0.7ex]\node[draw=popink,line width=1.6pt,fill=poporange,rounded corners=4pt,minimum size=22pt,inner sep=0]{\popdisplay\fontsize{12}{12}\selectfont\color{popcream}\thepopsec};%
  \hspace{8pt}{\popdisplay\fontsize{15}{17}\selectfont\color{popink}\MakeUppercase{#1}}\par
  \vspace{4pt}\noindent{\color{poporange}\rule{36pt}{3.6pt}}
  \end{minipage}\par\nopagebreak\vspace{8pt}
}
\newcounter{popsub}
\newcommand{\courssub}[1]{%
  \stepcounter{popsub}%
  \par\vspace{9pt}\noindent{\popxbold\fontsize{11.5}{13}\selectfont\color{popink}{\color{poporange}\thepopsec.\thepopsub}\hspace{6pt}#1}\par\nopagebreak\vspace{4pt}
}

% ============================================================
% Boîtes pédagogiques — même recette : fond blanc, bordure épaisse colorée,
% ombre dure encre, badge plein. Titre optionnel [#1].
% ============================================================
\tcbset{popbase/.style={breakable,enhanced,colback=white,boxrule=2.3pt,arc=9pt,
  shadow={3pt}{-3pt}{0pt}{popink},left=14pt,right=14pt,top=11pt,bottom=11pt,before skip=11pt,after skip=15pt,
  fontupper=\popsemi\fontsize{10.6}{14.5}\selectfont\color{popink},
  fontlower=\popsemi\fontsize{10.6}{14.5}\selectfont\color{popink}}}
% Coche dessinée (le glyphe ✓ n'existe pas dans les polices du document)
\newcommand{\popcheck}[1]{\tikz[baseline=-0.5ex]\draw[#1,line width=1.6pt,line cap=round,line join=round](-0.13,0.0)--(-0.04,-0.09)--(0.13,0.1);}
\providecommand{\coche}{\popcheck{popgreen}}
\providecommand{\resultat}[1]{\par\smallskip\noindent\fcolorbox{popgreen}{popgreenL}{\parbox{\dimexpr\linewidth-2\fboxsep-2\fboxrule\relax}{\textbf{Résultat.} #1}}\par\smallskip}
\newcommand{\popboxhead}[3]{\poppill{#1}{#2}{white}\ifx\relax#3\relax\else\hspace{6pt}{\popxbold\fontsize{10.6}{12}\selectfont\color{popink}#3}\fi\par\vspace{7pt}}

\newtcolorbox{aretenir}[1][]{popbase,colframe=popgreen,before upper={\popboxhead{À retenir}{popgreen}{#1}}}
\newtcolorbox{exemplebox}[1][]{popbase,colframe=poporange,before upper={\popboxhead{Exemple}{poporange}{#1}}}
\newtcolorbox{definition}[1][]{popbase,colframe=popblue,before upper={\popboxhead{Définition}{popblue}{#1}}}
\newtcolorbox{propriete}[1][]{popbase,colframe=poppurple,before upper={\popboxhead{Propriété}{poppurple}{#1}}}
\newtcolorbox{methode}[1][]{popbase,colframe=popgold,before upper={\popboxhead{Méthode}{popgold}{#1}}}
\newtcolorbox{attention}[1][]{popbase,colframe=poppink,before upper={\popboxhead{Attention}{poppink}{#1}}}
\newtcolorbox{experience}[1][]{popbase,colframe=popblue,colback=popblueL,before upper={\popboxhead{Expérience}{popblue}{#1}}}
% « Le savais-tu ? » : carte violette pleine (contexte, culture, Cameroun)
\newtcolorbox{savaistu}[1][]{popbase,colback=poppurple,colframe=popink,
  fontupper=\popsemi\fontsize{10.4}{14.2}\selectfont\color{white},
  before upper={\poppill{Le savais-tu\,?}{white}{poppurple}\ifx\relax#1\relax\else\hspace{6pt}{\popxbold\color{white}#1}\fi\par\vspace{7pt}}}
% Exercice résolu : énoncé en haut, \tcblower puis la solution sur fond crème
\newcounter{popexo}\newcounter{popexr}
\newtcolorbox{exoresolu}[1][]{popbase,colframe=poporange,colbacklower=popcream,
  segmentation style={popink,line width=1.2pt,dashed},
  before upper={\stepcounter{popexr}\popboxhead{Exercice résolu \thepopexr}{poporange}{#1}},
  before lower={\poppill{Solution}{popink}{popcream}\par\vspace{6pt}}}
% Exercice (non résolu en place) — la correction est dans la partie Corrigés
\newtcolorbox{exercice}[1][]{popbase,colframe=popink,
  before upper={\stepcounter{popexo}\popboxhead{Exercice \thepopexo}{popink}{#1}}}
% Corrigé
\newtcolorbox{corrige}[1][]{popbase,colframe=popgreen,colback=popgreenL,
  before upper={\popboxhead{Corrigé}{popgreen}{#1}}}
% Objectifs / prérequis / bilan
\newtcolorbox{objectifs}{popbase,colframe=popink,colback=poporangeL,
  before upper={\popboxhead{Objectifs de la leçon}{popink}{}}}
\newtcolorbox{prerequis}{popbase,colframe=popink,colback=popblueL,
  before upper={\popboxhead{Prérequis}{popblue}{}}}
\newtcolorbox{bilan}{popbase,colframe=popink,colback=white,boxrule=2.8pt,
  before upper={\popboxhead{Fiche bilan}{poporange}{L'essentiel à connaître pour l'examen}}}
% Figure encadrée avec légende
\newtcolorbox{popfigure}[1][]{enhanced,breakable=false,colback=white,colframe=popline,boxrule=1.4pt,arc=8pt,
  left=10pt,right=10pt,top=10pt,bottom=8pt,before skip=10pt,after skip=12pt,halign=center,
  lower separated=false,halign lower=center,
  fontlower=\popsemi\fontsize{9}{11.5}\selectfont\color{popmuted},
  #1}
\newcommand{\legende}[1]{\par\vspace{6pt}{\popsemi\fontsize{9}{11.5}\selectfont\color{popmuted}#1}}

% ============================================================
% Signature OnBuch+
% ============================================================
\newcommand{\poplogoinline}[1]{{\poplogo\fontsize{#1}{#1}\selectfont\color{popink}OnBuch\color{poporange}+}}
\newcommand{\signatureonbuch}{%
  \begin{tcolorbox}[enhanced,colback=popink,colframe=popink,boxrule=2.2pt,arc=10pt,
      drop shadow=poporange,left=14pt,right=14pt,top=10pt,bottom=10pt,before skip=0pt,after skip=0pt]
    \tikz[baseline=-0.7ex]\node[circle,fill=popgreen,draw=popcream,line width=1.3pt,minimum size=22pt,inner sep=0]{\popcheck{white}};%
    \hspace{9pt}\begin{minipage}[c]{0.62\linewidth}
      {\popxbold\fontsize{12}{13}\selectfont\color{popcream}Signé\ \textbullet\ {\poplogo OnBuch\color{poporange}+}}\par\vspace{2pt}
      {\raggedright\popsemi\fontsize{8}{10.5}\selectfont\color{popline}\MakeUppercase{Équipe pédagogique OnBuch+}\\\MakeUppercase{Conforme au programme officiel MINESEC}\par}
    \end{minipage}\hfill
    {\poplogo\fontsize{22}{22}\selectfont\color{popcream}OnBuch\color{poporange}+}
  \end{tcolorbox}%
}

% ============================================================
% Page de garde
% #1 titre · #2 matière · #3 niveau · #4 module · #5 n° leçon · #6 durée ·
% #7 description / objectifs (texte ou itemize)
% ============================================================
\newcommand{\pagegarde}[7]{%
  \thispagestyle{empty}
  \newgeometry{a4paper,margin=1.7cm,top=1.6cm,bottom=1.4cm}%
  \begin{tikzpicture}[remember picture,overlay]
    \fill[poporange!18] ([xshift=-3.2cm,yshift=-3.6cm]current page.north east) circle (4.6cm);
    \fill[poppurple!14] ([xshift=2.4cm,yshift=7.6cm]current page.south west) circle (3.8cm);
  \end{tikzpicture}%
  {\poplogo\fontsize{30}{30}\selectfont\color{popink}OnBuch\color{poporange}+}\hfill
  \raisebox{0.6ex}{\poppill{Cours signé \textbullet\ Premium}{popink}{popcream}}\par
  \vspace{1pt}\popsquiggle{poppurple}\par
  \vspace{8pt}
  \begin{tcolorbox}[enhanced,colback=poporange,colframe=popink,boxrule=2.8pt,arc=13pt,
      drop shadow=popink,left=18pt,right=18pt,top=12pt,bottom=13pt,after skip=10pt]
    \poppill{#2 \textbullet\ #3}{popink}{popcream}\hspace{5pt}\poppill{#4}{white}{popink}\par\vspace{8pt}
    {\popdisplay\fontsize{14}{14}\selectfont\color{popink}LEÇON #5}\par\vspace{4pt}
    {\raggedright\hyphenpenalty=10000\exhyphenpenalty=10000\popdisplay\fontsize{25}{27}\selectfont\color{popcream}\MakeUppercase{#1}\par}
    \vspace{8pt}
    {\popxbold\fontsize{9.5}{9.5}\selectfont\color{popink}\MakeUppercase{Durée conseillée : #6}}
  \end{tcolorbox}
  \begin{tcolorbox}[enhanced,colback=white,colframe=popink,boxrule=2.2pt,arc=10pt,
      drop shadow=popink,left=16pt,right=16pt,top=10pt,bottom=10pt,after skip=8pt]
    \poppill{Dans cette leçon}{poppurple}{white}\par\vspace{5pt}
    \popsemi\fontsize{9.3}{12.6}\selectfont\color{popink}#7
  \end{tcolorbox}
  \noindent\begin{minipage}[t]{0.487\linewidth}
    \begin{tcolorbox}[enhanced,colback=popgreen,colframe=popink,boxrule=2.2pt,arc=10pt,
        drop shadow=popink,left=11pt,right=11pt,top=10pt,bottom=10pt,height=3.1cm,valign=center]
      \tcbox[colback=white,colframe=popink,boxrule=1.3pt,arc=5pt,boxsep=0pt,nobeforeafter,
        left=3pt,right=3pt,top=3pt,bottom=3pt,valign=center]{\qrcode[height=1.9cm]{https://play.google.com/store/apps/details?id=com.luvvix.onbuch&hl=fr}}%
      \hspace{8pt}\begin{minipage}[c]{0.5\linewidth}
        {\popdisplay\fontsize{11}{12.5}\selectfont\color{white}\MakeUppercase{Télécharge OnBuch}}\par\vspace{4pt}
        {\raggedright\popsemi\fontsize{8.4}{11}\selectfont\color{white}Gratuit \textbullet\ Google Play\\Tuteur IA, annales, concours, résultats\par}
      \end{minipage}
    \end{tcolorbox}
  \end{minipage}\hfill
  \begin{minipage}[t]{0.487\linewidth}
    \begin{tcolorbox}[enhanced,colback=poppurple,colframe=popink,boxrule=2.2pt,arc=10pt,
        drop shadow=popink,left=11pt,right=11pt,top=10pt,bottom=10pt,height=3.1cm,valign=center]
      \tikz[baseline=-0.7ex]\node[circle,fill=white,draw=popink,line width=1.3pt,minimum size=1.6cm,inner sep=0]{\popdisplay\fontsize{24}{24}\selectfont\color{poppurple}+};%
      \hspace{8pt}\begin{minipage}[c]{0.6\linewidth}
        {\popdisplay\fontsize{11}{12.5}\selectfont\color{white}\MakeUppercase{Passe à OnBuch+}}\par\vspace{4pt}
        {\raggedright\popsemi\fontsize{8.4}{11}\selectfont\color{white}Tous les cours signés, Tuteur IA illimité, fiches PDF — onglet OnBuch+ de l'app\par}
      \end{minipage}
    \end{tcolorbox}
  \end{minipage}
  \vspace{8pt}
  \popcontenu
  \vfill
  \signatureonbuch
  \restoregeometry
  \newpage
}

% Rangée « Ce cours contient » (page de garde)
\newcommand{\poptile}[3]{% #1 couleur #2 gros texte #3 légende
  \begin{tcolorbox}[enhanced,colback=#1,colframe=popink,boxrule=2pt,arc=9pt,shadow={2.5pt}{-2.5pt}{0pt}{popink},
      left=6pt,right=6pt,top=9pt,bottom=9pt,halign=center,valign=center,height=1.75cm,width=\linewidth,nobeforeafter]
    {\popdisplay\fontsize{12.5}{13}\selectfont\color{white}\MakeUppercase{#2}}\par\vspace{3pt}
    {\centering\popsemi\fontsize{7.8}{9.5}\selectfont\color{white}#3\par}
  \end{tcolorbox}}
\newcommand{\popcontenu}{%
  \par\noindent{\popxboldls\fontsize{7.5}{7.5}\selectfont\color{popmuted}CE COURS SIGNÉ CONTIENT}\par\vspace{7pt}
  \noindent\begin{minipage}[t]{0.235\linewidth}\poptile{popblue}{Cours}{complet, illustré, pas à pas}\end{minipage}\hfill
  \begin{minipage}[t]{0.235\linewidth}\poptile{poporange}{Méthodes}{et exercices résolus}\end{minipage}\hfill
  \begin{minipage}[t]{0.235\linewidth}\poptile{poppink}{Exercices}{type examen + corrigés}\end{minipage}\hfill
  \begin{minipage}[t]{0.235\linewidth}\poptile{popgreen}{Fiche bilan}{l'essentiel à retenir}\end{minipage}\par}

% Sommaire
\newtcolorbox{sommaire}{enhanced,colback=white,colframe=popink,boxrule=2.2pt,arc=10pt,
  drop shadow=popink,left=18pt,right=18pt,top=13pt,bottom=13pt,before skip=6pt,after skip=20pt,
  before upper={\poppill{Sommaire}{poppurple}{white}\par\vspace{9pt}\popsemi\fontsize{10.6}{14.5}\selectfont\color{popink}}}

% Signature de fin de cours — poussée tout en bas de la dernière page
\newcommand{\findecours}{%
  \par\vfill\nopagebreak
  \begin{center}\popsquiggle{poporange}\end{center}\vspace{4pt}
  \signatureonbuch
}
\AtBeginDocument{\def\llbracket{\mathopen{[\mkern-3.5mu[}}\def\rrbracket{\mathclose{]\mkern-3.5mu]}}} % intervalles d'entiers (glyphe ⟦ absent de la police)
% Glyphes absents de Fira Math : redéfinis à partir de glyphes disponibles
\AtBeginDocument{%
  \def\setminus{\mathbin{\backslash}}\let\smallsetminus\setminus
  \def\vdots{\mathord{\vbox{\baselineskip3.2pt\lineskiplimit0pt\kern4pt\hbox{.}\hbox{.}\hbox{.}}}}%
  \def\ddots{\mathinner{\mkern1mu\raise6.5pt\hbox{.}\mkern2mu\raise3.6pt\hbox{.}\mkern2mu\raise0.7pt\hbox{.}\mkern1mu}}%
  \def\triangle{\mathord{\scalebox{0.9}{$\symup{\Delta}$}}}%
  \def\nabla{\mathord{\raisebox{\depth}{\scalebox{1}[-1]{$\symup{\Delta}$}}}}%
  \def\checkmark{\popcheck{popdark}}%
  \def\star{\mathbin{\ast}}\def\bigstar{\mathord{\ast}}%
  \def\diamond{\mathbin{\scalebox{0.6}{$\Diamond$}}}%
  \def\bigcup{\mathop{\mathchoice{\raisebox{-0.3em}{\scalebox{1.7}{$\cup$}}}{\scalebox{1.25}{$\cup$}}{\cup}{\cup}}\displaylimits}%
  \def\bigcap{\mathop{\mathchoice{\raisebox{-0.3em}{\scalebox{1.7}{$\cap$}}}{\scalebox{1.25}{$\cap$}}{\cap}{\cap}}\displaylimits}%
  \def\lesseqqgtr{\mathrel{\lessgtr}}\def\bigoplus{\mathop{\oplus}\displaylimits}%
}
\providecommand{\ssi}{si et seulement si} % abréviation fréquente
\AtBeginDocument{\providecommand{\wideparen}{\overparen}} % arc de cercle

\begin{document}

\pagegarde{Les villes du Cameroun}{Géographie}{Tle E}{Module 1 : Le Cameroun}{N03}{2 h}{Cette leçon étudie les villes du Cameroun comme espaces d'intégration nationale et centres d'impulsion économique, tout en analysant les problèmes urbains qui freinent leur développement. À travers l'observation de cartes, de photographies et de données chiffrées, l'élève apprend à décrire, classer et représenter les phénomènes urbains camerounais.
\vspace{4pt}
\begin{multicols}{2}\small
\begin{itemize}
\item À la fin de cette leçon, je sais définir la ville et le quartier cosmopolite à partir d'exemples camerounais
\item À la fin de cette leçon, je sais observer, localiser et décrire les principales villes du Cameroun sur une carte
\item À la fin de cette leçon, je sais expliquer le rôle de la ville comme espace d'intégration des populations, des cultures et des religions
\item À la fin de cette leçon, je sais montrer que la ville est un centre d'impulsion économique et un lieu de polarisation des services
\item À la fin de cette leçon, je sais inventorier et classer les problèmes des villes camerounaises à partir de photographies et de documents
\item À la fin de cette leçon, je sais légender une carte de l'urbanisation du Cameroun
\end{itemize}\end{multicols}}

\begin{sommaire}
\begin{multicols}{2}
\begin{enumerate}
\item Généralités sur la ville et l'urbanisation au Cameroun
\item La ville, espace d'intégration nationale
\item La ville, centre d'impulsion économique et de polarisation des services
\item Les problèmes des villes camerounaises
\item Solutions et perspectives pour les villes camerounaises
\item Travaux pratiques : représenter l'urbanisation du Cameroun
\item Activité d'intégration
\item Méthodes et exercices résolus
\item Exercices
\item Corrigés des exercices
\item Fiche bilan
\end{enumerate}
\end{multicols}
\end{sommaire}

\begin{objectifs}
\begin{itemize}
\item À la fin de cette leçon, je sais définir la ville et le quartier cosmopolite à partir d'exemples camerounais
\item À la fin de cette leçon, je sais observer, localiser et décrire les principales villes du Cameroun sur une carte
\item À la fin de cette leçon, je sais expliquer le rôle de la ville comme espace d'intégration des populations, des cultures et des religions
\item À la fin de cette leçon, je sais montrer que la ville est un centre d'impulsion économique et un lieu de polarisation des services
\item À la fin de cette leçon, je sais inventorier et classer les problèmes des villes camerounaises à partir de photographies et de documents
\item À la fin de cette leçon, je sais légender une carte de l'urbanisation du Cameroun
\item À la fin de cette leçon, je sais construire un croquis d'un quartier cosmopolite et un diagramme de la population urbaine
\item À la fin de cette leçon, je sais interpréter des tableaux statistiques sur la croissance urbaine du Cameroun
\item À la fin de cette leçon, je sais proposer des solutions aux problèmes des villes camerounaises
\end{itemize}
\end{objectifs}

\begin{prerequis}
\begin{itemize}
\item Notion de population et de répartition spatiale (classe de Première)
\item Notion de migration et d'exode rural (classe de Première)
\item Localisation des dix régions et des chefs-lieux du Cameroun
\item Lecture d'une carte et construction d'une légende (classes antérieures)
\item Notion d'activités économiques (primaire, secondaire, tertiaire)
\end{itemize}
\end{prerequis}

\newpage

\coursec{Généralités sur la ville et l'urbanisation au Cameroun}

Chaque matin, des milliers de jeunes quittent les villages de l'Ouest ou de l'Extrême-Nord pour tenter leur chance à Douala ou Yaoundé, convaincus que la ville offre travail, éducation et modernité. Pourtant, une fois sur place, beaucoup découvrent des quartiers sans eau courante, des rues anarchiques et une insécurité croissante. Comment expliquer cette contradiction entre la ville rêvée et la ville vécue ? Quel rôle jouent réellement les villes camerounaises dans l'intégration nationale et le développement économique du pays ? Cette leçon nous invite à observer, décrire et cartographier les villes du Cameroun pour répondre à ces questions.

\courssub{Qu'est-ce qu'une ville ? Définition et critères}

\begin{definition}[La ville et l'urbanisation]
Au sens géographique, une \cle{ville} est un espace concentré de population, d'activités économiques (industrielles, commerciales, tertiaires) et de services (administration, santé, éducation, culture), par opposition au \cle{milieu rural} où dominent les activités primaires (agriculture, élevage, pêche) et où la densité de population est plus faible.

L'\cle{urbanisation} désigne le processus d'accroissement de la part de la population urbaine dans la population totale d'un pays. Elle résulte de trois mécanismes combinés : l'exode rural (migration campagne $\to$ ville), l'accroissement naturel de la population urbaine (excédent des naissances sur les décès) et la reclassification de localités rurales en communes urbaines du fait de leur croissance.
\end{definition}

Au Cameroun, l'Institut National de la Statistique (INS) et la loi n°2004/018 du 22 juillet 2004 fixent les critères officiels pour qu'une localité accède au statut de \cle{commune urbaine} (et donc de ville) :
\begin{itemize}
    \item \textbf{Seuil démographique :} une population résidente d'au moins \num{5000} habitants (seuil souvent dépassé dans la pratique pour les chefs-lieux).
    \item \textbf{Fonction administrative :} être chef-lieu de région, de département ou d'arrondissement.
    \item \textbf{Présence de services urbains :} existence d'équipements collectifs (marché structuré, centre de santé/hôpital, lycée/collège, réseau d'eau et d'électricité, voirie bitumée, commissariat/gendarmerie).
\end{itemize}

\begin{attention}[Erreur fréquente : confondre croissance urbaine et urbanisation]
Ne pas écrire : « L'urbanisation de Douala est de 5 \% par an. »
\begin{itemize}
    \item La \cle{croissance urbaine} (ou croissance de la population urbaine) est l'augmentation absolue ou relative du nombre d'habitants dans les villes (taux en \%/an).
    \item L'\cle{urbanisation} est la variation de la \textbf{part} de la population urbaine dans la population totale (taux d'urbanisation en \%). Un pays peut avoir une forte croissance urbaine mais une urbanisation lente si la population rurale croît aussi vite.
\end{itemize}
Au Bac, précise toujours : « Le taux d'urbanisation \textbf{augmente} » (part relative) et non « La population urbaine \textbf{augmente} » (effectif absolu).
\end{attention}

\courssub{L'urbanisation au Cameroun : évolution et ampleur}

\begin{definition}[Taux d'urbanisation]
Le \cle{taux d'urbanisation} (noté $T_u$) s'exprime en pourcentage :
\[
T_u = \frac{\text{Population urbaine}}{\text{Population totale}} \times 100
\]
Il mesure le poids démographique des villes dans le pays. Un taux supérieur à 50 \% signifie que la majorité de la population vit en ville (seuil de la « transition urbaine »).
\end{definition}

Le Cameroun a connu une transition urbaine spectaculaire depuis l'indépendance.
\begin{itemize}
    \item \textbf{1960 (indépendance) :} taux d'urbanisation $\approx$ \num{10} \%. La population était massivement rurale.
    \item \textbf{1976 (1er RGPH) :} $\approx$ \num{20} \%. Début de l'exode rural massif vers Douala et Yaoundé (planification, emplois publics, plantations).
    \item \textbf{1987 (2e RGPH) :} $\approx$ \num{35} \%. Crise économique des années 1980 $\to$ retour à la terre freiné, mais villes refuges.
    \item \textbf{2005 (3e RGPH) :} $\approx$ \num{48} \%. Quasi parité.
    \item \textbf{2020 (estimations INS/ONU) :} $\approx$ \num{57} \%. Le Cameroun est devenu un pays \textbf{majoritairement urbain}.
\end{itemize}

\begin{popfigure}
\begin{tikzpicture}
\begin{axis}[
    ybar,
    bar width=12pt,
    width=\linewidth,
    height=9cm,
    ymin=0, ymax=70,
    ylabel={Taux d'urbanisation (\%)},
    xlabel={Années de recensement / estimation},
    symbolic x coords={1960,1976,1987,2005,2010,2020},
    xtick=data,
    xticklabel style={rotate=45, anchor=east, font=\small},
    ymajorgrids=true,
    grid style=dashed,
    nodes near coords,
    nodes near coords align={vertical},
    every node near coord/.append style={font=\scriptsize, popdark},
    title style={font=\bfseries, popdark},
    title={Évolution du taux d'urbanisation au Cameroun (1960--2020)}
]
\addplot[popblue, fill=popblueL] coordinates {
    (1960,12)
    (1976,20)
    (1987,35)
    (2005,48)
    (2010,52)
    (2020,57)
};
\end{axis}
\end{tikzpicture}
\legende{Figure 1 : Évolution du taux d'urbanisation au Cameroun. Sources : RGPH 1976, 1987, 2005 ; Estimations INS et ONU World Urbanization Prospects 2018. Le franchissement du seuil des 50 \% intervient vers 2010.}
\end{popfigure}

\begin{exemplebox}[Douala et Yaoundé : deux métropoles primatiques]
Les deux principales villes concentrent à elles seules près de \textbf{40 \% de la population urbaine totale} (estimation 2020).
\begin{itemize}
    \item \textbf{Douala} (capitale économique, Région du Littoral) : $\approx$ \num{3,8} millions d'habitants (agglomération). Port autonome, aéroport international, zone industrielle de Bassa/Bonabéri, siège des grandes entreprises (SABC, SOCAPALM, banques, télécoms).
    \item \textbf{Yaoundé} (capitale politique, Région du Centre) : $\approx$ \num{3,5} millions d'habitants (agglomération). Siège des institutions (Présidence, Assemblée, Ministères), ambassades, universités, hôpitaux de référence, zone industrielle de Mvan/Olembé.
\end{itemize}
Cette \cle{primatie} (domination de la plus grande ville sur les suivantes) est forte : Douala est environ 4 fois plus peuplée que la 3e ville (Garoua ou Bamenda selon les définitions d'agglomération).
\end{exemplebox}

\begin{savaistu}[Yaoundé, capitale depuis 1922]
C'est en \textbf{1922}, sous l'administration française (mandat SDN), que Yaoundé (alors « Jaunde ») devient officiellement la capitale du Cameroun français, remplaçant Douala (capitale allemande puis française provisoire). Le choix répondait à des raisons sanitaires (paludisme moins endémique en altitude), stratégiques (centrage territorial) et symboliques (rupture avec l'héritage allemand côtière). Douala est restée la capitale économique et portuaire.
\end{savaistu}

\courssub{Hiérarchie et localisation des villes camerounaises}

Au-delà des deux métropoles, le réseau urbain camerounais s'organise en une hiérarchie fonctionnelle à plusieurs niveaux.

\begin{propriete}[Hiérarchie urbaine du Cameroun (schéma simplifié)]
\begin{enumerate}
    \item \textbf{Métropoles nationales (rang 1) :} \textbf{Douala}, \textbf{Yaoundé}. Fonctions de commandement supra-régional (économie, politique, international, enseignement supérieur, santé de pointe). Population $> 3$ millions.
    \item \textbf{Villes régionales / Métropoles régionales (rang 2) :} \textbf{Garoua} (Nord), \textbf{Bamenda} (Nord-Ouest), \textbf{Bafoussam} (Ouest), \textbf{Maroua} (Extrême-Nord), \textbf{Ngaoundéré} (Adamaoua), \textbf{Bertoua} (Est), \textbf{Ebolowa} (Sud), \textbf{Buea} (Sud-Ouest, capitale régionale). Chefs-lieux de région, fonctions administratives, commerciales, universitaires régionales. Population \num{200 000} -- \num{600 000}.
    \item \textbf{Villes départementales / Petites villes moyennes (rang 3) :} Chefs-lieux de département (ex. : \textbf{Kribi}, \textbf{Limbé}, \textbf{Foumban}, \textbf{Dschang}, \textbf{Yagoua}, \textbf{Meiganga}, \textbf{Batouri}, \textbf{Abong-Mbang}). Fonctions de relais administratif et de marché pour l'arrière-pays rural. Population \num{20 000} -- \num{150 000}.
    \item \textbf{Bourgs et localités urbaines émergentes (rang 4) :} Chefs-lieux d'arrondissement ou localités franchissant le seuil de 5 000 hab. (ex. : \textbf{Banyo}, \textbf{Tibati}, \textbf{Mokolo}, \textbf{Magba}, \textbf{Akono}). Fonctions de proximité (marché hebdomadaire, dispensaire, CEG).
\end{enumerate}
\end{propriete}

\begin{popfigure}
\begin{center}
\includegraphics[width=\linewidth,height=10cm,keepaspectratio]{N03/cmr_regions.jpg}
\end{center}
\legende{Figure 2 : Carte des villes du Cameroun (noms en anglais), fond pour lire la hiérarchie urbaine. Les deux métropoles, Douala (grand point noir, sur la côte) et Yaoundé (étoile, capitale), dominent ; viennent ensuite les capitales régionales (Bamenda, Bafoussam, Garoua, Maroua, Ngaoundéré, Bertoua...), puis les villes secondaires. Sur le croquis demandé, on dessine des cercles proportionnels à la population d'agglomération. \\ Source : Wikimedia Commons, Peter Fitzgerald, CC BY 3.0.}
\end{popfigure}

\begin{methode}[Lire et légender une carte de répartition urbaine]
\begin{enumerate}
    \item \textbf{Identifier le fond de carte :} limites nationales, régionales, cours d'eau majeurs (Sanaga, Bénoué, Logone), axes routiers structurants (N1, N2, N3, N10).
    \item \textbf{Analyser la symbologie :} type de signes (cercles proportionnels, carrés, pastilles de couleur), légende des tailles (seuils démographiques), code couleur (fonctions, rang hiérarchique).
    \item \textbf{Repérer les pôles majeurs :} localiser Douala et Yaoundé (souvent surimposés), puis les capitales régionales alignées sur les axes Nord-Sud (N1, N2, N10) et Est-Ouest (N3, N6).
    \item \textbf{Décrire la distribution :} forte concentration côtière et dans les hautes terres de l'Ouest (bassin du Moungo, région de l'Ouest) ; maillage plus lâche au Nord (soudanien) et à l'Est (forestier) ; rôle des axes de transport (chemin de fer Transcam, routes bitumées) comme lignes d'urbanisation.
    \item \textbf{Rédiger le commentaire :} structurer en (a) Hiérarchie (métropoles, villes moyennes, petites villes), (b) Spatialisation (littoral, axe Yaoundé-Garoua, axe Douala-Bafoussam-Bamenda, vides du Nord-Est et Sud-Est), (c) Facteurs (histoire coloniale, ressources, transports, administration).
\end{enumerate}
\end{methode}

\courssub{La ville, moteur de l'intégration et du développement : premiers constats}

Les villes camerounaises ne sont pas de simples concentrations humaines ; ce sont les \cle{nœuds} du territoire national.
\begin{itemize}
    \item \textbf{Intégration nationale :} Brassage ethnique, linguistique (français, anglais, langues nationales, camfranglais), religieux (mosquées, églises, temples côte à côte). Exemple : le quartier \textbf{Briqueterie} à Yaoundé ou \textbf{New Bell} à Douala sont des \cle{quartiers cosmopolites} par excellence.
    \item \textbf{Impulsion économique :} Concentration du PIB (Douala + Yaoundé $\approx$ 50 \% du PIB national), siège des banques (BEAC, banques commerciales), ports (Douala, Kribi), aéroports, zones industrielles, marchés de gros.
    \item \textbf{Polarisation des services :} Hôpitaux de référence (Hôpital Central, Hôpital Général, CHU), grandes écoles et universités (Université de Yaoundé I/II, Université de Douala, Ngaoundéré, Buea, Maroua, Dschang), administrations centrales.
\end{itemize}

Cependant, cette concentration génère les \textbf{problèmes urbains} (logement, transport, eau, électricité, déchets, insécurité, anarchie foncière) qui feront l'objet des sections 4 et 5. Comprendre la structure urbaine, c'est déjà disposer des clés pour analyser ces dysfonctionnements.

\begin{exoresolu}[Commentaire guidé : Figure 1 (Évolution du taux d'urbanisation)]
\textbf{Énoncé :} À l'aide de la Figure 1, décrire l'évolution du taux d'urbanisation au Cameroun entre 1960 et 2020 et proposer deux facteurs explicatifs.

\tcblower
\textbf{Solution détaillée :}
\begin{enumerate}
    \item \textbf{Description chiffrée :} Le taux d'urbanisation passe d'environ \num{12} \% en 1960 à \num{57} \% en 2020, soit une multiplication par \num{4,75} en 60 ans. La courbe montre une accélération nette entre 1976 (\num{20} \%) et 2005 (\num{48} \%), puis un ralentissement de la progression (mais poursuite de la hausse) après 2010.
    \item \textbf{Facteur 1 : L'exode rural massif.} La crise du secteur agricole (baisse des cours café/cacao années 1980, enclavement des bassins de production), la pression foncière dans les Hautes Terres de l'Ouest et l'insécurité dans le Nord (Boko Haram) et le Nord-Ouest/Sud-Ouest (crise anglophone depuis 2016) poussent les ruraux vers les villes.
    \item \textbf{Facteur 2 : L'attractivité urbaine (pôle d'impulsion).} Concentration des emplois salariés (fonction publique, industries, services), des infrastructures (écoles, hôpitaux, électricité, eau) et de la vie culturelle. Les villes sont perçues comme des espaces de « modernité » et de mobilité sociale.
\end{enumerate}
\textbf{Note de rédaction :} Utilise toujours des chiffres lus sur le graphique (dates, valeurs). Lie description et explication. N'oublie pas de citer le franchissement du seuil des 50 \% vers 2010 (transition urbaine).
\end{exoresolu}

\begin{exercice}[Entraînement : Hiérarchie urbaine]
\begin{enumerate}
    \item Classer les villes suivantes dans l'ordre hiérarchique (Rang 1 à 4) : \textbf{Kribi, Douala, Bafoussam, Magba, Yaoundé, Bertoua, Ngaoundéré, Mokolo}.
    \item Sur une carte muette du Cameroun (fournie à part), localiser ces 8 villes et représenter leur rang hiérarchique par des cercles proportionnels (choisis une échelle : 1 mm pour 100 000 hab).
    \item Expliquer pourquoi la ville de \textbf{Kribi} (chef-lieu de département, port en eau profonde en construction) pourrait voir son rang hiérarchique évoluer dans les prochaines décennies.
\end{enumerate}
\end{exercice}

\begin{corrige}[Exercice n°1 : Hiérarchie urbaine]
\begin{enumerate}
    \item \textbf{Rang 1 (Métropoles) :} Douala, Yaoundé. \textbf{Rang 2 (Capitales régionales) :} Bafoussam (Ouest), Bertoua (Est), Ngaoundéré (Adamaoua). \textbf{Rang 3 (Villes départementales) :} Kribi (Océan). \textbf{Rang 4 (Bourgs/Chefs-lieux d'arrondissement) :} Magba (Noun), Mokolo (Mayo-Tsanaga).
    \item \textit{Attendu cartographique :} Cercles : Douala/Yaoundé $\approx$ 35-38 mm ; Bafoussam/Bertoua/Ngaoundéré $\approx$ 3-5 mm ; Kribi $\approx$ 1-2 mm ; Magba/Mokolo $\approx$ 0.5-1 mm. Légende obligatoire : Titre, Échelle des cercles, Orientation, Source.
    \item \textbf{Évolution potentielle de Kribi :} Le projet de \textbf{port en eau profonde de Kribi} (phase 1 opérationnelle depuis 2014, phase 2 en cours), le pipeline Tchad-Cameroun, la zone industrielle et le terminal gazier (GNL) transforment Kribi en \cle{pôle de croissance} d'envergure sous-régionale. Cela attire des investissements, des emplois, des populations $\to$ croissance démographique rapide $\to$ montée en rang (vers rang 2 voire 1 à long terme si agglomération avec Campo/Lolodorf). C'est un cas d'\cle{urbanisation impulsée par un grand équipement infrastructurel}.
\end{enumerate}
\end{corrige}

\coursec{La ville, espace d'intégration nationale}

Dans la section précédente, nous avons vu que le Cameroun connaît une urbanisation rapide : plus d'un Camerounais sur deux vit aujourd'hui en ville (taux d'urbanisation estimé à environ 58 \% en 2023, source INS/ONU-Habitat), et les deux métropoles, Douala et Yaoundé, concentrent à elles seules plusieurs millions d'habitants (environ 3,5 à 4 millions chacune, ordres de grandeur 2023). Mais d'où viennent tous ces citadins ? Ils viennent des dix régions du pays, de tous les villages, de toutes les ethnies, de toutes les confessions religieuses. La ville camerounaise n'est donc pas seulement un espace peuplé : c'est un \cle{espace d'intégration nationale}, c'est-à-dire un lieu où se rencontrent, se mélangent et apprennent à vivre ensemble des populations que tout séparait à l'origine. C'est ce phénomène fondamental que nous allons étudier maintenant.

\courssub{La ville, cadre de brassage des populations}

\textbf{Une observation pour commencer.} Entre dans un taxi-brousse à l'agence de voyage de Bafoussam, de Maroua ou de Buéa : presque tous les passagers descendent à Douala ou à Yaoundé. Interroge-les : l'un va y chercher du travail, l'autre y poursuivre des études, un troisième y rejoindre un frère déjà installé. Cette scène quotidienne illustre le mécanisme essentiel : la ville attire, et cette attraction provoque un \cle{brassage des populations}.

\begin{definition}[Brassage des populations]
Le \cle{brassage des populations} désigne le mélange, sur un même espace urbain, de personnes originaires de régions, d'ethnies et de milieux sociaux différents. Il résulte principalement des \cle{migrations} (rurales-urbaines et interurbaines) vers les grandes villes, qui offrent emplois, écoles, hôpitaux et services.
\end{definition}

\textbf{Le mécanisme.} Les migrations vers les villes s'expliquent par un double jeu de facteurs : les \emph{facteurs de répulsion} (pauvreté rurale, manque de terres, absence d'écoles secondaires ou d'hôpitaux dans certains villages, exode des jeunes) et les \emph{facteurs d'attraction} (emplois salariés, universités, administrations, marchés, loisirs). Douala, capitale économique, attire surtout pour l'emploi et le commerce ; Yaoundé, capitale politique, attire pour l'administration, les grandes écoles et les universités. Mais les villes secondaires (Bafoussam, Bamenda, Garoua, Maroua, Ngaoundéré) accueillent aussi des migrants de leurs arrière-pays et d'ailleurs.

\begin{exemplebox}[Yaoundé, mosaïque nationale]
On estime que plus de \cle{200 groupes ethniques} sont représentés dans la ville de Yaoundé, alors que le Cameroun en compte environ 250 au total (ordre de grandeur couramment admis). Autrement dit, presque toute la diversité humaine du pays se retrouve concentrée dans sa capitale. On y entend parler ewondo, bamiléké, fulfuldé, duala, bassa, haoussa, pidgin english, français et anglais, parfois dans la même rue.
\end{exemplebox}

\begin{popfigure}
\begin{tikzpicture}[scale=1, every node/.style={font=\small}]
% Tronc : villes d'accueil
\node[draw, fill=popblue, text=white, rounded corners, minimum width=3.4cm, minimum height=1.1cm, align=center, font=\small\bfseries] (douala) at (0,0) {DOUALA\\ (capitale économique)};
\node[draw, fill=poppurple, text=white, rounded corners, minimum width=3.4cm, minimum height=1.1cm, align=center, font=\small\bfseries] (yaounde) at (0,-3.2) {YAOUNDÉ\\ (capitale politique)};
% Régions sources
\node[draw, fill=popgreenL, rounded corners, align=center] (ouest) at (-6.2,1.6) {Ouest, Nord-Ouest,\\ Sud-Ouest, Littoral};
\node[draw, fill=popgreenL, rounded corners, align=center] (nord) at (6.2,1.6) {Adamaoua, Nord,\\ Extrême-Nord};
\node[draw, fill=poporangeL, rounded corners, align=center] (sud) at (-6.2,-4.8) {Centre, Sud, Est};
\node[draw, fill=poporangeL, rounded corners, align=center] (diaspora) at (6.2,-4.8) {Autres villes,\\ étranger (Tchad, Nigeria...)};
% Flèches
\draw[-{Stealth[length=3mm]}, line width=2.2pt, popblue] (ouest.east) -- (douala.west);
\draw[-{Stealth[length=3mm]}, line width=1.6pt, popblue] (nord.west) -- (douala.east);
\draw[-{Stealth[length=3mm]}, line width=2.2pt, poppurple] (sud.east) -- (yaounde.west);
\draw[-{Stealth[length=3mm]}, line width=1.6pt, poppurple] (diaspora.west) -- (yaounde.east);
\draw[-{Stealth[length=3mm]}, line width=1pt, popmuted] (ouest.south east) -- (yaounde.north west);
\draw[-{Stealth[length=3mm]}, line width=1pt, popmuted] (nord.south west) -- (yaounde.north east);
% Motifs
\node[align=left, font=\scriptsize, text=popdark] at (-6.2,0.4) {Motifs : commerce,\\ emploi, études};
\node[align=left, font=\scriptsize, text=popdark] at (6.2,0.4) {Motifs : emploi,\\ élevage, commerce};
\node[align=left, font=\scriptsize, text=popdark] at (-6.2,-6.2) {Motifs : administration,\\ universités, fonctions};
\node[align=left, font=\scriptsize, text=popdark] at (6.2,-6.2) {Motifs : études,\\ travail, famille};
\end{tikzpicture}
\legende{Schéma en arbre des flux migratoires vers Douala et Yaoundé. L'épaisseur des flèches est proportionnelle à l'importance relative des flux (schématique). Toutes les régions du Cameroun alimentent les deux métropoles, ce qui en fait des espaces de brassage national.}
\end{popfigure}

\courssub{Le brassage des cultures}

Le brassage des populations entraîne mécaniquement un \cle{brassage des cultures} : dans les quartiers des grandes villes coexistent des langues, des cuisines, des musiques et des modes de vie venus de tout le pays.

\begin{itemize}
\item \textbf{Les langues.} À côté du français et de l'anglais (langues officielles), on parle dans les rues de Douala ou de Yaoundé des dizaines de langues nationales. Le \emph{camfranglais} et le pidgin english, nés du mélange urbain, sont eux-mêmes des produits du brassage.
\item \textbf{Les cuisines.} Dans un même quartier de Yaoundé, on peut manger l'eru du Sud-Ouest, le ndolé du Littoral, le poulet DG, le koki du Centre, le foléré et le couscous de mil du Nord. Les restaurants « à la carte régionale » sont une spécialité des grandes villes.
\item \textbf{Les musiques et les fêtes.} Makossa, bikutsi, assiko, mbansah, musiques saheliennes se côtoient dans les bars et les maquis. Les mariages urbains réunissent souvent des familles de deux régions différentes : les \emph{mariages mixtes} interethniques, de plus en plus fréquents en ville, sont un indicateur fort de l'intégration.
\item \textbf{Les modes de vie.} L'habillement (boubou, kaba ngondo, gandoura, tenue occidentale), les associations (tontines, réunions de ressortissants) et les fêtes traditionnelles célébrées en ville témoignent d'une coexistence culturelle quotidienne.
\end{itemize}

\courssub{Le brassage des religions}

La ville camerounaise est également un lieu de \cle{cohabitation religieuse}. Chrétiens (catholiques, protestants, pentecôtistes), musulmans et adeptes des religions traditionnelles y vivent côte à côte. À Yaoundé comme à Douala, une mosquée peut se dresser à quelques centaines de mètres d'une église, et les fêtes religieuses (Noël, Pâques, Ramadan, fête du mouton, rites traditionnels) rythment la vie de tous les quartiers. Les écoles confessionnelles (catholiques, protestantes, islamiques) scolarisent des enfants de toutes origines. Cette coexistence, globalement pacifique au Cameroun, est un acquis précieux de la vie urbaine : la ville apprend la tolérance par la pratique du voisinage quotidien.

\courssub{Le quartier cosmopolite : définition et exemples}

\begin{definition}[Quartier cosmopolite]
Un \cle{quartier cosmopolite} est un quartier urbain qui accueille des populations de \textbf{toutes les régions du pays} (et parfois de l'étranger), de cultures, de langues et de religions différentes, vivant et travaillant ensemble dans un même espace. Il se caractérise par la diversité de ses habitants, de ses activités (commerces, ateliers, services) et de ses lieux de culte.
\end{definition}

\begin{exemplebox}[Quatre quartiers cosmopolites camerounais]
\begin{itemize}
\item \textbf{La Briqueterie} (Yaoundé) : quartier populaire et commerçant, réputé pour ses marchés, ses ateliers de couture et de mécanique, sa forte communauté musulmane venue du Nord cohabitant avec des populations du Centre, de l'Ouest et d'ailleurs.
\item \textbf{Nkoulouloun} (Yaoundé) : quartier de forte densité où se côtoient des habitants originaires de presque toutes les régions.
\item \textbf{New-Bell} (Douala) : vaste quartier historique, cosmopolite depuis l'époque coloniale, où vivent Duala, Bamiléké, ressortissants du Nord, anglophones et étrangers.
\item \textbf{Le quartier Haoussa} (Garoua) : témoigne de l'ancienneté des migrations commerciales transsahariennes et de la cohabitation entre Haoussa, Peuls, et populations venues du sud du pays (fonctionnaires, commerçants, étudiants).
\end{itemize}
\end{exemplebox}

\begin{savaistu}[D'où vient le nom « New-Bell » ?]
Le quartier New-Bell de Douala doit son nom à la \textbf{dynastie Bell}, l'une des grandes familles royales duala. À l'époque coloniale allemande (début du XX\textsuperscript{e} siècle), les autorités ont déplacé les populations africaines hors de la ville européenne (Joss, Bonanjo) vers de nouveaux quartiers : le « nouveau Bell » (New-Bell) fut créé pour accueillir les sujets du roi Bell et, très vite, des migrants venus de tout le pays. New-Bell est ainsi l'un des plus anciens quartiers cosmopolites d'Afrique centrale.
\end{savaistu}

\begin{definition}[Voirie urbaine]
La \cle{voirie urbaine} désigne l'ensemble des voies de circulation (rues, avenues, boulevards, places, ponts, ronds-points) et de leurs équipements (trottoirs, éclairage public, signalisation, réseaux enterrés) qui structurent l'espace de la ville et assurent la mobilité des personnes et des biens. Son maillage (orthogonal, radioconcentrique, irrégulier) reflète l'histoire de la croissance urbaine.
\end{definition}

\begin{popfigure}
\begin{tikzpicture}[scale=0.92, every node/.style={font=\scriptsize}]
% Fond / voirie
\draw[fill=popcream, draw=popline] (-0.4,-0.4) rectangle (12.4,7.4);
% Avenue principale (horizontale)
\draw[fill=popmuted!40, draw=none] (0,3.2) rectangle (12,4.2);
\node[rotate=0, text=popdark, font=\scriptsize\bfseries] at (2.2,3.7) {Avenue commerçante};
% Rue secondaire (verticale)
\draw[fill=popmuted!40, draw=none] (5.6,0) rectangle (6.4,7);
\node[rotate=90, text=popdark, font=\scriptsize\bfseries] at (6.0,1.2) {Rue};
% 1. Marché
\draw[fill=poporangeL, draw=poporange, thick] (0.4,4.6) rectangle (3.4,7.0);
\node[align=center, font=\scriptsize] at (1.9,6.5) {\textbf{1. Marché}};
\node[align=center, font=\tiny] at (1.9,5.6) {vivres, tissus,\\ épices du Nord,\\ poisson fumé};
% 2. Mosquée
\draw[fill=popgreenL, draw=popgreen, thick] (3.8,4.6) rectangle (5.4,7.0);
\node[align=center, font=\scriptsize] at (4.6,6.4) {\textbf{2. Mosquée}};
\node[align=center, font=\tiny] at (4.6,5.4) {communauté\\ musulmane};
% 3. Église
\draw[fill=poppurpleL, draw=poppurple, thick] (6.6,4.6) rectangle (8.4,7.0);
\node[align=center, font=\scriptsize] at (7.5,6.4) {\textbf{3. Église}};
\node[align=center, font=\tiny] at (7.5,5.4) {paroisse\\ chrétienne};
% 4. École
\draw[fill=popblueL, draw=popblue, thick] (8.8,4.6) rectangle (11.8,7.0);
\node[align=center, font=\scriptsize] at (10.3,6.4) {\textbf{4. École}};
\node[align=center, font=\tiny] at (10.3,5.4) {enfants de toutes\\ origines, français\\ et langues locales};
% 5. Ateliers
\draw[fill=popgoldL, draw=popgold, thick] (0.4,0.4) rectangle (3.4,2.8);
\node[align=center, font=\scriptsize] at (1.9,2.3) {\textbf{5. Ateliers}};
\node[align=center, font=\tiny] at (1.9,1.3) {couture, mécanique,\\ menuiserie,\\ cordonnerie};
% 6. Habitat
\draw[fill=poppinkL, draw=poppink, thick] (3.8,0.4) rectangle (5.4,2.8);
\node[align=center, font=\scriptsize] at (4.6,2.3) {\textbf{6. Habitat}};
\node[align=center, font=\tiny] at (4.6,1.3) {cours communes,\\ maisons de briques\\ (d'où le nom)};
% 7. Maquis / restos
\draw[fill=poporangeL, draw=poporange, thick] (6.6,0.4) rectangle (8.4,2.8);
\node[align=center, font=\scriptsize] at (7.5,2.3) {\textbf{7. Maquis}};
\node[align=center, font=\tiny] at (7.5,1.3) {cuisines de toutes\\ les régions,\\ musiques variées};
% 8. Stade / Terrain
\draw[fill=popgreenL, draw=popgreen, thick] (8.8,0.4) rectangle (11.8,2.8);
\node[align=center, font=\scriptsize] at (10.3,2.3) {\textbf{8. Terrain de sport}};
\node[align=center, font=\tiny] at (10.3,1.3) {tournois de quartier,\\ jeunes de toutes\\ ethnies};
% Flèche Nord
\draw[-{Stealth[length=3mm]}, thick, popdark] (11.5,6.8) -- (11.5,7.2);
\node[font=\scriptsize\bfseries, anchor=south] at (11.5,7.25) {N};
\end{tikzpicture}
\legende{Croquis schématique d'un quartier cosmopolite (inspiré de la Briqueterie, Yaoundé). 1 : marché multi-produits ; 2 et 3 : lieux de culte voisins (cohabitation religieuse) ; 4 : école, creuset du brassage ; 5 : ateliers du secteur informel ; 6 : habitat dense ; 7 : restauration et musiques de toutes les régions ; 8 : sport, facteur de mixité. Échelle indicative : le quartier s'étend sur quelques centaines de mètres.}
\end{popfigure}

\begin{methode}[Construire un croquis de quartier à partir d'une photographie]
Au Baccalauréat, on peut te demander de transformer une photographie de quartier en croquis. Voici la démarche :
\begin{enumerate}
\item \textbf{Observer et délimiter} : identifier le premier plan, l'arrière-plan, et les grandes lignes de l'espace (routes, places, alignements de maisons).
\item \textbf{Choisir un plan} : découper l'espace en 3 ou 4 zones fonctionnelles (habitat, commerce, équipements, voirie).
\item \textbf{Placer les repères} : dessiner d'abord la \cle{voirie urbaine} (elle structure le croquis), puis les bâtiments et équipements par des formes simples (rectangles, points, symboles).
\item \textbf{Légender par des chiffres} : ne jamais écrire les noms en toutes lettres sur le croquis ; placer des chiffres renvoyant à une légende numérotée sous le croquis.
\item \textbf{Ajouter titre, orientation (flèche du nord) et échelle indicative} : un croquis sans titre ni légende perd des points.
\end{enumerate}
\end{methode}

\courssub{Le rôle des institutions urbaines dans l'intégration}

Le brassage ne se fait pas tout seul : ce sont les \cle{institutions urbaines} qui le transforment en véritable intégration.

\begin{itemize}
\item \textbf{L'école} : elle accueille des enfants de toutes origines, leur enseigne les mêmes programmes en français ou en anglais, et fabrique une culture commune (hymne national, valeurs civiques, amitiés interethniques).
\item \textbf{Le marché} : lieu d'échanges économiques mais aussi linguistiques et culturels ; on y négocie dans plusieurs langues, on y découvre les produits des autres régions.
\item \textbf{Le stade et le sport} : les tournois de quartier et les matches du championnat national (Canon, Tonnerre de Yaoundé, Union de Douala...) créent des appartenances qui dépassent l'ethnie.
\item \textbf{Les lieux de culte} : églises et mosquées rassemblent des fidèles de toutes régions autour d'une même foi.
\item \textbf{Les administrations et entreprises} : la fonction publique affecte les agents dans tout le pays ; un fonctionnaire de l'Extrême-Nord peut servir à Yaoundé aux côtés de collègues du Sud-Ouest. Cette mobilité administrative est un puissant instrument d'intégration nationale.
\end{itemize}

\courssub{Les limites de l'intégration urbaine}

Il serait faux de croire que la ville efface toutes les différences. L'intégration a des \cle{limites} qu'un bon candidat au Bac doit savoir mentionner avec nuance.

\begin{itemize}
\item \textbf{Les regroupements communautaires} : les migrants s'installent souvent près de leurs « frères » de village ou d'ethnie, ce qui crée des îlots à dominante régionale dans certains quartiers. Les associations de ressortissants (réunions d'originaires d'un même village) maintiennent les liens avec la région d'origine.
\item \textbf{Les tensions ponctuelles} : des conflits intercommunautaires peuvent éclater (disputes foncières, rivalités commerciales, incidents entre communautés), même si le Cameroun a globalement préservé la paix sociale dans ses villes.
\item \textbf{Les inégalités sociales} : la ville intègre les cultures, mais elle sépare aussi les riches des pauvres (quartiers résidentiels contre quartiers précaires) : l'intégration culturelle ne supprime pas la ségrégation sociale.
\end{itemize}

\begin{attention}[Erreur fréquente au Bac]
Ne rédige jamais que « le brassage urbain supprime les identités régionales et ethniques ». C'est \textbf{faux} : en ville, les identités d'origine \emph{coexistent} avec l'identité nationale ; elles se recomposent (associations de ressortissants, fêtes villageoises célébrées en ville) sans disparaître. La formulation correcte est : « la ville favorise l'intégration nationale par le brassage, tout en préservant les identités particulières ».
\end{attention}

\begin{exoresolu}[Commenter la phrase : « La ville camerounaise est un creuset de l'intégration nationale »]
Énoncé. En t'appuyant sur des exemples précis, montre en une dizaine de lignes que la ville camerounaise est un espace d'intégration nationale, puis nuance ce constat.
\tcblower
\textbf{Solution détaillée.} La ville camerounaise est un espace d'intégration nationale parce qu'elle y brasse les populations : les migrations venues des dix régions convergent vers Douala et Yaoundé, où l'on compte plus de 200 groupes ethniques représentés. Ce brassage est d'abord culturel : dans un quartier comme New-Bell à Douala ou la Briqueterie à Yaoundé, coexistent des dizaines de langues, des cuisines de toutes les régions (ndolé, eru, couscous de mil) et des musiques variées (makossa, bikutsi). Il est aussi religieux : chrétiens, musulmans et adeptes des religions traditionnelles y vivent en voisins pacifiques. Les institutions urbaines renforcent ce processus : l'école forge une culture commune, le marché multiplie les échanges, le stade crée des appartenances supra-ethniques, et la fonction publique mélange les agents de toutes origines. Cependant, cette intégration a des limites : les regroupements communautaires, les associations de ressortissants et de rares tensions intercommunautaires montrent que les identités régionales ne disparaissent pas ; elles coexistent avec l'identité nationale que la ville contribue à construire.
\end{exoresolu}

\begin{aretenir}[L'essentiel de la section]
\begin{itemize}
\item La ville camerounaise est un \cle{espace d'intégration nationale} : elle brasse les populations (migrations de toutes les régions), les cultures (langues, cuisines, musiques) et les religions (chrétiens, musulmans, religions traditionnelles).
\item Le \cle{quartier cosmopolite} (Briqueterie, Nkoulouloun, New-Bell, quartier Haoussa) est la forme spatiale la plus visible de ce brassage.
\item La \cle{voirie urbaine} structure l'espace du quartier et facilite la mixité des usages.
\item Les institutions urbaines (école, marché, stade, lieux de culte, administration) transforment le brassage en intégration durable.
\item L'intégration a des limites : regroupements communautaires, tensions ponctuelles, ségrégation sociale. Le brassage ne supprime pas les identités, il les fait coexister.
\end{itemize}
\end{aretenir}

\coursec{La ville, centre d'impulsion économique et de polarisation des services}

Dans la section précédente, nous avons montré que la ville camerounaise est un espace d'intégration nationale : elle accueille, brasse et mélange des populations venues de toutes les régions, de toutes les cultures et de toutes les religions. Mais pourquoi ces populations quittent-elles leurs villages pour s'installer en ville ? La réponse tient en grande partie à ce que la ville \emph{offre} : des emplois, des commerces, des banques, des hôpitaux, des écoles, des routes. La ville n'est donc pas seulement un creuset humain : c'est aussi le moteur économique du pays et le lieu où se concentrent les services. C'est ce double rôle que nous allons maintenant étudier.

\courssub{Une concentration exceptionnelle des activités économiques}

\textbf{Intuition et observation.} Fais une expérience simple : parcours mentalement le trajet entre le village de Bangangté et la ville de Douala. Au village, tu croises quelques boutiques, un marché hebdomadaire, un dispensaire. En approchant de Douala, tout change : usines, entrepôts, banques à chaque carrefour, supermarchés, agences de voyage, sièges d'entreprises, grues du port. Cette différence visible à l'œil nu traduit une loi fondamentale de la géographie urbaine : \cle{les activités économiques modernes se concentrent dans les villes}.

\begin{definition}[Polarisation]
La \cle{polarisation} désigne la capacité d'une ville à attirer vers elle les hommes, les capitaux, les biens et les activités de tout l'espace environnant (son \cle{aire d'influence}), puis à redistribuer vers cet espace des biens manufacturés et des services. Une ville polarise parce qu'elle concentre les fonctions de commandement (administrations, sièges d'entreprises), de production (industries) et de services (banques, hôpitaux, universités).
\end{definition}

\textbf{Le mécanisme de la concentration.} Pourquoi les activités se concentrent-elles en ville ? Le raisonnement s'établit en trois étapes.

\begin{enumerate}
\item \textbf{Les entreprises cherchent les villes} parce qu'elles y trouvent une main-d'œuvre nombreuse et qualifiée, des clients solvables, des infrastructures (électricité, eau, routes, télécommunications) et la proximité des administrations qui délivrent les autorisations.
\item \textbf{Cette concentration attire de nouveaux habitants}, qui viennent chercher du travail et des services ; la population urbaine augmente.
\item \textbf{Cette population croissante attire à son tour de nouvelles entreprises}, qui veulent profiter de ce marché. C'est un cercle cumulatif : la ville grossit parce qu'elle est déjà grosse.
\end{enumerate}

Au Cameroun, ce mécanisme explique que l'essentiel du tissu économique moderne soit concentré dans un petit nombre de villes, et d'abord dans deux métropoles : Douala et Yaoundé. On y trouve :
\begin{itemize}
\item la quasi-totalité des \cle{industries} de transformation (brasseries, cimenteries, huileries, conserveries, textile, bois) ;
\item le siège des \cle{banques} et des compagnies d'assurance (BICEC, SGBC, Afriland First Bank, UBA, etc.) ;
\item les grands \cle{commerces} : marchés centraux, supermarchés, grossistes, importateurs ;
\item les \cle{administrations} : ministères, directions générales, services déconcentrés, qui emploient des dizaines de milliers de fonctionnaires.
\end{itemize}

\begin{attention}[Erreur fréquente au Bac]
Beaucoup de candidats limitent le rôle économique des villes au seul \emph{commerce} (« les villes servent à acheter et à vendre »). C'est une réduction grave. Une ville comme Douala ou Yaoundé joue \textbf{trois rôles} qu'il faut distinguer dans ta copie : un rôle \textbf{industriel} (production et transformation), un rôle \textbf{commercial} (échanges, import-export) et un rôle de \textbf{services} (banques, administrations, santé, éducation, transports). Une copie qui n'évoque que les marchés perd la moitié des points.
\end{attention}

\courssub{Douala, capitale économique du Cameroun}

\textbf{Situation concrète.} Un container de cacao parti de Sangmélima, un camion de bananes venu de Penja, une cargaison de pétrole de Kribi : tous convergent vers le même point, le port de Douala, sur l'estuaire du Wouri. Inversement, les motos, les téléviseurs, les médicaments et les ciments importés débarquent à Douala avant d'irriguer tout le pays. Douala est la porte d'entrée et de sortie du Cameroun.

\begin{exemplebox}[Le port autonome de Douala, poumon du commerce extérieur]
Le \cle{port autonome de Douala} (PAD) traite \textbf{plus de 90 \% du commerce extérieur du Cameroun} (ordre de grandeur régulièrement cité par les sources officielles ; le trafic annuel est de l'ordre de 10 à 13 millions de tonnes). Il dessert aussi, par voie terrestre, le Tchad et la Centrafrique, pays enclavés qui dépendent de Douala pour leurs importations. Autour du port s'est développée la \cle{zone industrielle de Bonabéri}, sur la rive droite du Wouri : installations de stockage de la SONARA (raffinerie de Limbé), cimenteries (CIMENCAM), huileries, brasseries (SABC), usines de bois et de savonnerie. Douala concentre ainsi, avec sa banlieue, la majorité des emplois industriels du pays.
\end{exemplebox}

\textbf{Les autres fonctions économiques de Douala.} Au-delà du port et des usines, Douala est :
\begin{itemize}
\item le \textbf{premier marché du pays} : le marché Mboppi, le marché des fleurs, le marché Congo, les quartiers commerçants d'Akwa drainent grossistes et détaillants de tout le Cameroun et de la sous-région ;
\item la \textbf{place financière} nationale : sièges des principales banques, de la Bourse des valeurs mobilières de l'Afrique centrale (BVMAC), des compagnies d'assurance ;
\item le \textbf{siège des grandes entreprises} : la plupart des sociétés nationales et des filiales des firmes multinationales ont leur direction à Douala, même lorsque leur activité se situe ailleurs ;
\item un \textbf{nœud de transport} : l'aéroport international de Douala est le plus fréquenté du pays, et la ville est le point de départ des axes vers Yaoundé, l'Ouest et le Nord.
\end{itemize}

\begin{savaistu}[D'où vient le nom « Douala » ?]
\emph{Douala} signifie littéralement « les Duala », du nom du peuple riverain du Wouri, pêcheurs et commerçants qui servirent d'intermédiaires entre les Européens et l'intérieur du pays. La ville, fondée autour de comptoirs, fut la \textbf{première capitale du Kamerun allemand} (de 1884 à 1901, avant le transfert de l'administration à Yaoundé). C'est dire l'ancienneté de sa fonction de commandement : Douala commande l'économie camerounaise depuis près d'un siècle et demi.
\end{savaistu}

\begin{popfigure}
\begin{center}
\includegraphics[width=\linewidth,height=10.5cm,keepaspectratio]{N03/douala_map.jpg}
\end{center}
\legende{Plan de Douala (fond OpenStreetMap) pour lire les fonctions économiques et la voirie. On y repère : le Wouri et son estuaire ; le centre d'affaires (Akwa, Bonanjo) sur la rive gauche près du fleuve ; la zone industrielle de Bassa à l'est ; la zone industrielle de Bonabéri sur la rive droite, reliée à la ville par le pont sur le Wouri (route N3) ; l'aéroport au sud ; les grandes routes (en rouge et orange) qui partent vers l'intérieur du pays. Les ponts sur le Wouri sont des goulots d'étranglement entre les deux rives. \\ Source : Wikimedia Commons, contributeurs d'OpenStreetMap, CC BY-SA 3.0.}
\end{popfigure}

\courssub{Yaoundé, capitale politique et administrative}

Si Douala est le cœur économique, Yaoundé est le cerveau politique du pays. Capitale du Cameroun, elle concentre :
\begin{itemize}
\item les \cle{ministères} et la présidence de la République, l'Assemblée nationale et le Sénat : toutes les grandes décisions publiques s'y prennent ;
\item les \cle{sièges des institutions} : Cour suprême, ambassades, représentations des organisations internationales ;
\item les grandes \cle{universités} et écoles : l'Université de Yaoundé I et II, l'École nationale d'administration et de magistrature (ENAM), l'École militaire interarmées (EMIA), l'École nationale supérieure polytechnique (ENSP) de Yaoundé ;
\item un tissu économique alimenté par la fonction administrative : banques, hôtels, commerces, BTP (la construction des immeubles ministériels et des logements est un secteur majeur).
\end{itemize}

La complémentarité Douala--Yaoundé est un trait fondamental de l'organisation du territoire camerounais : l'une produit et échange, l'autre décide et administre. L'axe routier Douala--Yaoundé (environ 270 km par la route) est de ce fait l'artère la plus fréquentée du pays, et sa mise à deux fois deux voies illustre l'effort de l'État pour fluidifier cette relation vitale.

\courssub{La polarisation des services}

Au-delà de l'économie, les villes concentrent les \cle{services} dont la population a besoin, et ce d'autant plus que ces services sont rares et spécialisés. On parle de \emph{hiérarchie des services} : plus un service est rare (chirurgie spécialisée, université, ambassade), plus il se situe dans une grande ville et plus son aire d'influence est vaste.

\begin{itemize}
\item \textbf{Santé} : les hôpitaux de référence (Centre hospitalier universitaire de Yaoundé, hôpital Laquintinie de Douala, hôpital général de Douala) reçoivent des malades évacués de tout le pays.
\item \textbf{Éducation} : les universités et grandes écoles sont presque toutes urbaines ; un jeune de l'Extrême-Nord qui veut faire médecine doit descendre à Yaoundé ou Douala.
\item \textbf{Médias et télécommunications} : sièges de la CRTV, des journaux et des opérateurs de téléphonie (MTN, Orange) dans les deux métropoles.
\item \textbf{Transports} : gares, aéroports internationaux, agences de voyage et gares routières (les « agences » du boulevard de la Réunification à Douala) concentrent les départs vers toutes les régions.
\end{itemize}

\begin{popfigure}
\begin{center}
\begin{tikzpicture}[font=\small]
% Ville centrale
\fill[popblue] (0,0) circle (1.15);
\node[white, font=\small\bfseries, align=center] at (0,0) {VILLE\\(services,\\industries,\\administration)};
% Villages
\foreach \x/\y/\n in {-4.5/2.6/{Village A}, -4.8/-0.2/{Village B}, -4.2/-2.8/{Village C}, 4.5/2.6/{Village D}, 4.8/-0.2/{Village E}, 4.2/-2.8/{Village F}}{
  \fill[popgreen] (\x,\y) circle (0.42);
  \node[font=\footnotesize] at (\x,\y) {\n};
}
% Flèches convergentes (produits vivriers, main-d'oeuvre, patients, élèves)
\foreach \x/\y in {-4.5/2.6, -4.8/-0.2, -4.2/-2.8, 4.5/2.6, 4.8/-0.2, 4.2/-2.8}{
  \draw[-{Stealth[length=2.6mm]}, line width=1.6pt, poporange] (\x,\y) -- (0,0);
}
% Flèches divergentes (biens manufacturés, services)
\foreach \x/\y in {-4.5/2.6, -4.8/-0.2, -4.2/-2.8, 4.5/2.6, 4.8/-0.2, 4.2/-2.8}{
  \draw[-{Stealth[length=2.6mm]}, line width=1.1pt, poppurple, dashed] (0,0) -- (\x,\y);
}
% Légende des flux
\node[font=\footnotesize, poporange, align=left] at (0,4.1) {Flèches pleines : produits vivriers, main-d'œuvre, malades, élèves $\longrightarrow$ la ville};
\node[font=\footnotesize, poppurple, align=left] at (0,-4.1) {Flèches pointillées : biens manufacturés, services, informations $\longrightarrow$ les campagnes};
\end{tikzpicture}
\end{center}
\legende{Schéma de la polarisation : la ville attire les hommes et les produits des campagnes (flux convergents) et redistribue biens manufacturés et services (flux divergents).}
\end{popfigure}

\textbf{Les relations villes--campagnes.} Ce schéma montre que la polarisation n'est pas un pillage à sens unique : c'est un \emph{échange}. Les campagnes alimentent les villes en \cle{produits vivriers} (plantains, manioc, maïs, arachides, poissons) et en \cle{main-d'œuvre} (les migrants qui alimentent l'exode rural) ; en retour, les villes approvisionnent les campagnes en \cle{biens manufacturés} (savons, tissus, motos, téléphones, matériaux), en services (santé, éducation, administration) et en informations. Les revenus gagnés en ville financent aussi, par les transferts des migrants, les écoles et les maisons des villages : la ville irrigue la campagne, la campagne nourrit la ville.

\courssub{La voirie urbaine, support de la circulation des hommes et des biens}

Aucune de ces fonctions ne serait possible sans un support matériel : les rues, avenues, boulevards et ponts qui structurent la ville.

\begin{definition}[Voirie urbaine]
La \cle{voirie urbaine} désigne l'ensemble des voies de circulation aménagées dans une ville : rues, avenues, boulevards, ronds-points, ponts, trottoirs et leurs équipements (éclairage, drainage, signalisation). Elle est le support de la circulation des hommes, des biens et des services, et elle structure l'organisation de l'espace urbain : c'est le long des grands axes que s'installent commerces, banques et administrations.
\end{definition}

\textbf{Rôle et état des réseaux.} À Douala, la voirie s'organise autour d'axes majeurs : le boulevard de la Liberté, le boulevard de la Réunification, l'axe Douala--Bonabéri qui franchit le Wouri par deux ponts (l'ancien pont de 1984 et le nouveau pont de la Réunification inauguré en 2021), qui restent toutefois un goulot d'étranglement aux heures de pointe. À Yaoundé, les grands axes (avenue Kennedy, boulevard du 20 mai, axe Yaoundé--Nsimalen) desservent le centre administratif et les quartiers résidentiels perchés sur les collines. Dans les deux villes, l'état de la voirie est contrasté : artères principales bitumées mais saturées aux heures de pointe, rues secondaires souvent dégradées ou non bitumées, encombrements chroniques qui ralentissent la circulation des biens et renchérissent les coûts de transport. L'entretien et l'extension de la voirie sont donc un enjeu économique majeur, pas seulement un confort.

\courssub{Le rayonnement des villes régionales}

La polarisation ne profite pas qu'à Douala et Yaoundé. Des \cle{villes régionales} relaient les deux métropoles et organisent l'espace national :
\begin{itemize}
\item \textbf{Garoua} (Nord) et \textbf{Maroua} (Extrême-Nord) : capitales régionales, marchés de collecte du coton et du bétail, industries (cotonneries, tanneries), relais commerciaux vers le Tchad et le Nigeria ;
\item \textbf{Bafoussam} (Ouest) : carrefour commercial de la haute terre de l'Ouest, marché de vivres et de café, plaque tournante des transports entre Douala, Yaoundé et le Nord-Ouest ;
\item \textbf{Ngaoundéré} (Adamaoua) : \cle{carrefour} stratégique, terminus de la ligne ferroviaire du Transcamerounais (exploitée par Camrail) venue de Douala ; c'est là que les marchandises du Nord basculent du rail à la route vers Garoua, Maroua, le Tchad et la Centrafrique.
\end{itemize}

\begin{methode}[Analyser une photographie de paysage urbain]
Au Bac, on te donnera souvent une photo d'une ville camerounaise à commenter. Procède en trois étapes :
\begin{enumerate}
\item \textbf{Décrire} objectivement ce que tu vois, du premier plan vers l'arrière-plan : bâtiments (hauts ou bas, modernes ou précaires), voirie (bitumée ou non, encombrée ou fluide), activités visibles (commerces, ateliers, port), population.
\item \textbf{Localiser} : identifie la ville et le quartier si possible, et situe-les dans l'espace national (capitale économique, ville régionale, quartier d'affaires ou quartier spontané).
\item \textbf{Interpréter} : relie chaque élément observé à une notion du cours (concentration des services, polarisation, voirie, industrialisation) et dégage le problème ou la dynamique que la photo illustre. Conclus par une phrase de synthèse.
\end{enumerate}
Piège à éviter : ne jamais rester au stade de la description sans interprétation géographique.
\end{methode}

\begin{exoresolu}[Commenter les fonctions économiques de Douala]
\textbf{Énoncé.} À partir de tes connaissances et du croquis de la leçon, montre en une quinzaine de lignes que Douala est la capitale économique du Cameroun.
\tcblower
\textbf{Solution détaillée.} Introduction : Douala, plus grande ville du Cameroun, concentre les fonctions économiques de commandement du pays. Développement en trois points : (1) une fonction \emph{portuaire et commerciale} dominante : le port autonome de Douala traite plus de 90 \% du commerce extérieur national et dessert le Tchad et la Centrafrique ; les grands marchés (Mboppi) et les sièges de banques en font la place financière du pays ; (2) une fonction \emph{industrielle} : la zone industrielle de Bonabéri regroupe cimenteries, huileries, brasseries et unités de stockage pétrolier, soit la majeure partie des emplois industriels nationaux ; (3) une fonction de \emph{nœud de transports et de services} : aéroport le plus fréquenté, voirie structurée autour des ponts sur le Wouri, redistribution des importations vers tout l'intérieur. Conclusion : par cette triple concentration, Douala polarise l'ensemble du territoire camerounais et mérite son titre de capitale économique, complémentaire de Yaoundé, capitale politique.
\end{exoresolu}

\begin{aretenir}[L'essentiel de la section]
La ville camerounaise est un \cle{centre d'impulsion économique} : elle concentre industries, banques, commerces et administrations. Douala, capitale économique (port, zone industrielle de Bonabéri, marchés, sièges d'entreprises), et Yaoundé, capitale politique et administrative (ministères, universités, institutions), polarisent l'espace national. La ville concentre aussi les services rares (hôpitaux de référence, universités, médias, transports) et entretient avec les campagnes des échanges complémentaires : vivres et main-d'œuvre contre biens manufacturés et services. La \cle{voirie urbaine} est le support de toutes ces circulations. Enfin, des villes régionales (Garoua, Maroua, Bafoussam, Ngaoundéré) relaient les métropoles et structurent les régions.
\end{aretenir}

\begin{exercice}[Pour t'entraîner]
1. Définis les termes \emph{polarisation} et \emph{voirie urbaine}. 2. Construis un tableau à deux colonnes comparant les fonctions de Douala et de Yaoundé. 3. Explique, en dix lignes, pourquoi Ngaoundéré est qualifiée de « carrefour ». 4. Sur un croquis simplifié du Cameroun, place Douala, Yaoundé, Bafoussam, Ngaoundéré, Garoua et Maroua, puis trace par des flèches les flux de vivres et de biens manufacturés entre ces villes et leurs arrière-pays.
\end{exercice}

\coursec{Les problèmes des villes camerounaises}

La ville camerounaise, que nous venons d'étudier comme un puissant \emph{centre d'impulsion économique} et un lieu de \emph{polarisation des services} (administrations, santé, éducation, marchés), paie aujourd'hui le prix fort de son succès. L'attractivité qu'elle exerce, combinée à une croissance démographique fulgurante, a engendré une urbanisation \emph{par défaut} : la ville s'étend plus vite qu'elle ne s'équipe. Derrière la vitalité des grands marchés comme Mfoundi à Yaoundé ou Sandaga à Douala, derrière l'effervescence des zones industrielles de Bassa ou de Bonabéri, se cachent des dysfonctionnements quotidiens qui minent la qualité de vie des citadins et freinent la compétitivité du pays. Comprendre ces problèmes, c'est saisir les limites actuelles de l'intégration nationale par l'urbain, là où le \cle{quartier cosmopolite} idéal se heurte à la fragmentation spatiale.

\courssub{Causes générales : une urbanisation rapide et mal maîtrisée}

L'origine commune à la quasi-totalité des maux urbains réside dans le décalage structurel entre la \textbf{vitesse de croissance de la population urbaine} et la \textbf{capacité des pouvoirs publics à planifier, équiper et gérer l'espace}.

\begin{itemize}
    \item \textbf{L'exode rural massif et continu :} Poussés par la pauvreté rurale, l'insécurité dans certaines régions (Extrême-Nord, Nord-Ouest, Sud-Ouest) et l'attrait des salaires urbains, des milliers de Camerounais rejoignent chaque année Douala, Yaoundé, Bafoussam, Bamenda, Garoua ou Maroua. Le taux d'urbanisation dépasse \num{58}\,\% (ordre de grandeur 2023) et la population urbaine double tous les \num{15} à \num{20} ans.
    \item \textbf{L'absence de planification foncière préventive :} Les documents d'urbanisme (Plans Directeurs d'Aménagement, Plans d'Occupation des Sols) existent souvent sur le papier mais sont rarement respectés ou mis à jour. L'État ne parvient pas à viabiliser les terrains (voirie, eau, électricité, assainissement) \emph{avant} l'installation des populations.
    \item \textbf{La faiblesse des recettes locales :} Les communes (Communes d'Arrondissement, Communautés Urbaines) manquent de ressources fiscales propres pour financer l'entretien des voiries, la collecte des ordures ou l'extension des réseaux d'eau et d'électricité, dont la gestion technique relève souvent d'opérateurs nationaux (Camwater, Eneo) aux logiques centralisées.
\end{itemize}

\begin{attention}[Erreur fréquente : dissocier les problèmes de leur cause structurelle]
Au Bac, évitez de lister l'insalubrité, les embouteillages, le manque d'eau et l'insécurité comme des « malheurs » isolés ou imputables uniquement à l'incivilité des habitants. \textbf{Reliez systématiquement chaque problème à la croissance urbaine rapide et au déficit de planification/gestion.} Par exemple : « L'insalubrité résulte de l'insuffisance des infrastructures de collecte (HYSACAM) face à une production de déchets multipliée par l'explosion démographique. » C'est cette mise en relation cause/effet qui vaut des points.
\end{attention}

\courssub{L'incivisme et l'anarchie urbaine : quand la norme recule}

L'incivisme et l'anarchie sont les deux faces d'une même médaille : l'incapacité de la régulation publique à s'imposer dans l'espace public.

\begin{definition}[Incivisme urbain]
Ensemble des comportements individuels ou collectifs qui contreviennent aux règles de vie collective en ville : dépôts sauvages d'ordures ménagères dans les caniveaux ou sur les trottoirs, dégradation volontaire du mobilier urbain (bancs, lampadaires, panneaux de signalisation), non-respect du code de la route, occupation illicite des emprises routières pour le commerce ambulant.
\end{definition}

\begin{definition}[Anarchie urbaine (ou urbanisation spontanée)]
Mode de production de l'espace urbain caractérisé par l'occupation et la construction sans titre foncier légal, sans permis de construire, en dehors de tout tracé de voirie et de tout réseau d'assainissement. Elle génère des quartiers aux ruelles étroites, impraticables aux engins de secours et de collecte.
\end{definition}

\begin{itemize}
    \item \textbf{Dépôts sauvages et caniveaux bouchés :} À Douala comme à Yaoundé, les caniveaux servent trop souvent de poubelles. Lors des fortes pluies saisonnières, l'eau ne s'écoule plus, provoquant des inondations rapides (rues transformées en torrents aux quartiers Mokolo, Nkolbisson, New-Bell, Akwa).
    \item \textbf{Occupation anarchique du domaine public :} Les trottoirs disparaissent sous les étals des « bayam-sellam » (commerçants ambulants) et les motos-taxis stationnées. La chaussée se rétrécit, piétons et véhicules se partagent un espace conflictuel.
    \item \textbf{Constructions sans permis :} Sur les pentes instables de Yaoundé (collines de Mvog-Ada, Essos) ou dans les zones marécageuses de Douala (Yassa, Logbaba), des milliers de maisons en parpaings ou en banco poussent sans étude de sol, exposant leurs occupants aux glissements de terrain et aux effondrements.
\end{itemize}

\begin{exoresolu}[Analyse d'une photographie aérienne d'un quartier spontané (type « quartier populaire » de Yaoundé ou Douala)]
\textbf{Énoncé :} Vous observez une photographie aérienne verticale d'un quartier périphérique. On y voit : un tissu bâti très dense, aux toits en tôle ondulée de couleurs diverses ; un réseau de ruelles sinueuses, étroites (moins de 2 m), sans hiérarchie visible ; l'absence totale d'espaces verts ou d'équipements collectifs (école, centre de santé) ; des caniveaux à ciel ouvert longeant les ruelles, certains comblés par des déchets ; une érosion visible sur les pentes dénudées.

\textbf{Solution détaillée :}
\begin{enumerate}
    \item \textbf{Identification du type d'espace :} Quartier d'habitat spontané / précaire (bidonville en cours de consolidation).
    \item \textbf{Preuves de l'anarchie urbaine :} Absence de trame orthogonale (pas de plan d'aménagement), ruelles non carrossables (pas de voirie), densité extrême (surpeuplement), non-respect des normes de constructibilité (alignements, recul, surface au sol).
    \item \textbf{Conséquences environnementales :} Ruissellement concentré sur sols nus $\rightarrow$ érosion hydrique (ravines) ; stagnation des eaux usées dans les caniveaux bouchés $\rightarrow$ gîtes larvaires (anophèles/paludisme) et odeurs nauséabondes.
    \item \textbf{Risques sanitaires et sécuritaires :} Inaccessibilité aux camions de pompiers et bennes à ordures ; promiscuité favorisant la propagation des maladies (choléra, COVID-19, tuberculose) ; insécurité liée à l'absence d'éclairage public et de « yeux sur la rue ».
    \item \textbf{Conclusion :} Ce quartier illustre l'échec de la politique foncière préventive. Sa régularisation (opération « ville propre », réaménagement) coûtera infiniment plus cher qu'une viabilisation \emph{amont}.
\end{enumerate}
\end{exoresolu}

\courssub{L'insalubrité : un défi sanitaire majeur}

L'insalubrité est la manifestation la plus visible et la plus dangereuse de la crise urbaine. Elle touche l'eau, l'air, les sols et directement la santé des populations.

\begin{definition}[Insalubrité urbaine]
État de dégradation de l'environnement urbain (air, eau, sols, cadre bâti) résultant de l'accumulation de déchets solides et liquides non collectés, de l'absence d'assainissement, de la pollution industrielle et automobile, créant des conditions favorables au développement de maladies vectorielles et hydriques.
\end{definition}

\begin{itemize}
    \item \textbf{Gestion des déchets solides :} Le système repose largement sur l'entreprise HYSACAM (Hygiène et Salubrité du Cameroun), concessionnaire dans les grandes villes. Mais la collecte ne couvre pas tous les quartiers, surtout les zones non viabilisées inaccessibles aux bennes.
    \item \textbf{Assainissement liquide :} Quasi-absence de réseaux d'égouts collectifs. La majorité des ménages utilise des fosses septiques (souvent non étanches, vidangées manuellement) ou des latrines à ciel ouvert. Les eaux usées ruissellent en surface.
    \item \textbf{Maladies corrélées :} Paludisme (eaux stagnantes = gîtes larvaires), choléra (eau de boisson contaminée par les matières fécales), fièvre typhoïde, diarrhées aiguës, infections respiratoires (brûlage des ordures à ciel ouvert).
\end{itemize}

\begin{exemplebox}[Chiffres clés : la gestion des déchets à Douala]
Douala, capitale économique ($\approx$ 3,5 millions d'habitants), produit environ \SI{2500}{tonnes} de déchets ménagers par jour (estimation). La capacité de collecte effective d'HYSACAM tourne autour de \SI{1000}{tonnes/jour} à \SI{1200}{tonnes/jour} (soit \textbf{moins de la moitié} du volume produit). Le reste finit dans les caniveaux, les cours d'eau (Wouri, Dibamba), les terrains vagues ou est brûlé sur place. Le site d'enfouissement technique de \textbf{Pk 10} (à l'est de la ville) arrive à saturation, posant le problème urgent d'un nouveau site.
\end{exemplebox}

\begin{savaistu}[Choléra et insalubrité : un lien direct]
Le Cameroun fait face à des épidémies de choléra récurrentes (2010-2011, 2014, 2018-2020, 2021-2023). L'épidémie de 2021-2023 a touché plus de \num{15000} cas pour plus de \num{300} décès (chiffres Ministère de la Santé Publique). Les foyers épidémiques se déclarent systématiquement dans les quartiers insalubres des grandes villes (Nouvelle Route Bonabéri, Bepanda à Douala ; Mvog-Ada, Etoudi à Yaoundé) où l'accès à l'eau potable est aléatoire et l'assainissement inexistant. La lutte anti-choléra (chloration de l'eau, sensibilisation, vaccination orale) ne peut être durable sans investissement massif dans l'assainissement structurel.
\end{savaistu}

\courssub{L'insécurité et les risques quotidiens}

L'insécurité urbaine au Cameroun revêt deux dimensions : la \textbf{délinquance} (vols, agressions, « coupeurs de route » urbains) et l'\textbf{insécurité routière/technique} liée à l'état de la voirie.

\begin{itemize}
    \item \textbf{Délinquance juvénile et chômage :} Le taux de chômage des jeunes diplômés (estimé > \num{30}\,\% en milieu urbain) alimente une économie de survie illicite : arrachage de sacs (souvent en moto), cambriolages, cybercriminalité (« arnaqueurs » ou « brouteurs »).
    \item \textbf{Éclairage public défaillant :} De nombreux quartiers, y compris centraux, sont plongés dans le noir la nuit (vol de câbles, factures impayées des communes auprès d'Eneo, matériel vétuste). L'obscurité favorise les agressions.
    \item \textbf{Accidents de la voirie :} L'état des routes (nids-de-poule, absence de marquage au sol, de feux tricolores fonctionnels, de passages piétons) combiné à l'indiscipline routière (excès de vitesse, alcool, surcharge) fait des accidents de la circulation une cause majeure de mortalité et d'invalidité. Douala et Yaoundé concentrent plus de \num{60}\,\% des accidents corporels du pays.
\end{itemize}

\courssub{Le déficit de logement et l'habitat précaire}

Le logement est le premier besoin insatisfait du citadin camerounais. L'offre formelle (promoteurs immobiliers, SIC, MAETUR) est hors de portée de la majorité ; l'offre informelle domine.

\begin{definition}[Habitat précaire]
Logement qui ne répond pas aux normes minimales de salubrité, de sécurité et de durabilité : matériaux de récupération (tôles rouillées, planches, bâches, banco non stabilisé), absence de fondations, surpeuplement (plus de 3 personnes par pièce), absence d'eau courante, d'électricité légale, d'assainissement individuel ou collectif, insécurité foncière (risque d'expulsion).
\end{definition}

\begin{itemize}
    \item \textbf{Cherté des loyers :} À Yaoundé ou Douala, une chambre en dur (parpaings, tôle) dans un quartier périphérique (Nsam, Damas, Nkolbisson, Logbaba) se loue \SI{25000}{FCFA/mois} à \SI{50000}{FCFA/mois} ; un studio « correct » dépasse \SI{100000}{FCFA/mois}. Avec un SMIG à \SI{41875}{FCFA} (depuis 2023), un ménage monoparental consacre \num{60}\,\% à \num{80}\,\% de son revenu au loyer.
    \item \textbf{Surpeuplement :} Dans les quartiers spontanés (Nylon à Douala, Mimboman à Yaoundé), la densité dépasse souvent \num{500} habitants/ha. La promiscuité favorise les violences domestiques, l'échec scolaire (pas d'espace pour étudier) et la propagation des maladies.
    \item \textbf{Opérations « Coup de balai » / « Ville propre » :} Les démolitions périodiques d'habitats précaires (souvent sans relogement décent ni indemnisation préalable) créent un climat d'insécurité foncière et de méfiance envers l'administration.
\end{itemize}

\courssub{La crise des transports urbains : la loi du « benskineur »}

Le transport urbain est le miroir grossissant de la dysfonctionnalité urbaine : distances domicile-travail allongées par l'étalement, voirie saturée, absence de transport en commun de masse.

\begin{itemize}
    \item \textbf{L'hégémonie du transport informel :} Le taxi (peinture jaune à Douala, verte/rouge à Yaoundé) et surtout la \textbf{moto-taxi} (appelée « \emph{benskin} » ou « \emph{benskineur} ») assurent \textbf{plus de \num{80}\,\%} des déplacements motorisés. Le benskin est rapide, porte-à-porte, mais dangereux (accidents graves, absence de casque passager, inhalation de gaz d'échappement) et source de conflits (stationnement anarchique).
    \item \textbf{Insuffisance du transport en commun structuré :} La SOTUC (Société de Transport Urbain du Cameroun) et les bus « Carconfort » / « Transurb » ne couvrent que quelques axes majeurs avec des fréquences aléatoires et une flotte vieillissante. Le projet de Bus à Haut Niveau de Service (BHNS) à Douala (ligne Nord-Sud) et le tramway envisagé à Yaoundé tardent à se concrétiser.
    \item \textbf{Embouteillages chroniques :} Les points de congestion (rond-point Deido, Carrefour Warda, Pont Wouri, Rond-point Nlongkak, Carrefour Mvog-Ada) paralysent la ville aux heures de pointe (6h-9h ; 16h-20h). Temps de trajet multiplié par 3 ou 4. Coût économique énorme (perte d'heures de travail, surconsommation carburant, pollution).
\end{itemize}

\begin{popfigure}
\begin{tikzpicture}
\node[anchor=east, font=\small] at (0,-0.0) {Marche à pied};
\fill[popblue] (0.1,-0.28) rectangle (7.24,0.28);
\node[anchor=west, font=\small\bfseries] at (7.34,-0.0) {42\,\%};
\node[anchor=east, font=\small] at (0,-0.8) {Moto-taxi (benskin)};
\fill[poporange] (0.1,-1.08) rectangle (5.20,-0.52);
\node[anchor=west, font=\small\bfseries] at (5.30,-0.8) {30\,\%};
\node[anchor=east, font=\small] at (0,-1.6) {Taxi collectif};
\fill[popgreen] (0.1,-1.8800000000000001) rectangle (3.16,-1.32);
\node[anchor=west, font=\small\bfseries] at (3.26,-1.6) {18\,\%};
\node[anchor=east, font=\small] at (0,-2.4000000000000004) {Bus et minibus};
\fill[poppurple] (0.1,-2.6800000000000006) rectangle (1.29,-2.12);
\node[anchor=west, font=\small\bfseries] at (1.39,-2.4000000000000004) {7\,\%};
\node[anchor=east, font=\small] at (0,-3.2) {Véhicule personnel};
\fill[poppink] (0.1,-3.4800000000000004) rectangle (0.61,-2.92);
\node[anchor=west, font=\small\bfseries] at (0.71,-3.2) {3\,\%};
\end{tikzpicture}
\legende{Répartition modale estimée des déplacements quotidiens à Douala/Yaoundé (ordre de grandeur basé sur études CODATU/Banque Mondiale). La marche domine, le benskin est le 1er mode motorisé. Le transport en commun structuré reste marginal.}
\end{popfigure}

\courssub{Les difficultés d'approvisionnement en eau et en électricité : le paradoxe des ressources}

Le Cameroun dispose d'un potentiel hydroélectrique énorme (2\textsuperscript{e} d'Afrique après la RDC) et de ressources en eau abondantes. Pourtant, les villes subissent pénurie et délestages.

\begin{itemize}
    \item \textbf{Eau (Camwater / filiales) :} Le réseau de distribution est vétuste (fuites \textgreater \num{30}\,\% de l'eau produite), sous-dimensionné par rapport à la demande, et ne dessert pas les quartiers périphériques. Les coupures sont quotidiennes (rationnement). Les populations recourent aux forages privés (coûteux, qualité non contrôlée), aux bornes-fontaines (files d'attente, coût au jerrican), ou aux vendeurs d'eau en sachets (« pure water ») et bidons.
    \item \textbf{Électricité (Eneo / Sonatrel) :} La capacité installée ($\approx$ \SI{1500}{MW} dont $\approx$ \SI{700}{MW} hydrauliques au fil de l'eau) peine à suivre la pointe de demande ($\approx$ \SI{1300}{MW} et croissante). La saison sèche (janvier-avril) voit le niveau des barrages (Lagdo, Song Loulou, Edéa) baisser $\rightarrow$ \textbf{délestages tournants} programmés ou imprévus, parfois \num{8} à \num{12} h/jour. Coût pour les entreprises (groupes électrogènes) et les ménages (appareils endommagés, insécurité nocturne).
    \item \textbf{Coût élevé :} Les tarifs sociaux existent mais les compteurs prépayés (Cash Power) et les factures estimées pèsent lourd sur les budgets modestes. Le raccordement initial reste un obstacle financier majeur (poteau, câble, compteur à la charge de l'abonné).
\end{itemize}

\courssub{Synthèse : classement des problèmes urbains par catégorie}

Pour structurer votre analyse (dissertation, commentaire de document), classez les problèmes en quatre dimensions interdépendantes :

\begin{center}
\begin{tabularx}{\linewidth}{|l|X|X|X|X|}
\hline
\rowcolor{popblueL}
\textbf{Dimension} & \textbf{Problèmes principaux} & \textbf{Causes structurelles} & \textbf{Acteurs concernés} & \textbf{Conséquences sur l'intégration} \\
\hline
\textbf{Sociale} & Habitat précaire, surpeuplement, exclusion des pauvres, difficultés d'accès aux services (santé, éducation), ségrégation spatiale & Pauvreté, chômage, coût du foncier, absence de politique sociale du logement & Maires, MINHDU, Promoteurs, ONG & Fragmentation spatiale (ville riche / ville pauvre), tensions identitaires, échec du brassage cosmopolite \\
\hline
\textbf{Environnementale} & Insalubrité (déchets, eaux usées), inondations, érosion, pollution air/bruit, perte biodiversité urbaine & Absence assainissement, collecte défaillante, occupation zones à risque (lits de rivière, pentes), brûlage déchets & HYSACAM, MINEPDED, Communes, Population & Risques sanitaires (choléra, paludisme), dégradation cadre de vie, changement climatique local \\
\hline
\textbf{Économique} & Embouteillages (perte temps), coût transport/logement, délestages (perte production), secteur informel dominant & Voirie insuffisante, absence TCSP, réseau électrique saturé, fiscalité locale faible & MINTRANSPORT, Eneo, Camwater, Entreprises, MINEPAT & Perte compétitivité villes, frein investissement, précarisation emplois \\
\hline
\textbf{Sécuritaire} & Délinquance, accidents routiers, risques industriels/technologiques, insécurité foncière (expulsions) & Chômage jeunes, éclairage défaillant, voirie dangereuse, impunité, non-respect normes construction & MINAT, DGSN, Maires, Justice & Peur, repli communautaire, perte confiance institutions, instabilité politique potentielle \\
\hline
\end{tabularx}
\end{center}

\begin{methode}[Inventorier et classer les problèmes urbains à partir de photographies (TP / Épreuve pratique)]
\textbf{Objectif :} Transformer une observation visuelle en analyse géographique structurée.
\begin{enumerate}
    \item \textbf{Observation systématique (lecture d'image) :} Balayez la photo par plans (premier plan, plan moyen, arrière-plan). Notez \textbf{tout} ce qui est visible : type d'habitat (matériaux, état), voirie (revêtement, largeur, état), présence/absence d'équipements (éclairage, bornes d'eau, poubelles, écoles), activités (commerce, artisanat, jeux), végétation, déchets, eau stagnante, véhicules.
    \item \textbf{Inventaire codé :} Créez un tableau à deux colonnes : \textit{Élément observé} / \textit{Problème géographique associé}. Exemple : « Ruelles en terre battue, ornières, flaques » $\rightarrow$ « Absence de voirie / drainage $\rightarrow$ Insalubrité, inaccessibilité, érosion ».
    \item \textbf{Classification thématique :} Regroupez les problèmes identifiés dans les 4 catégories du tableau ci-dessus (Social, Environnemental, Économique, Sécuritaire). Un même élément peut relever de plusieurs catégories (ex: déchets = Environnemental + Sanitaire/Social + Économique/coût santé).
    \item \textbf{Hiérarchisation et argumentation :} Identifiez le \textbf{problème central} (souvent l'absence de planification/viabilisation \emph{amont}) dont découlent les autres. Rédigez un paragraphe synthétique : « La photographie du quartier X illustre l'urbanisation spontanée : l'habitat précaire (Social) s'installe sur zone marécageuse sans drainage (Environnemental), loin des emplois (Économique), dans l'insécurité foncière et physique (Sécuritaire). »
    \item \textbf{Proposition d'aménagements (croquis de réaménagement) :} Sur un calque ou un croquis à part, tracez la voirie structurante, les espaces de détente, les emplacements d'équipements, les zones de densification contrôlée. Légendez.
\end{enumerate}
\end{methode}

\begin{aretenir}[Les problèmes des villes camerounaises : l'essentiel pour le Bac]
\begin{itemize}
    \item \textbf{Cause racine :} Croissance urbaine rapide (exode rural, natalité) + Planification défaillante + Faiblesse financière des communes.
    \item \textbf{4 piliers des dysfonctionnements :} Insalubrité (déchets, assainissement, choléra) ; Habitat précaire (surpeuplement, cherté loyers, insécurité foncière) ; Transport (hégémonie benskin, embouteillages, absence TCSP) ; Énergie/Eau (délestages Eneo, coupures Camwater, réseau vétuste).
    \item \textbf{Incivisme / Anarchie :} Symptômes et amplificateurs (dépôts sauvages, occupation trottoirs, constructions sans permis).
    \item \textbf{Mots-clés à maîtriser :} \cle{Urbanisation spontanée}, \cle{Habitat précaire}, \cle{Insalubrité}, \cle{Benskin}, \cle{Délestage}, \cle{HYSACAM}, \cle{Camwater}, \cle{Eneo}, \cle{Viabilisation}, \cle{TCSP (Transport Collectif en Site Propre)}, \cle{Quartier cosmopolite}.
    \item \textbf{Compétence clé :} Savoir passer de la description d'une photo/ carte à l'explication structurelle (lien cause/effet) et à la proposition d'aménagement (croquis de réaménagement).
\end{itemize}
\end{aretenir}

\coursec{Solutions et perspectives pour les villes camerounaises}

Dans la section précédente, nous avons dressé un constat sévère : croissance démographique galopante, habitat précaire, insalubrité, embouteillages, insécurité, déficits d'eau et d'électricité... Les villes camerounaises souffrent de maux nombreux et imbriqués. Mais un diagnostic, même lucide, ne suffit pas : la géographie est aussi une science de l'action. Cette section répond donc à une question simple et exigeante : \emph{que faire, concrètement, pour rendre nos villes plus vivables ?} Nous verrons que les solutions existent — aménagement, assainissement, voirie, approvisionnement, sécurité, éducation civique, décentralisation — à condition d'être réalistes, coordonnées et portées par tous : État, communes, entreprises et habitants.

\courssub{L'aménagement urbain : organiser la croissance de la ville}

\textbf{Intuition.} Imagine un marché qui s'étend chaque semaine sans aucun plan : les étals envahissent la route, les camions ne passent plus, les bagarres éclatent pour les meilleures places. C'est exactement ce qui arrive à une ville qui grandit sans règles. L'aménagement urbain, c'est le « règlement intérieur » de la ville, pensé à l'avance.

\begin{definition}[Aménagement urbain et assainissement]
L'\cle{aménagement urbain} est l'ensemble des actions de planification, d'organisation et d'équipement de l'espace urbain (plans d'urbanisme, lotissements, zonage, construction d'infrastructures) visant à rendre la ville fonctionnelle, harmonieuse et vivable.
L'\cle{assainissement} désigne l'ensemble des techniques et des services destinés à évacuer et à traiter les déchets (solides et liquides) et les eaux usées, afin de préserver la salubrité et la santé publique en milieu urbain.
\end{definition}

Concrètement, l'aménagement urbain au Cameroun repose sur trois outils complémentaires :

\begin{itemize}
\item \textbf{Les plans d'urbanisme} (plans directeurs, plans de zonage) : documents officiels qui fixent à l'avance la vocation de chaque espace (habitat, industrie, équipements publics, espaces verts). Douala et Yaoundé disposent de plans directeurs, mais leur application reste inégale, faute de moyens de contrôle et de pression foncière.
\item \textbf{Les lotissements} : découpage de terrains en parcelles viabilisées (desservies par des voies, l'eau et l'électricité) pour encadrer l'habitat. Des opérations comme celles de la \textbf{SIC} (Société immobilière du Cameroun) ou de \textbf{MAETUR} (Mission d'aménagement et d'équipement des terrains urbains et ruraux) ont produit des quartiers planifiés (Cité Verte et quartier du Lac à Yaoundé, Bonamoussadi à Douala).
\item \textbf{La réhabilitation des quartiers précaires} : plutôt que de démolir systématiquement (ce qui déplace le problème), on améliore sur place les quartiers spontanés — ouverture de voies, adduction d'eau, latrines, drainage. C'est la logique des projets de « restructuration » menés dans des quartiers comme Nylon à Douala ou Mvog-Ada à Yaoundé.
\end{itemize}

\begin{attention}[Piège fréquent : les solutions « de papier »]
Au Baccalauréat, beaucoup de copies proposent des solutions irréalistes : « raser les bidonvilles », « construire des immeubles pour tous », « déplacer la capitale ». Une bonne solution tient compte des \cle{moyens réels des communes} (budgets limités), du \cle{coût foncier}, des \cle{droits des habitants} et du temps nécessaire. Une réhabilitation progressive avec participation des populations est souvent plus efficace qu'une démolition brutale, qui jette les familles à la rue et recrée le précaire ailleurs.
\end{attention}

\courssub{L'assainissement : la bataille quotidienne contre les déchets}

\textbf{Observation.} Un matin de pluie à Douala : les caniveaux débordent, les ordures flottent, l'eau stagne dans les rues de Deïdo. Quelques heures plus tard, les camions orange de l'\textbf{Hysacam} (Hygiène et salubrité du Cameroun) passent dans les artères. Cette scène résume l'enjeu : une ville produit chaque jour des montagnes de déchets ; sans collecte organisée, elle s'étouffe.

\begin{exemplebox}[Hysacam : des milliers de tonnes par jour]
La société \textbf{Hysacam}, en partenariat avec les communautés urbaines, assure la collecte des ordures ménagères dans les grandes villes. À Douala et Yaoundé réunies, ce sont \textbf{plusieurs milliers de tonnes de déchets par jour} qui sont ramassées (ordre de grandeur : environ \num{2000} à \num{3000} tonnes/jour pour les deux métropoles ; les chiffres varient selon les années et les sources). Malgré cela, une part importante des déchets échappe encore à la collecte, surtout dans les quartiers spontanés non desservis par la voirie carrossable.
\end{exemplebox}

La gestion des déchets urbains fonctionne comme une \cle{chaîne} : si un maillon est faible, tout le système s'effondre.

\begin{popfigure}
\begin{center}
\begin{tikzpicture}[
box/.style={draw=popblue, fill=popblueL, rounded corners=3pt, minimum width=2.7cm, minimum height=1.4cm, align=center, font=\small, text=popdark},
arr/.style={-{Stealth[length=3mm]}, very thick, poporange},
node distance=0.5cm]
\node[box, align=center] (prod){\textbf{1. Production}\\ ménages, marchés,\\ industries};
\node[box, right=of prod, align=center] (tri){\textbf{2. Précollecte}\\ poubelles, points\\ d'apport, tri};
\node[box, right=of tri, align=center] (coll){\textbf{3. Collecte}\\ camions (Hysacam),\\ bennes};
\node[box, right=of coll, align=center] (trait){\textbf{4. Traitement}\\ décharges contrôlées,\\ compostage};
\node[box, right=of trait, align=center] (valo){\textbf{5. Valorisation}\\ recyclage, biogaz,\\ compost};
\draw[arr] (prod) -- (tri);
\draw[arr] (tri) -- (coll);
\draw[arr] (coll) -- (trait);
\draw[arr] (trait) -- (valo);
\draw[arr, popgreen, dashed] (valo.south west) .. controls +(0,-0.8) and +(0,-1.5) .. node[below, font=\scriptsize, text=popgreen]{retour à l'économie} (prod.south);
\end{tikzpicture}
\end{center}
\legende{Schéma de la chaîne de gestion des déchets urbains au Cameroun. 1 : production (les ménages et marchés sont les premiers producteurs) ; 2 : précollecte (souvent assurée par des associations de jeunes à vélo ou à charrette dans les quartiers sans voirie) ; 3 : collecte mécanisée (Hysacam) ; 4 : traitement (décharges, encore trop souvent sauvages) ; 5 : valorisation (recyclage du plastique et des métaux, compostage des déchets organiques), maillon encore embryonnaire mais prometteur d'emplois.}
\end{popfigure}

Au-delà des ordures, l'assainissement comprend le \textbf{curage régulier des caniveaux} (indispensable avant la saison des pluies pour éviter les inondations), les \textbf{campagnes de salubrité} (opérations « ville propre », journées de nettoyage citoyen) et la construction de \textbf{réseaux d'eaux usées et de latrines améliorées}, car une grande partie des villes camerounaises n'est pas raccordée à un égout collectif.

\begin{exemplebox}[Avant / après : l'assainissement dans une ville camerounaise]
\begin{center}
\begin{tabularx}{\linewidth}{|l|X|X|}
\hline
\rowcolor{popblueL} \textbf{Indicateur} & \textbf{Avant les actions} & \textbf{Après les actions} \\
\hline
Collecte des ordures & Irrégulière, dépôts sauvages dans les rues et caniveaux & Passages réguliers des camions, bennes dans les marchés \\
\hline
Caniveaux & Colmatés, inondations à chaque pluie & Curés avant la saison des pluies, écoulement amélioré \\
\hline
Santé & Recrudescence du paludisme et du choléra & Baisse des épidémies et des eaux stagnantes \\
\hline
Cadre de vie & Odeurs, mouches, image dégradée du quartier & Rues propres, attractivité commerciale retrouvée \\
\hline
\end{tabularx}
\end{center}
\emph{Lecture :} ce tableau comparatif, construit à partir de l'observation d'un quartier (par exemple autour du marché Mokolo à Yaoundé), montre que l'assainissement produit des effets en cascade : propreté, santé, économie. À l'examen, sache construire un tel tableau à deux colonnes à partir d'un document photographique.
\end{exemplebox}

\courssub{L'amélioration de la voirie urbaine}

La \cle{voirie urbaine} désigne l'ensemble des voies de circulation d'une ville (rues, avenues, boulevards, ponts, échangeurs) ainsi que leurs équipements annexes (trottoirs, caniveaux, signalisation, éclairage). Une voirie de qualité fluidifie le trafic, ouvre les quartiers à la collecte des déchets et aux secours, et réduit le coût des transports.

Les actions menées au Cameroun sont de deux ordres :

\begin{itemize}
\item \textbf{La construction et la modernisation} : bitumage de rues secondaires, élargissement d'artères, construction de ponts et d'échangeurs. L'\textbf{échangeur de la gare à Yaoundé} (inauguré en 2017) illustre cette logique : en dénivelant le carrefour, il a considérablement fluidifié un nœud de circulation autrefois paralysé aux heures de pointe. À Douala, le second pont sur le Wouri (inauguré 2017) et la réhabilitation du pont ancien jouent le même rôle de désenclavement de la rive droite (Bonabéri).
\item \textbf{L'entretien} : reprofilage, colmatage des nids-de-poule, curage des buses. Une route non entretenue se dégrade vite sous l'effet des pluies et des charges lourdes ; l'entretien régulier coûte bien moins cher que la reconstruction.
\end{itemize}

\begin{attention}[Ne pas confondre construction et entretien]
Dans les copies, on lit souvent : « il faut construire plus de routes ». C'est insuffisant : sans \textbf{entretien}, une route neuve devient impraticable en quelques années. La bonne réponse articule toujours construction \emph{et} maintenance, et précise qui paie (communes, État, redevances).
\end{attention}

\courssub{Renforcer l'approvisionnement en eau et en électricité}

\textbf{Intuition.} Une ville sans eau courante oblige des milliers d'enfants à parcourir chaque matin des kilomètres avec des bidons : c'est du temps scolaire perdu et un frein à l'hygiène. L'accès à l'eau et à l'électricité est donc une condition de tous les autres progrès urbains.

\begin{itemize}
\item \textbf{Pour l'eau} : forages équipés de bornes-fontaines dans les quartiers non raccordés, extension des réseaux de la \textbf{Camwater} (Cameroun Water Utilities), réhabilitation des stations de pompage (par exemple celles alimentant Douala à partir du fleuve Wouri et de la nappe de la Japoma), lutte contre les fuites et les branchements frauduleux qui font perdre une part importante de l'eau produite (taux de rendement du réseau perfectible).
\item \textbf{Pour l'électricité} : extension et densification des réseaux d'\textbf{Eneo}, renforcement de la production (barrages hydroélectriques de Songloulou et d'Édéa sur la Sanaga, centrales thermiques, projet de Nachtigal sur la Sanaga, d'une capacité de \num{420} MW, mis en service progressif depuis 2023-2024, ordre de grandeur à retenir), et développement de solutions solaires pour les quartiers excentrés.
\end{itemize}

\courssub{La sécurisation des villes}

L'insécurité (agressions, cambriolages, banditisme nocturne) se combat par une combinaison de moyens :

\begin{itemize}
\item \textbf{L'éclairage public} : une rue éclairée décourage les agresseurs et redonne vie au quartier le soir (petits commerces, déplacements). Les programmes d'installation de lampadaires, parfois solaires, transforment la vie nocturne des villes.
\item \textbf{La police de proximité} : postes de quartier, patrouilles à pied, qui rétablissent un lien de confiance entre forces de l'ordre et habitants.
\item \textbf{Les comités de vigilance} : habitants organisés qui surveillent leur quartier et collaborent avec la police et la gendarmerie. Leur efficacité dépend d'un encadrement strict pour éviter les dérives (justice populaire).
\end{itemize}

\courssub{L'éducation civique : la solution la moins chère et la plus difficile}

Beaucoup de problèmes urbains naissent de comportements : ordures jetées dans les caniveaux, constructions anarchiques sur les voies, dégradation des équipements publics, corruption aux contrôles. Aucun camion ne suffira si chacun jette ses déchets par la fenêtre. L'\cle{éducation civique} vise donc à transformer les habitants en acteurs de leur ville : campagnes de sensibilisation (radios, écoles, églises et mosquées), éducation à l'environnement dès le primaire, implication des populations dans la gestion du quartier (comités de quartier, associations de précollecte), et sanctions effectives contre l'incivisme. Les \textbf{communes}, au plus près des habitants, sont le meilleur relais de cette pédagogie. C'est aussi le ciment des \textbf{quartiers cosmopolites} : le respect des règles communes permet le « vivre-ensemble » entre populations d'origines, de cultures et de religions diverses.

\begin{savaistu}[Les communautés urbaines : des compétences propres]
Au Cameroun, les grandes villes sont organisées en \textbf{communautés urbaines} (Communauté urbaine de Douala, de Yaoundé, de Bafoussam, de Bamenda, de Garoua, de Maroua...), qui regroupent plusieurs communes d'arrondissement. La loi sur la décentralisation leur confie des compétences propres en matière de \textbf{voirie urbaine}, d'\textbf{hygiène et salubrité}, d'éclairage public. Elles perçoivent pour cela des recettes propres (taxes, redevances, centimes additionnels communaux). C'est donc souvent la communauté urbaine — et non l'État central — qui contracte avec Hysacam ou programme le bitumage d'une rue : un détail qui peut valoriser ta copie au Bac.
\end{savaistu}

\courssub{Décentralisation et désenclavement : traiter le problème à la racine}

Toutes les solutions précédentes soignent les symptômes \emph{dans} les villes. Mais pourquoi Douala et Yaoundé grossissent-elles si vite ? Parce que l'exode rural et les migrations alimentent en permanence leur croissance : on y vient chercher emplois, universités, hôpitaux, marchés. La solution de fond consiste donc à \textbf{rééquilibrer le territoire national} :

\begin{itemize}
\item \textbf{La décentralisation} : transférer des compétences et des ressources aux régions et communes (régionalisation engagée depuis la loi de 1996, accélérée avec l'élection des conseils régionaux en décembre 2020) pour que les services publics de qualité ne soient plus concentrés dans les deux métropoles.
\item \textbf{Le désenclavement régional} : routes et chemins de fer reliant les régions (axe Yaoundé--Ngaoundéré, route Bamenda--Enugu, corridor Douala--N'Djamena), qui permettent aux villes secondaires (Bafoussam, Ngaoundéré, Maroua, Bertoua, Kribi) de devenir des pôles attractifs.
\item \textbf{Le développement des villes moyennes} : y installer universités, hôpitaux, zones industrielles (comme le port en eau profonde de Kribi et sa zone industrielle) pour qu'elles retiennent leurs jeunes au lieu de les pousser vers Douala et Yaoundé.
\end{itemize}

\begin{methode}[Rédiger un paragraphe argumenté proposant des solutions à un problème urbain]
\begin{enumerate}
\item \textbf{Rappeler le problème} en une phrase précise et localisée (ex. : « À Douala, l'insalubrité provoque inondations et épidémies »).
\item \textbf{Annoncer la solution principale} avec son acteur (ex. : « la communauté urbaine peut renforcer la collecte des déchets »).
\item \textbf{Expliquer le mécanisme} : comment la solution agit sur le problème (curage des caniveaux $\rightarrow$ écoulement des eaux $\rightarrow$ fin des inondations).
\item \textbf{Illustrer par un exemple camerounais chiffré} (Hysacam, échangeur de la gare, forages Camwater...).
\item \textbf{Nuancer} : citer une limite (coût, entretien, comportements) et une condition de réussite (participation des habitants, financement communal).
\end{enumerate}
Un bon paragraphe argumenté fait donc toujours : problème $\rightarrow$ solution $\rightarrow$ mécanisme $\rightarrow$ exemple $\rightarrow$ limite.
\end{methode}

\begin{exoresolu}[Paragraphe argumenté : lutter contre l'insalubrité à Yaoundé]
\textbf{Énoncé.} En un paragraphe argumenté de 10 à 15 lignes, propose des solutions réalistes pour lutter contre l'insalubrité dans la ville de Yaoundé.
\tcblower
\textbf{Solution rédigée.} Yaoundé, capitale politique du Cameroun, souffre d'une insalubrité chronique : dépôts sauvages d'ordures, caniveaux colmatés provoquant des inondations dans les quartiers de Mokolo ou Mvog-Ada, et risques épidémiques (choléra, paludisme). Pour y remédier, la Communauté urbaine de Yaoundé peut d'abord renforcer la collecte des déchets en augmentant la fréquence des passages d'Hysacam, qui ramasse déjà plusieurs centaines de tonnes par jour dans la ville, et en installant des bennes d'apport près des marchés. Ensuite, le curage systématique des caniveaux avant la saison des pluies permettrait l'écoulement des eaux et réduirait les inondations. Enfin, comme aucun service ne suffit face à l'incivisme, des campagnes d'éducation civique dans les écoles et les médias, relayées par les comités de quartier, doivent convaincre les habitants de ne plus jeter les ordures dans les rigoles. Ces solutions restent réalistes car elles s'appuient sur des structures existantes (communauté urbaine, Hysacam, comités de quartier) ; leur limite est financière, d'où la nécessité que les habitants paient effectivement la redevance d'enlèvement des ordures ménagères.
\end{exoresolu}

\begin{aretenir}[L'essentiel de la section]
\begin{itemize}
\item L'\textbf{aménagement urbain} (plans d'urbanisme, lotissements, réhabilitation des quartiers précaires) organise la croissance de la ville au lieu de la subir.
\item L'\textbf{assainissement} repose sur une chaîne : production, précollecte, collecte (Hysacam), traitement, valorisation ; le curage des caniveaux et les campagnes de salubrité le complètent.
\item La \textbf{voirie} exige construction \emph{et} entretien (ex. : échangeur de la gare à Yaoundé, second pont sur le Wouri à Douala).
\item L'approvisionnement en \textbf{eau} (forages, Camwater) et en \textbf{électricité} (Eneo, Nachtigal \num{420} MW) conditionne tous les autres progrès.
\item La \textbf{sécurité} combine éclairage public, police de proximité et comités de vigilance encadrés.
\item L'\textbf{éducation civique} et la participation des communes rendent les solutions durables et consolident le vivre-ensemble dans les quartiers cosmopolites.
\item La \textbf{décentralisation} (conseils régionaux 2020) et le \textbf{désenclavement régional} freinent l'exode vers Douala et Yaoundé en développant les villes moyennes.
\item Au Bac : propose toujours des solutions \textbf{réalistes}, reliées à un acteur, un mécanisme et un exemple camerounais.
\end{itemize}
\end{aretenir}

\begin{exercice}[Pour t'entraîner]
\begin{enumerate}
\item Définis les termes suivants : aménagement urbain ; assainissement ; voirie urbaine ; lotissement.
\item Construis un tableau à deux colonnes (« problème » / « solution adaptée ») pour cinq problèmes urbains étudiés dans la section précédente.
\item Rédige un paragraphe argumenté de 12 lignes maximum montrant comment la décentralisation peut freiner la croissance excessive de Douala et Yaoundé.
\item À partir du schéma de la chaîne de gestion des déchets, explique pourquoi la faiblesse du maillon « précollecte » dans les quartiers sans voirie compromet toute la chaîne.
\end{enumerate}
\end{exercice}

\begin{corrige}[Exercice]
\textbf{1.} Aménagement urbain : ensemble des actions de planification et d'équipement de l'espace urbain pour le rendre fonctionnel et vivable. Assainissement : ensemble des techniques d'évacuation et de traitement des déchets et des eaux usées pour préserver la salubrité. Voirie urbaine : ensemble des voies de circulation d'une ville et de leurs équipements annexes. Lotissement : découpage d'un terrain en parcelles viabilisées destinées à la construction.

\textbf{2.} Exemple de tableau attendu : insalubrité $\rightarrow$ collecte Hysacam + curage des caniveaux ; embouteillages $\rightarrow$ échangeurs et entretien de la voirie ; insécurité $\rightarrow$ éclairage public et police de proximité ; pénurie d'eau $\rightarrow$ forages et extension des réseaux Camwater ; habitat précaire $\rightarrow$ lotissements et réhabilitation des quartiers.

\textbf{3.} Le paragraphe doit suivre la méthode : problème (concentration des activités et migrations vers les deux métropoles), solution (décentralisation et développement des villes moyennes), mécanisme (transfert de services et ressources vers les régions $\rightarrow$ attractivité des villes secondaires $\rightarrow$ ralentissement de l'exode), exemple (conseils régionaux élus en 2020, pôle de Kribi), limite (lenteur des transferts de moyens financiers).

\textbf{4.} La précollecte rassemble les déchets des habitations vers les points d'apport. Dans les quartiers sans voirie carrossable, les camions ne passent pas : sans précollecte (charrettes, tricycles associatifs), les ordures restent dans les cours et les caniveaux, la collecte mécanisée devient impossible en aval, et les déchets finissent en dépôts sauvages. Un maillon faible bloque donc toute la chaîne.
\end{corrige}

\coursec{Travaux pratiques : représenter l'urbanisation du Cameroun}

Après avoir analysé les \textbf{solutions et perspectives} pour faire face aux défis urbains (décentralisation, planification, assainissement, habitat), il est indispensable de passer de l'analyse théorique à la \textbf{pratique cartographique et statistique}. En géographie, «~faire parler les chiffres et les cartes~» est une compétence centrale évaluée au Baccalauréat. Cette section te guide pas à pas dans la réalisation de quatre productions types : une carte choroplèthe/proportionnelle, un diagramme en barres, un croquis de quartier et une interprétation statistique. Chaque TP suit une \textbf{méthode rigoureuse} (cadre, titre, légende, échelle, orientation) que tu dois maîtriser parfaitement.

\courssub{TP 1 : Légender une carte muette du Cameroun — Villes et métropoles}

\begin{exemplebox}[Objectif]
Savoir localiser les 10 chefs-lieux de région et les 2 métropoles (Douala, Yaoundé) sur un fond de carte régionalisé, en utilisant une \textbf{symbolique proportionnelle} (cercles ou carrés dont la surface est proportionnelle à la population) et une \textbf{hiérarchisation des symboles} (métropoles, chefs-lieux de région, autres villes).
\end{exemplebox}

\begin{methode}[Étapes pour légender une carte de peuplement urbain]
\begin{enumerate}
    \item \textbf{Analyser le fond de carte :} identifier le contour national, les 10 régions, le tracé des grands axes routiers (N1, N2, N3, N10) et le réseau hydrographique de repère (Sanaga, Bénoué, Logone).
    \item \textbf{Choisir la symbolique :} des \textbf{cercles proportionnels} (surface $\propto$ population) sont la norme pour les villes. Définir 3 à 4 classes de population (ex. $> 2\,000\,000$ ; $500\,000 - 2\,000\,000$ ; $100\,000 - 500\,000$ ; $< 100\,000$).
    \item \textbf{Calculer les rayons :} si $S = \pi r^2 \propto P$, alors $r \propto \sqrt{P}$. Fixer un rayon max (ex. 5 mm pour Douala/Yaoundé $\approx$ 3,8 M hab.) et en déduire les autres.
    \item \textbf{Positionner les symboles :} au centre du chef-lieu (point précis), sans masquer le nom de la ville qui sera écrit \emph{à côté} du symbole (en général à droite ou en bas à droite).
    \item \textbf{Construire la légende :} encadrée, en bas à droite ou gauche de la carte. Elle contient : le titre de la carte, l'échelle graphique, la flèche Nord, la source, et la \textbf{boîte de symboles} (cercles de tailles exactes avec leur valeur en habitants).
    \item \textbf{Soigner l'esthétique :} traits fins pour les régions, traits plus épais pour la frontière nationale, couleur unique pour les symboles (ex. rouge ou noir), écriture lisible (stylo feutre noir).
\end{enumerate}
\end{methode}

\begin{popfigure}
\begin{center}
\includegraphics[width=\linewidth,height=11cm,keepaspectratio]{N03/cmr_reg_muette.jpg}
\end{center}
\legende{Carte muette du Cameroun (limites des régions en trait épais, départements en trait fin) servant de base au TP 1. Le modèle attendu y place des symboles proportionnels : deux grands cercles pour les métropoles (Douala, Yaoundé) et dix cercles plus petits pour les chefs-lieux de région, avec une légende dont les rayons sont identiques à ceux de la carte. \\ Source du fond de carte : Wikimedia Commons, Flappiefh, CC BY-SA 4.0.}
\end{popfigure}

\begin{exoresolu}[TP 1 : Production attendue pour la carte]
\textbf{Énoncé :} À l'aide du fond de carte ci-dessus et du tableau de population ci-dessous, réaliser la carte des principales villes du Cameroun avec symboles proportionnels.

\begin{center}
\begin{tabularx}{\linewidth}{|l|c|X|}\hline
\textbf{Ville} & \textbf{Région} & \textbf{Population estimée (2023, agglomération)} \\ \hline
Douala & Littoral & 3\,800\,000 \\ \hline
Yaoundé & Centre & 3\,800\,000 \\ \hline
Bamenda & Nord-Ouest & 600\,000 \\ \hline
Bafoussam & Ouest & 500\,000 \\ \hline
Garoua & Nord & 500\,000 \\ \hline
Maroua & Extrême-Nord & 400\,000 \\ \hline
Ngaoundéré & Adamaoua & 300\,000 \\ \hline
Bertoua & Est & 150\,000 \\ \hline
Ebolowa & Sud & 150\,000 \\ \hline
Buéa & Sud-Ouest & 150\,000 \\ \hline
\end{tabularx}
\end{center}

\tcblower

\textbf{Solution détaillée (étapes de correction) :}
\begin{enumerate}
    \item \textbf{Calcul des rayons} (échelle : $r = 0,18 \times \sqrt{P_{\text{millions}}}$ cm) :
    \begin{itemize}
        \item Douala/Yaoundé (3,8 M) : $r = 0,18 \times \sqrt{3,8} \approx 0,35$ cm.
        \item Bamenda (0,6 M) : $r \approx 0,14$ cm.
        \item Bafoussam/Garoua (0,5 M) : $r \approx 0,13$ cm.
        \item Maroua (0,4 M) : $r \approx 0,11$ cm.
        \item Ngaoundéré (0,3 M) : $r \approx 0,10$ cm.
        \item Bertoua/Ebolowa/Buéa (0,15 M) : $r \approx 0,07$ cm.
    \end{itemize}
    \item \textbf{Traçage :} Dessiner les cercles au compas (ou gabarit) aux coordonnées précises des villes. Douala et Yaoundé en \textbf{rouge} (métropoles), les 4 grands chefs-lieux en \textbf{bleu}, les 3 petits en \textbf{vert}.
    \item \textbf{Nommage :} Écrire le nom de la ville en minuscules (sauf initiale) en police droite, horizontalement, à droite du symbole, sans le toucher.
    \item \textbf{Légende finale :} Reporter exactement les 3 cercles de la légende avec leurs valeurs (ex. «~$\bullet$ = 3,8 M hab.~»), l'échelle graphique (0-200 km), la flèche Nord, le titre \textbf{«~Cameroun : Répartition des principales villes (2023)~»} et la source.
\end{enumerate}
\end{exoresolu}

\courssub{TP 2 : Construire un diagramme en barres de la population des cinq plus grandes villes}

\begin{exemplebox}[Objectif]
Transformer un tableau statistique en \textbf{diagramme en barres verticales} (histogramme à barres disjointes) pour comparer visuellement le poids démographique des 5 premières villes. Le diagramme doit être rigoureux : axes gradués, unités, titre, source.
\end{exemplebox}

\begin{methode}[Construction d'un diagramme en barres à partir d'un tableau]
\begin{enumerate}
    \item \textbf{Trier les données :} classer les villes par population décroissante (Douala, Yaoundé, Bamenda, Bafoussam, Garoua).
    \item \textbf{Choisir l'échelle de l'axe des ordonnées (Y) :} valeur max $\approx$ 3,8 M. Prendre un arrondi supérieur commode : 4 000 000 (4 M). Graduation tous les 500 000 ou 1 000 000.
    \item \textbf{Tracer le cadre :} abscisse (X) = noms des villes (largeur égale, espacement régulier). Ordonnée (Y) = Population (en millions d'habitants).
    \item \textbf{Dessiner les barres :} largeur constante (ex. 1 cm), espacement constant (ex. 0,5 cm). Hauteur proportionnelle à la valeur exacte (ex. 3,8 M = 3.8cm si 1 cm = 1 M).
    \item \textbf{Habiller :} Titre centré au-dessus («~Population des 5 plus grandes villes du Cameroun en 2023~»), légende d'axe Y («~Population (en millions d'habitants)~»), source sous le graphique, valeurs au sommet de chaque barre.
    \item \textbf{Colorier :} Deux tons pour distinguer métropoles (rouge) et villes secondaires (bleu).
\end{enumerate}
\end{methode}

\begin{popfigure}
\begin{tikzpicture}
\begin{axis}[
    ybar,
    bar width=12pt,
    width=\linewidth,
    height=9cm,
    enlargelimits=0.15,
    ylabel={Population (en millions d'habitants)},
    ylabel style={font=\small},
    symbolic x coords={Douala, Yaoundé, Bamenda, Bafoussam, Garoua},
    xtick=data,
    xticklabel style={font=\small, rotate=15, anchor=east},
    ymin=0, ymax=4.2,
    ytick={0,1,2,3,4},
    yticklabel style={font=\small},
    nodes near coords,
    nodes near coords style={font=\small, anchor=south, yshift=2pt},
    title={Population des 5 plus grandes villes du Cameroun (2023)},
    title style={font=\normalsize\bfseries, yshift=0.3cm},
    grid=major,
    grid style={popline, dashed},
    axis lines=left,
    axis line style={popdark},
    tick style={popdark},
    legend style={at={(0.5,-0.15)}, anchor=north, legend columns=-1, font=\tiny},
]
% Métropoles (Rouge)
\addplot[popred!70, fill=popred!50] coordinates {(Douala,3.8) (Yaoundé,3.8)};
% Villes secondaires (Bleu)
\addplot[popblue!70, fill=popblue!50] coordinates {(Bamenda,0.6) (Bafoussam,0.5) (Garoua,0.5)};
\legend{Métropoles, Villes secondaires}
\end{axis}
\end{tikzpicture}
\legende{Diagramme en barres verticales comparant la population des 5 principales villes camerounaises en 2023. Douala et Yaoundé (métropoles, en rouge) dominent largement (environ 3,8 M hab. chacune), suivies de loin par Bamenda, Bafoussam et Garoua (environ 0,5-0,6 M hab., en bleu). Source : BUCREP/INS, estimations.}
\end{popfigure}

\begin{exoresolu}[TP 2 : Commentaire du diagramme réalisé]
\textbf{Énoncé :} À partir du diagramme ci-dessus, rédiger un commentaire structuré en 3 phrases maximum.

\tcblower

\textbf{Solution détaillée :}
\begin{enumerate}
    \item \textbf{Constat global :} Le diagramme met en évidence une \textbf{primatie macrocéphalique} extrême : les deux métropoles, Douala et Yaoundé, concentrent à elles seules près de 7,6 millions d'habitants, soit une population quasi égale à la somme des 8 autres chefs-lieux de région.
    \item \textbf{Rupture d'échelle :} Il existe une \textbf{discontinuité nette} (facteur 6 à 7) entre le duo de tête ($> 3,5$ M) et le trio suivant (Bamenda, Bafoussam, Garoua $\approx 0,5$ M). Le Cameroun ne possède pas de «~villes moyennes~» nationales puissantes (1 à 2 M hab.).
    \item \textbf{Conséquence spatiale :} Cette structure bipolaire (Douala/Yaoundé) renforce la \textbf{dorsale urbaine} (axe Douala-Yaoundé-Bafoussam-Bamenda) au détriment de l'axe septentrional (Garoua-Maroua) et des périphéries Est/Sud/Adamaoua, posant le problème de l'aménagement du territoire.
\end{enumerate}
\end{exoresolu}

\courssub{TP 3 : Réaliser le croquis d'un quartier cosmopolite (Briqueterie ou New-Bell)}

\begin{definition}[Quartier cosmopolite]
Un \cle{quartier cosmopolite} est un espace urbain dense, à l'habitat souvent spontané ou dégradé, où cohabitent des populations d'origines ethniques, régionales, nationales et religieuses très diverses (ex. \textbf{Briqueterie} à Yaoundé, \textbf{New-Bell} à Douala). Il illustre la ville comme \textbf{espace d'intégration nationale} (brassage) mais aussi les failles de l'urbanisme (insalubrité, insécurité).
\end{definition}

\begin{methode}[De la photographie aérienne/plan au croquis géographique]
\begin{enumerate}
    \item \textbf{Observer le document :} identifier le \textbf{cadre} (limites du quartier), la \textbf{trame viaire} (rues principales, secondaires, impasses), le \textbf{type d'habitat} (aligné, groupé, dispersé ; matériaux : banco, parpaing, tôle), les \textbf{équipements} (marché, mosquée, église, école, forage, lampadaire), les \textbf{espaces vides} (places, décharges, terrains vagues).
    \item \textbf{Simplifier (sélectionner) :} ne garder que les éléments \textbf{structurants} (axes principaux, grands équipements, limites nettes). Gommer le détail (chaque case, chaque arbre).
    \item \textbf{Choisir les symboles (ponctuels, linéaires, surfaciques) :}
    \begin{itemize}
        \item \textit{Surfaciques (hachures/couleurs) :} Habitat dense (hachures fines), Habitat lâche (points), Zone commerciale/marché (hachures croisées), Équipements publics (plein couleur).
        \item \textit{Linéaires :} Route goudronnée (trait double épais), Piste/ruelle (trait fin pointillé), Limite de quartier (trait mixte).
        \item \textit{Ponctuels :} Mosquée (croissant/étoile), Église (croix), Forage (goutte), Lampadaire (cercle + trait), École (carré).
    \end{itemize}
    \item \textbf{Dessiner le croquis :} Cadre $\rightarrow$ Trame viaire principale $\rightarrow$ Taches d'habitat $\rightarrow$ Équipements $\rightarrow$ Légende \textbf{ordonnée} (Surfaciques / Linéaires / Ponctuels).
    \item \textbf{Indiquer l'orientation (Nord) et l'échelle approximative} (ex. 1 cm = 100 m).
\end{enumerate}
\end{methode}

\begin{popfigure}
\begin{tikzpicture}[scale=0.8]
% Cadre du croquis
\draw[popdark, thick] (0,0) rectangle (10,7);

% Nord
\node[anchor=south, font=\small\bfseries] at (9.5, 6.5) {N $\uparrow$};

% Echelle
\draw[thick, popdark] (0.5, 0.3) -- (2.5, 0.3);
\node[anchor=north, font=\tiny] at (1.5, 0.1) {200 m};

% Trame viaire principale (Axes structurants)
\draw[popdark, line width=1.5pt] (1, 0) -- (1, 7); % Nord-Sud principal
\draw[popdark, line width=1.5pt] (0, 3.5) -- (10, 3.5); % Est-Ouest principal
\draw[popdark, line width=0.8pt, dashed] (3, 0) -- (3, 7); % Secondaire
\draw[popdark, line width=0.8pt, dashed] (7, 0) -- (7, 7); % Secondaire
\draw[popdark, line width=0.5pt, dotted] (2, 1.5) -- (4, 1.5); % Ruelle
\draw[popdark, line width=0.5pt, dotted] (6, 5) -- (8, 5); % Ruelle

% Habitat : Hachures (pattern) - Zones surfaciques
% Zone Nord-Ouest : Habitat dense ancien (Banco/Tôle) - Hachures croisées fines
\foreach \x in {0.2, 0.6, 1.0, 1.4, 1.8, 2.2, 2.6} 
  \foreach \y in {4.0, 4.4, 4.8, 5.2, 5.6, 6.2} 
    \draw[popblue!60, thin] (\x, \y) -- (\x+0.3, \y+0.3) (\x+0.3, \y) -- (\x, \y+0.3);

% Zone Nord-Est : Habitat dense récent (Parpaing) - Hachures régulières
\foreach \x in {3.2, 3.6, 4.0, 4.4, 4.8, 5.2, 5.6, 6.4, 6.8} 
  \foreach \y in {4.0, 4.4, 4.8, 5.2, 5.6, 6.2} 
    \draw[poporange!70, thin, pattern=north east lines] (\x, \y) rectangle (\x+0.3, \y+0.3);

% Zone Sud-Ouest : Habitat lâche / Terrain vague (points épars)
\foreach \x in {0.5, 1.5, 2.5} 
  \foreach \y in {0.5, 1.5, 2.5} 
    \draw[popgreen!50, thin] (\x, \y) rectangle (\x+0.4, \y+0.4);

% Zone Sud-Est : Zone commerciale / Marché (Hachures croisées rouges)
\fill[pattern=north west lines, pattern color=popred!60] (3.2, 0.2) rectangle (6.8, 3.0);
\node[font=\small\bfseries, text=popdark, rotate=0] at (5, 1.6) {MARCHÉ};

% Équipements ponctuels
% Mosquée (Nord-Ouest)
\draw[poppurple, thick] (0.8, 5.5) circle (0.15); \node[font=\tiny] at (0.8, 5.5) {$\star$};
% Église (Nord-Est)
\draw[poppurple, thick] (4.5, 6.0) -- (4.5, 6.3) (4.3, 6.15) -- (4.7, 6.15);
% École (Sud-Ouest)
\draw[popgreen, thick, fill=popgreen!20] (1.5, 1.0) rectangle (2.0, 1.4); \node[font=\tiny] at (1.75, 1.2) {École};
% Forage (Centre)
\draw[popblue, thick] (5, 3.5) circle (0.1); \node[font=\tiny, anchor=north] at (5, 3.35) {Forage};

% Légende (encadrée en bas à droite)
\begin{scope}[shift={(6.5, 0.2)}]
\draw[popline] (0,0) rectangle (3.3, 3.5);
\node[anchor=west, font=\small\bfseries] at (0.1, 3.2) {LÉGENDE};
% Surfaciques
\node[anchor=west, font=\tiny\bfseries] at (0.1, 2.9) {SURFACIQUES};
\fill[pattern=north east lines, pattern color=poporange!70] (0.2, 2.5) rectangle (0.6, 2.8); \node[anchor=west, font=\tiny] at (0.7, 2.5) {Habitat dense récent (parpaing)};
\draw[popblue!60, thin] (0.2, 2.1) -- (0.5, 2.4) (0.5, 2.1) -- (0.2, 2.4); \node[anchor=west, font=\tiny] at (0.7, 2.1) {Habitat dense ancien (banco/tôle)};
\draw[popgreen!50, thin] (0.2, 1.7) rectangle (0.6, 2.0); \node[anchor=west, font=\tiny] at (0.7, 1.7) {Habitat lâche / Terrain vague};
\fill[pattern=north west lines, pattern color=popred!60] (0.2, 1.3) rectangle (0.6, 1.6); \node[anchor=west, font=\tiny] at (0.7, 1.3) {Zone commerciale / Marché};
% Linéaires
\node[anchor=west, font=\tiny\bfseries] at (0.1, 1.0) {LINÉAIRES};
\draw[popdark, line width=1.5pt] (0.2, 0.8) -- (0.6, 0.8); \node[anchor=west, font=\tiny] at (0.7, 0.75) {Route goudronnée principale};
\draw[popdark, line width=0.8pt, dashed] (0.2, 0.5) -- (0.6, 0.5); \node[anchor=west, font=\tiny] at (0.7, 0.45) {Piste / Rue secondaire};
\draw[popdark, line width=0.5pt, dotted] (0.2, 0.2) -- (0.6, 0.2); \node[anchor=west, font=\tiny] at (0.7, 0.15) {Ruelle / Impasse};
% Ponctuels
\node[anchor=west, font=\tiny\bfseries] at (0.1, -0.1) {PONCTUELS};
\node[anchor=west, font=\tiny] at (0.2, -0.3) {$\star$ Mosquée \quad $\dagger$ Église \quad $\square$ École \quad $\bullet$ Forage};
\end{scope}
\end{tikzpicture}
\legende{Croquis type d'un quartier cosmopolite (inspiré de Briqueterie/New-Bell) : trame viaire hiérarchisée, taches d'habitat différenciées par densité et matériaux, équipements structurants (marché central, lieux de culte, école, forage). La légende est ordonnée par type de symbole (Surfaciques, Linéaires, Ponctuels).}
\end{popfigure}

\begin{attention}[Erreurs fréquentes au Bac — Croquis et Carte]
\begin{itemize}
    \item \textbf{Oubli du titre, du Nord ou de l'échelle :} \textbf{0 point} pour la carte/croquis même si le contenu est bon. Ces 3 éléments sont \textbf{obligatoires}.
    \item \textbf{Légende «~en vrac~» :} Mélanger surfaciques, linéaires et ponctuels sans regroupement. $\rightarrow$ Toujours ordonner : 1. Surfaciques, 2. Linéaires, 3. Ponctuels.
    \item \textbf{Symboles non conformes :} Dessiner des petites maisons en 3D au lieu de hachures surfaciques ; dessiner des arbres réalistes au lieu de points verts. $\rightarrow$ \textbf{Symboles géographiques normalisés} (abstraits, pas pictographiques).
    \item \textbf{Cercles proportionnels mal centrés :} Le centre du cercle doit tomber \textbf{exactement} sur la localisation de la ville (le point), pas à côté.
    \item \textbf{Diagramme sans unités :} Axe Y gradué «~0, 1, 2, 3~» sans écrire «~Population (en millions d'habitants)~».
    \item \textbf{Barres collées :} Dans un diagramme en barres (série discontinue), les barres ne se touchent \textbf{pas} (espace entre elles). L'histogramme (série continue) a des barres jointes.
\end{itemize}
\end{attention}

\courssub{TP 4 : Interpréter un tableau statistique de l'évolution de la population urbaine}

\begin{exemplebox}[Objectif]
Savoir lire un tableau d'évolution (années $\times$ indicateurs), calculer des \textbf{taux de variation} et des \textbf{parts relatives}, et en tirer \textbf{deux conclusions géographiques argumentées} (chiffres à l'appui).
\end{exemplebox}

\begin{exemplebox}[Tableau support : Évolution de la population urbaine du Cameroun (1987-2023)]
\begin{center}
\begin{tabular}{|c|ccc|}\hline
\textbf{Année} & \textbf{Population totale (M hab.)} & \textbf{Population urbaine (M hab.)} & \textbf{Taux d'urbanisation (\%)} \\ \hline
1987 (RGPH) & 10,5 & 3,7 & 35,2 \\ \hline
2005 (RGPH) & 17,5 & 8,5 & 48,6 \\ \hline
2010 (Estim.) & 19,6 & 10,5 & 53,6 \\ \hline
2023 (Proj.) & 28,5 & 16,5 & 57,9 \\ \hline
\end{tabular}
\end{center}
\textit{Source : BUCREP (RGPH 1987, 2005), INS (Projections 2010, 2023).}
\end{exemplebox}

\begin{exoresolu}[TP 4 : Calculs et Conclusions]
\textbf{Énoncé :} À partir du tableau ci-dessus, calculer le taux d'accroissement annuel moyen (TAAM) de la population urbaine entre 1987 et 2023, et formuler deux conclusions géographiques.

\tcblower

\textbf{Solution détaillée :}
\begin{enumerate}
    \item \textbf{Calcul du TAAM (Taux d'Accroissement Annuel Moyen) :}
    Formule : $TAAM = \left( \left( \frac{P_{final}}{P_{initial}} \right)^{\frac{1}{n}} - 1 \right) \times 100$
    \begin{itemize}
        \item $P_{initial} = 3,7$ M (1987)
        \item $P_{final} = 16,5$ M (2023)
        \item $n = 2023 - 1987 = 36$ ans
        \item $\frac{16,5}{3,7} \approx 4,459$
        \item $4,459^{\frac{1}{36}} \approx 1,0425$ (calculatrice : $x^{1/n}$)
        \item $TAAM \approx (1,0425 - 1) \times 100 = \textbf{4,25 \% par an}$.
    \end{itemize}
    \item \textbf{Conclusion 1 (Vitesse) :} La population urbaine croît à un rythme très soutenu (\textbf{+4,25 \%/an} en moyenne sur 36 ans), bien supérieur à la croissance totale du pays ($\approx 2,8 \%$/an). Cela traduit un \textbf{exode rural massif} et un accroissement naturel fort en ville (jeune population, fécondité encore élevée).
    \item \textbf{Conclusion 2 (Structure) :} Le taux d'urbanisation passe de \textbf{35 \% (1987) à 58 \% (2023)}. Le Cameroun est devenu un \textbf{pays à majorité urbaine} (seuil des 50 \% franchi vers 2005-2010). Cette transition rapide s'accompagne des problèmes étudiés (logement, emploi, voirie) car l'offre urbaine (infrastructures, services) n'a pas suivi la demande (croissance «~par le bas~»).
\end{enumerate}
\end{exoresolu}

\courssub{Critères d'évaluation des productions cartographiques et graphiques}

\begin{aretenir}[Grille d'auto-évaluation / Évaluation Bac (sur 4-6 pts typiques)]
\begin{tabularx}{\linewidth}{|l|X|}\hline
\textbf{Critère} & \textbf{Indicateurs de réussite} \\ \hline
\textbf{1. Exactitude (Données)} & Localisation précise des villes (coordonnées) ; valeurs respectées (populations, pourcentages) ; calculs justes (TAAM, rayons proportionnels). \\ \hline
\textbf{2. Lisibilité (Graphisme)} & Trait net (règle, compas) ; écriture horizontale, lisible, sans rature ; symboles distincts (couleurs/hachures contrastées) ; espace aéré (pas de surcharge). \\ \hline
\textbf{3. Légende complète et ordonnée} & \textbf{Titre} informatif (Quoi ? Où ? Quand ?) ; \textbf{Nord} (flèche) ; \textbf{Échelle} (graphique) ; \textbf{Source} ; \textbf{Boîte de symboles} groupés par type (Surfaciques / Linéaires / Ponctuels) avec libellés précis. \\ \hline
\textbf{4. Respect de l'échelle / Proportionnalité} & Rayons des cercles $\propto \sqrt{\text{Population}}$ ; Hauteur des barres $\propto$ Valeur ; Échelle graphique cohérente avec le cadre. \\ \hline
\textbf{5. Analyse / Commentaire} & Vocabulaire géographique précis (métropolisation, primatie, exode rural, dorsale, quartier cosmopolite, habitat spontané) ; Chiffres clés cités ; Liens cause-effet explicités. \\ \hline
\end{tabularx}
\end{aretenir}

\begin{savaistu}[Repères chiffrés clés pour le Bac — Ordres de grandeur 2023-2025]
\begin{itemize}
    \item \textbf{Population totale Cameroun :} $\approx$ \textbf{28 à 29 millions} d'habitants (INS, projections post-RGPH 2005).
    \item \textbf{Taux d'urbanisation :} $\approx$ \textbf{58 \%} (majorité urbaine depuis $\approx$ 2010).
    \item \textbf{Métropoles :} \textbf{Douala} (capitale économique, port, aéroport) $\approx$ \textbf{3,8 - 4,0 M} hab. (agglomération) ; \textbf{Yaoundé} (capitale politique, administrative) $\approx$ \textbf{3,8 - 4,0 M} hab. (agglomération).
    \item \textbf{Villes secondaires (Chefs-lieux de région $> 300\,000$ hab.) :} \textbf{Bamenda} (NO, $\approx$ 600 k), \textbf{Bafoussam} (Ouest, $\approx$ 500 k), \textbf{Garoua} (Nord, $\approx$ 500 k), \textbf{Maroua} (EN, $\approx$ 400 k), \textbf{Ngaoundéré} (Adamaoua, $\approx$ 300 k).
    \item \textbf{Petits chefs-lieux ($< 200\,000$ hab.) :} \textbf{Bertoua} (Est), \textbf{Ebolowa} (Sud), \textbf{Buéa} (SO).
    \item \textbf{Primatie :} Douala + Yaoundé $\approx$ 50 \% de la population urbaine totale. Facteur de primatie (Douala / Bamenda) $\approx$ 6 à 7.
    \item \textbf{TAAM urbain (1987-2023) :} $\approx$ \textbf{4,2 - 4,5 \%} / an (très élevé, double de la moyenne nationale).
    \item \textbf{Habitat spontané :} Représente \textbf{60 à 70 \%} du tissu urbain à Douala et Yaoundé (source : MAETUR / ONU-Habitat).
\end{itemize}
\emph{Note : Les chiffres varient selon les sources (BUCREP, INS, ONU-Habitat) et la définition (commune vs agglomération). Au Bac, utilisez les ordres de grandeur du manuel officiel ou fournis dans le sujet.}
\end{savaistu}

\begin{exercice}[Entraînement Bac — Sujet type]
\begin{enumerate}
    \item \textbf{Carte (4 pts) :} Sur le fond de carte muette fourni (frontières régionales), représenter par des cercles proportionnels la population des 10 chefs-lieux de région en 2023 (données du tableau TP1). Légender complètement la carte.
    \item \textbf{Diagramme (3 pts) :} Construire le diagramme en barres de l'évolution du taux d'urbanisation (1987, 2005, 2010, 2023) à partir du tableau TP4.
    \item \textbf{Croquis (3 pts) :} À partir de la photographie aérienne du quartier New-Bell (fournie), réaliser un croquis légendé mettant en évidence : la trame viaire, les deux types d'habitat (dense ancien / dense récent), le marché central, un lieu de culte et un équipement sanitaire.
    \item \textbf{Analyse (5 pts) :} «~L'urbanisation du Cameroun est une opportunité pour l'intégration nationale mais un défi pour l'aménagement du territoire.~» Discuter cette affirmation en vous appuyant sur l'ensemble de vos productions et vos connaissances.
\end{enumerate}
\end{exercice}

\begin{corrige}[Exercice d'entraînement — Éléments de correction]
\begin{enumerate}
    \item \textbf{Carte :} Vérifier : 10 cercles aux bons endroits, 3 classes de tailles (Métropoles rouges grands, 4 grands chefs-lieux bleus moyens, 3 petits chefs-lieux verts petits), Titre «~Cameroun : Population des chefs-lieux de région (2023)~», Nord, Échelle, Source, Légende ordonnée.
    \item \textbf{Diagramme :} 4 barres (1987: 35\%, 2005: 49\%, 2010: 54\%, 2023: 58\%). Axe Y : Taux d'urbanisation (\%). Barres espacées. Titre, Source.
    \item \textbf{Croquis :} Cadre, Nord, Échelle. Trame viaire hiérarchisée (2 axes principaux perpendiculaires type damier déformé). Hachures fines pour habitat ancien (banco/tôle) au centre, hachures plus larges/régulières pour habitat récent (parpaing) en périphérie. Symbole marché (hachures croisées) au carrefour principal. Symboles ponctuels (mosquée, église, centre de santé). Légende structurée Surfaciques/Linéaires/Ponctuels.
    \item \textbf{Analyse :} 
    \textbf{Thèse 1 (Intégration) :} Villes = brassage (quartiers cosmopolites Briqueterie/New-Bell), creuset national (écoles, marchés, administrations), mobilité interrégionale.
    \textbf{Thèse 2 (Défi aménagement) :} Primatie macrocéphalique (Douala/Yaoundé), croissance non maîtrisée (habitat spontané 60-70\%), déficit infrastructures (voirie, eau, électricité), étalement spatial, pollution.
    \textbf{Synthèse :} Nécessité d'une politique de villes intermédiaires (décentralisation, pôles régionaux Garoua, Bamenda, Bafoussam) pour rééquilibrer le territoire.
\end{enumerate}
\end{corrige}

\newpage

\coursec{Activité d'intégration}

\coursec{Leçon 03 : Les villes du Cameroun}

\courssub{Généralités sur la ville et l'urbanisation au Cameroun}

\begin{definition}[Ville et urbanisation au sens de l'INS du Cameroun]
Une \cle{ville} est une localité comptant au moins \num{5000} habitants, chef-lieu d'arrondissement ou de département, dont la population active est majoritairement non agricole. L'\cle{urbanisation} est le processus de croissance de la proportion de la population vivant dans les villes.
\end{definition}

\begin{aretenir}[Chiffres clés de l'urbanisation camerounaise (estimations INS / Banque mondiale, 2023-2024)]
\begin{itemize}
\item Taux d'urbanisation : \textbf{environ 60 \%} (contre 15 \% en 1960, 50 \% en 2010).
\item Population urbaine : \textbf{environ 17 millions d'habitants} sur 29 millions.
\item Croissance urbaine annuelle : \textbf{3,5 à 4 \%} (l'une des plus rapides d'Afrique), tirée par l'exode rural et l'accroissement naturel.
\item \cle{Macrocéphalie urbaine} : Douala (capitale économique, \textasciitilde 4 millions d'hab. en agglomération) et Yaoundé (capitale politique, \textasciitilde 4 millions) concentrent plus de 40 \% de la population urbaine totale.
\end{itemize}
\end{aretenir}

\begin{exemplebox}[Hiérarchie et morphologie des villes camerounaises]
\begin{itemize}
\item \textbf{Métropoles de premier rang :} Douala (port, industrie, commerce international), Yaoundé (administration, services, diplomatie).
\item \textbf{Capitales régionales :} Bafoussam, Garoua, Maroua, Bertoua, Ngaoundéré, Bamenda, Limbe, Ebolowa, Kribi (pôles régionaux d'équipements et de commerce).
\item \textbf{Villes moyennes et petites :} Chefs-lieux d'arrondissement, relais locaux (ex. : Foumban, Dschang, Kumba, Kousséri).
\item \textbf{Morphologie type (héritage colonial + croissance spontanée) :}
  \begin{itemize}
  \item \emph{Centre-ville (plateau) :} Damier colonial, administrations, banques, grands commerces (ex. : Bonanjo à Douala, Centre administratif à Yaoundé).
  \item \emph{Quartiers populaires / spontanés :} Habitat dense, voirie lacunaire, forte mixité sociale (ex. : New-Bell, Bépanda, Nkolbisson, Mvog-Ada).
  \item \emph{Quartiers résidentiels huppés :} Villas, voirie soignée, équipements de standing (ex. : Akwa, Bonapriso, Bastos, Elig-Essono).
  \item \emph{Zones industrielles et portuaires :} Bassa, Bonabéri (Douala) ; Mvan, Nsam (Yaoundé).
  \end{itemize}
\end{itemize}
\end{exemplebox}

\courssub{La ville, espace d'intégration nationale}

\begin{propriete}[La ville comme creuset de l'unité nationale]
La ville camerounaise est le lieu privilégié du \cle{brassage des populations}. Elle accueille des migrants de \textbf{toutes les régions} (Nord, Sud, Est, Ouest, Grand Nord, Grand Sud) et des \cle{ressortissants étrangers} (Nigéria, Tchad, RCA, Mali, Guinée, France, Chine, Liban...). Ce brassage se manifeste par :
\begin{itemize}
\item La \cle{cohabitation religieuse} : mosquées, églises (catholiques, protestantes, réveil), temples traditionnels souvent mitoyens.
\item La \cle{diversité linguistique} : le \cle{camfranglais} (argot urbain mélangeant français, anglais, langues locales, pidgin) comme langue de communication interethnique ; le français et l'anglais comme langues officielles.
\item Les \cle{mariages mixtes} et les associations de développement villageois (\cle{groupements d'élites}) qui maintiennent le lien avec le terroir d'origine.
\end{itemize}
\end{propriete}

\begin{definition}[Quartier cosmopolite]
Un \cle{quartier cosmopolite} est un espace urbain où résident de manière stable et mélangée des populations d'origines géographiques, culturelles, religieuses et sociales très diverses, sans ségrégation spatiale marquée. New-Bell (Douala), Nkolbisson (Yaoundé), Haussaré (Maroua) en sont des archétypes.
\end{definition}

\begin{savaistu}[Le rôle intégrateur des marchés et gares routières]
Les grands marchés (Mboppi, Mfoundi, Central, Nkoulouloun) et les gares routières (Bassa, Mvan, Bépanda) sont des \cle{espaces de rencontre} quotidiens où s'opèrent les échanges économiques, culturels et informationnels entre Camerounais de toutes origines. Ils sont le « poumon » de l'intégration par le bas.
\end{savaistu}

\courssub{La ville, centre d'impulsion économique et de polarisation des services}

\begin{propriete}[Fonctions économiques : Douala et Yaoundé, moteurs du pays]
\begin{itemize}
\item \textbf{Douala :} Premier port d'Afrique centrale (80 \% du trafic maritime Cameroun, Tchad, RCA), zone industrielle de Bassa/Bonabéri (brasseries, aluminium, textiles, agro-alimentaire), place bancaire (sièges BEAC, banques commerciales), aéroport international.
\item \textbf{Yaoundé :} Administration centrale, ambassades, organisations internationales, enseignement supérieur (Université de Yaoundé I \& II, grandes écoles), services de santé de référence (CHU, Hôpital Général), commerce de gros et de détail.
\item \textbf{Villes secondaires :} Polarisation régionale (abattoirs, marchés de gros, usines de transformation agricole : coton à Garoua, thé à Ndu, bois à Bertoua, bauxite à Ngaoundal).
\end{itemize}
\end{propriete}

\begin{propriete}[Polarisation des services urbains]
La ville concentre les \cle{équipements structurants} :
\begin{itemize}
\item \textbf{Santé :} Hôpitaux de référence (CHU, Hôpitaux Généraux, Hôpitaux Centraux), cliniques privées, pharmacies.
\item \textbf{Éducation :} Universités, écoles normales supérieures, instituts universitaires de technologie, lycées techniques, collèges, écoles primaires.
\item \textbf{Administration :} Préfectures, sous-préfectures, délégations régionales/minisérielles, tribunaux, centres d'état civil.
\item \textbf{Réseaux :} Eau (CAMWATER), Électricité (ENEO), Télécoms (MTN, Orange, Camtel, Viettel), Voirie (routes bitumées, ponts, échangeurs).
\end{itemize}
\end{propriete}

\begin{attention}[Polarisation \emph{subie} vs Polarisation \emph{maîtrisée}]
La croissance urbaine rapide (\textgreater 4 \%/an) dépasse souvent la capacité d'investissement public. La polarisation devient \emph{subie} : les équipements (écoles, santé, eau, électricité) sont saturés ou absents dans les quartiers périphériques spontanés. C'est le cœur du \cle{défi urbain camerounais}.
\end{attention}

\courssub{Les problèmes des villes camerounaises}

\begin{methode}[Analyser les problèmes urbains : une grille de lecture en 4 dimensions]
Pour tout commentaire de document ou étude de cas au Bac, classe les problèmes dans ce tableau :
\begin{center}
\begin{tabularx}{\linewidth}{|l|X|}\hline
\rowcolor{popblueL} \textbf{Dimension} & \textbf{Manifestations concrètes au Cameroun} \\ \hline
\cle{Sociale} & Déficit de logements (bidonvilles, habitat précaire en planches/banco), surpeuplement, manque d'écoles (effectifs \textgreater 80/élèves/classe), insuffisance centres de santé, chômage des jeunes, exclusion sociale. \\ \hline
\cle{Environnementale} & \cle{Insalubrité} (ordures non ramassées, dépôts sauvages), \cle{inondations} récurrentes (caniveaux obstrués, occupation des zones non aedificandi), pollution de l'air (vieux véhicules, industries), érosion des sols, perte de biodiversité urbaine. \\ \hline
\cle{Économique} & \cle{Voirie urbaine} dégradée (rues non bitumées, nids-de-poule, embouteillages chroniques), transport informel dominant (mototaxis \emph{benskin}, taxis vétustes), coût de la vie élevé, précarité de l'économie informelle (marchands ambulants, artisans sans local). \\ \hline
\cle{Sécuritaire} & \cle{Insécurité} (vols, agressions, \emph{feymen}, gangs), absence d'éclairage public, incivisme (non-respect code de la route, constructions anarchiques), conflits fonciers. \\ \hline
\end{tabularx}
\end{center}
\end{methode}

\begin{savaistu}[Acteurs de la gestion urbaine]
\begin{itemize}
\item \textbf{État :} MINHDU (Ministère de l'Habitat et du Développement Urbain), MINEPAT (Planification), MAETUR (Aménagement du territoire - structures déconcentrées), CAMWATER, ENEO (concessions).
\item \textbf{Collectivités Territoriales Décentralisées (CTD) :} Communautés Urbaines (CUD, CUY), Communes d'arrondissement, Communes rurales. Compétences : voirie communale, marchés, état civil, hygiène, éclairage, Plan Local de Développement (PLD).
\item \textbf{Privé / Société civile :} HYSACAM (ramassage ordures Douala/Yaoundé), ONG (assainissement, habitat), groupements de quartier, syndicats de transporteurs.
\end{itemize}
\end{savaistu}

\courssub{Solutions et perspectives pour les villes camerounaises}

\begin{exemplebox}[Réponses institutionnelles et opérationnelles récentes]
\begin{itemize}
\item \textbf{Planification :} Schémas Directeurs d'Aménagement Urbain (SDAU), Plans d'Urbanisme (PU), Plans Locaux d'Urbanisme (PLU) — obligation légale pour toute commune.
\item \textbf{Programmes structurants :} \emph{Programme d'Urgence Triennal (PUT)}, \emph{Plan d'Urgence pour la Résorption des Bidonvilles}, \emph{Programme d'Asphaltage des Villes} (bitumage voirie secondaire), \emph{Projet d'Assainissement de Douala (PAD)}.
\item \textbf{Logement :} Création de la \cle{SIC} (Société Immobilière du Cameroun), programmes de logements sociaux (ex. : 10 000 logements, cités SIC à Douala/Yaoundé/Bafoussam), réforme du foncier (titres fonciers, certificats d'urbanisme).
\item \textbf{Transport :} Projet de Bus Rapid Transit (BRT) à Douala, réhabilitation du chemin de fer (Camrail) pour le fret et voyageurs banlieue, régulation du \emph{benskin} (casques, assurances, cartes grises).
\item \textbf{Gouvernance :} Budget participatif dans certaines communes (ex. : Douala 1er, Yaoundé 2e), numérisation des services fonciers (e-foncier), police municipale renforcée.
\end{itemize}
\end{exemplebox}

\begin{aretenir}[Leviers pour une ville durable et intégrée]
\begin{enumerate}
\item \cle{Maîtrise de l'étalement} : densification contrôlée, renouvellement urbain (résorption bidonvilles \emph{in situ} ou relocalisation concertée).
\item \cle{Investissement dans les réseaux} : eau, électricité, assainissement, voirie hiérarchisée (structurante, secondaire, desserte).
\item \cle{Gestion des déchets} : tri à la source, valorisation (compost, biogaz), extension couverture HYSACAM / entreprises locales.
\item \cle{Éducation à la citoyenneté} : lutte contre l'\cle{incivisme} (dépôts sauvages, stationnement anarchique, raccordements frauduleux eau/électricité).
\item \cle{Financement local} : amélioration de la fiscalité foncière et des patentes, partenariats public-privé (PPP).
\end{enumerate}
\end{aretenir}

\courssub{Travaux pratiques : Représenter l'urbanisation du Cameroun — Étude de cas « Mon quartier de demain : New-Bell »}

\courssub{Objectifs de l'activité}

À l'issue de cette activité d'intégration, tu seras capable de :

\begin{itemize}
\item \cle{observer} et \cle{décrire} un paysage urbain camerounais à partir d'une photographie et d'une carte ;
\item \cle{localiser} un quartier dans la ville et \cle{légender} une carte muette de façon ordonnée ;
\item \cle{inventorier}, \cle{classer} et \cle{interpréter} des documents géographiques (photographie, tableau statistique, article de presse) ;
\item \cle{construire un croquis} de quartier avec titre, orientation et légende numérotée ;
\item \cle{proposer} des solutions d'aménagement réalistes et les justifier, en mobilisant les savoirs de la leçon (ville, espace d'intégration ; centre d'impulsion économique ; problèmes des villes).
\end{itemize}

Cette activité mobilise l'ensemble des savoir-faire du programme officiel : observer, localiser et décrire ; représenter, schématiser et cartographier ; inventorier, classer et interpréter des documents.

\courssub{Situation}

\textbf{Communiqué de la Communauté Urbaine de Douala (CUD).} Dans le cadre de la célébration de la Journée de la Jeunesse, la Communauté Urbaine de Douala lance un concours intitulé \emph{« Mon quartier de demain »}, ouvert aux classes de Terminale de la ville. Chaque équipe candidate reçoit un dossier documentaire portant sur le quartier \cle{New-Bell}, situé dans l'arrondissement de Douala II, à l'est du centre-ville (Bonanjo), entre le marché Mboppi et les berges du Wouri. Le dossier comprend quatre documents :

\begin{itemize}
\item \textbf{Document 1 — Photographie aérienne décrite du quartier New-Bell.} On y voit un habitat dense et hétérogène : des maisons en matériaux définitifs (parpaings, tôles) côtoient des constructions en planches et en briques de terre ; les rues, non goudronnées pour la plupart, sont envahies par les eaux stagnantes en saison des pluies ; des amoncellements d'ordures jonchent les caniveaux ; de nombreux petits commerces (boutiques, ateliers de couture, kiosques à crédit téléphonique, gargotes) animent les carrefours ; on aperçoit une mosquée, une église et une école primaire publique surchargée.
\item \textbf{Document 2 — Carte muette de la ville de Douala} (contour de la ville, fleuve Wouri, pont Wouri, principaux axes routiers, à l'échelle indicative 1/100 000).
\item \textbf{Document 3 — Tableau statistique du quartier} (données fictives mais réalistes, à but pédagogique) :
\end{itemize}

\begin{center}
\begin{tabularx}{\linewidth}{|l|X|}\hline
\rowcolor{popblueL} \textbf{Indicateur} & \textbf{Valeur} \\ \hline
Population estimée du quartier & environ \num{120000} habitants \\ \hline
Superficie approximative & 6 km\textsuperscript{2} \\ \hline
Part des moins de 20 ans & environ 55 \% \\ \hline
Rues goudronnées & environ 15 \% de la voirie \\ \hline
Ménages raccordés à l'eau potable (CAMWATER) & environ 40 \% \\ \hline
Ménages raccordés à l'électricité (ENEO) & environ 65 \% \\ \hline
Écoles primaires publiques & 4 (pour un besoin estimé de 12) \\ \hline
Centres de santé & 2 \\ \hline
Origine des habitants & toutes les régions du Cameroun, plus des ressortissants du Nigeria, du Tchad et du Mali \\ \hline
\end{tabularx}
\end{center}

\begin{itemize}
\item \textbf{Document 4 — Extrait de presse} (\emph{Le Messager}, édition fictive) : « À New-Bell, les montagnes de déchets non ramassés étouffent les caniveaux. À chaque pluie, l'eau monte dans les maisons. Les habitants dénoncent l'insalubrité et l'insécurité nocturne, faute d'éclairage public. »
\end{itemize}

\textbf{Tâche.} Ta classe, organisée en équipes, doit produire un dossier de candidature comprenant : la localisation et la légende du quartier sur la carte muette, la description de la photographie, un tableau de classification des problèmes, un croquis du quartier, et trois propositions d'aménagement justifiées. La production sera présentée au tableau et évaluée selon la grille commune : exactitude des informations (4 points), qualité du croquis et de sa légende (8 points), pertinence et justification des solutions (8 points).

\courssub{Consignes}

\begin{enumerate}
\item Sur la carte muette de Douala (document 2), localise le quartier New-Bell, place le fleuve Wouri, le port, le centre-ville (Bonanjo) et l'aéroport, puis rédige une \cle{légende ordonnée} (numérotée, du plus important au plus détaillé).
\item Décris la photographie (document 1) en suivant la démarche : description générale, puis description détaillée du premier plan vers l'arrière-plan. Relève les indices du \cle{brassage des populations} (langues, commerces, lieux de culte) qui font du quartier un espace d'intégration nationale.
\item À partir des documents 1, 3 et 4, inventorie les problèmes du quartier et classe-les dans un tableau à quatre colonnes : problèmes sociaux, problèmes environnementaux, problèmes économiques, problèmes sécuritaires.
\item Calcule la \cle{densité de population} du quartier et le déficit en écoles primaires. Que conclus-tu sur la polarisation insuffisante des services ?
\item Construis un \cle{croquis} du quartier New-Bell : contour simplifié, orientation (flèche du nord), titre, et légende numérotée hiérarchisée (voirie, habitat, équipements, points noirs).
\item Propose \textbf{trois solutions} d'aménagement réalistes (une par catégorie de problèmes au minimum) et justifie chacune en une phrase, en précisant l'acteur concerné (commune, État, habitants, entreprises).
\end{enumerate}

\courssub{Ressources à mobiliser}

\begin{aretenir}[Ce qu'il faut se rappeler]
\begin{itemize}
\item La \cle{ville} est un espace d'\cle{intégration nationale} : elle brasse des populations de toutes les régions, des cultures et des religions différentes ; un quartier comme New-Bell est un \cle{quartier cosmopolite}.
\item La ville est un \cle{centre d'impulsion économique} et un lieu de \cle{polarisation des services} (écoles, santé, eau, électricité, marchés).
\item Les problèmes des villes camerounaises : \cle{insalubrité}, \cle{anarchie} urbaine, \cle{incivisme}, \cle{insécurité}, déficits de logement, de transport, d'eau et d'électricité, \cle{voirie urbaine} dégradée.
\item Densité de population : $D = \dfrac{\text{population}}{\text{superficie}}$, exprimée en habitants/km\textsuperscript{2}.
\item Un bon croquis exige : un \textbf{titre}, une \textbf{orientation}, une \textbf{légende numérotée et hiérarchisée} (du général au particulier), placée hors du dessin.
\end{itemize}
\end{aretenir}

\courssub{Démarche et corrigé commenté}

\begin{exoresolu}[Mon quartier de demain : dossier de candidature corrigé]
Énoncé : à partir du dossier documentaire sur New-Bell, produire la localisation, la description, la classification des problèmes, le croquis et trois solutions d'aménagement.

\tcblower

\textbf{Étape 1 — Localisation et légende de la carte.}
New-Bell se situe dans l'est de Douala (arrondissement Douala II), entre le centre-ville historique (Bonanjo, sur la rive gauche du Wouri à l'ouest) et les quartiers périphériques (Bépanda, PK8). Sur la carte muette, on place et on légende dans l'ordre hiérarchique :
\begin{enumerate}
\item \textbf{Hydrographie structurante :} Fleuve Wouri (coule du Nord-Est au Sud-Ouest), Pont sur le Wouri (liaison Bonanjo -- Bonabéri).
\item \textbf{Pôles majeurs :} Port autonome de Douala (rive gauche, au sud de Bonanjo), Centre-ville / Bonanjo (administrations, banques), Aéroport international de Douala (sud-est, hors centre dense).
\item \textbf{Quartier d'étude :} New-Bell (zone dense à l'est de Bonanjo, bordée au nord par l'axe vers Bépanda/PK8, au sud par le Wouri).
\item \textbf{Réseau viaire principal :} Axes bitumés (Boulevard de la Liberté, Rue de l'Hôpital, Route de Bépanda).
\end{enumerate}
\emph{Astuce :} La légende est numérotée du plus structurant (fleuve, port) au plus local (quartier, axes secondaires).

\textbf{Étape 2 — Description de la photographie.}
\emph{Description générale :} Il s'agit d'un paysage urbain dense et contrasté, typique d'un quartier populaire camerounais en croissance spontanée, où l'habitat, l'activité économique et les équipements se mêlent dans un espace saturé.
\emph{Description détaillée (premier plan $\rightarrow$ arrière-plan) :}
\begin{itemize}
\item \textbf{Premier plan :} Rues en terre battue, dégradées, bordées de caniveaux ouverts obstrués par des déchets solides (plastiques, biodégradables) ; eaux stagnantes témoins d'un drainage défaillant.
\item \textbf{Deuxième plan :} Habitat hétérogène et mitoyen : maisons en parpaings et tôles (matériaux définitifs) côtoyant des cases en planches, briques de terre crue ou tôles récupérées (habitat précaire) ; toitures souvent en tôles ondulées.
\item \textbf{Troisième plan / Arrière-plan :} Silhouettes d'équipements collectifs : une mosquée (minaret), une église (clocher/croix) voisines, une école primaire publique (bâtiment long, cour récréative bondée) ; animation commerciale intense aux carrefours (étals, boutiques, ateliers, kiosques \emph{call-box}).
\end{itemize}
\emph{Indices du brassage (intégration nationale) :} La coexistence immédiate d'une mosquée et d'une église symbolise la tolérance religieuse ; la diversité des commerces (restauration variée, ateliers, téléphonie) attire une clientèle multiethnique ; le document 3 confirme la présence d'habitants originaires des 10 régions du Cameroun et de 3 pays voisins. New-Bell est un \cle{quartier cosmopolite} archétypal.

\textbf{Étape 3 — Classification des problèmes.}

\begin{center}
\begin{tabularx}{\linewidth}{|l|X|}\hline
\rowcolor{popblueL} \textbf{Catégorie} & \textbf{Problèmes relevés (Documents 1, 3, 4)} \\ \hline
\cle{Sociaux} & Déficit scolaire criant (4 écoles pour 12 nécessaires = 8 manquantes), effectifs pléthoriques ; seulement 2 centres de santé pour 120 000 hab. ; habitat précaire (planches, banco) ; accès à l'eau potable limité (40 \% raccordés CAMWATER) ; jeunesse nombreuse (55 \% \textless 20 ans) sans encadrement suffisant. \\ \hline
\cle{Environnementaux} & \cle{Insalubrité} massive : déchets non ramassés (Doc 1, 4), caniveaux obstrués, dépôts sauvages ; \cle{inondations} saisonnières (eaux stagnantes, montée des eaux dans les maisons Doc 4) ; absence d'espaces verts. \\ \hline
\cle{Économiques} & \cle{Voirie urbaine} très dégradée (seulement 15 \% bitumés) : frein aux échanges, usure véhicules, coût transport élevé ; précarité de l'économie informelle (locaux inadaptés, pas de sécurité juridique) ; coupures d'électricité récurrentes (35 \% non raccordés ENEO). \\ \hline
\cle{Sécuritaires} & \cle{Insécurité} nocturne (Doc 4) favorisée par l'absence d'éclairage public ; incivisme (dépôts d'ordures, stationnement anarchique) ; risques sanitaires (eaux souillées, promiscuité). \\ \hline
\end{tabularx}
\end{center}

\textbf{Étape 4 — Calculs et interprétation.}
\begin{align*}
\text{Densité } D &= \frac{\text{Population}}{\text{Superficie}} = \frac{120\,000\ \text{hab.}}{6\ \text{km}^2} = 20\,000\ \text{hab./km}^2 \\[0.5em]
\text{Déficit scolaire} &= \text{Besoin estimé} - \text{Existant} = 12 - 4 = 8\ \text{écoles primaires publiques manquantes}
\end{align*}
\emph{Conclusion :} Une densité de \num{20000} hab/km\textsuperscript{2} est très élevée (comparable aux densités des centres-villes parisiens), signe d'un fort peuplement spontané non accompagné. Le déficit de 8 écoles (67 \% du besoin) et le ratio d'un centre de santé pour 60 000 hab. prouvent que la \cle{polarisation des services} est \textbf{insuffisante et débordée} par la croissance démographique rapide. La ville joue mal son rôle d'intégration par les services publics.

\textbf{Étape 5 — Croquis du quartier New-Bell.}

\begin{popfigure}
\begin{center}
\includegraphics[width=\linewidth,height=10.5cm,keepaspectratio]{N03/douala_map.jpg}
\end{center}
\legende{Plan de Douala (fond OpenStreetMap). On y lit le Wouri et son estuaire à l'ouest, le centre historique (Bonanjo, Akwa) sur la rive gauche, le quartier New Bell, au sud de la gare centrale, près de l'ancien aéroport, l'ancien aéroport, la zone industrielle de Bassa, les grands axes (en rouge et orange) et, sur l'autre rive, Bonabéri. Le croquis attendu du quartier New-Bell, à construire à partir de ce fond, représente par une légende hiérarchisée : 1. le fleuve Wouri (limite ouest) ; 2. la voirie principale ; 3. la voirie secondaire ; 4. l'habitat dense et hétérogène ; 5. les écoles ; 6. les lieux de culte ; 7. les pôles commerciaux ; 8. les points noirs environnementaux. Nord en haut. \\ Source : Wikimedia Commons, contributeurs d'OpenStreetMap, CC BY-SA 3.0.}
\end{popfigure}

\textbf{Étape 6 — Trois solutions d'aménagement justifiées.}
\begin{enumerate}
\item \textbf{Dimension environnementale / Salubrité :} \textbf{Mise en place d'un ramassage régulier des ordures ménagères par la CUD (via HYSACAM) avec installation de points de tri et de pré-collecte dans chaque îlot, couplée à une campagne de sensibilisation contre l'incivisme (dépôts sauvages).} \emph{Justification :} L'insalubrité (Doc 1, 3, 4) cause les inondations et les maladies hydriques ; la commune est compétente en matière d'hygiène et de salubrité publique (Code général des CTD).
\item \textbf{Dimension économique / Voirie et Sécurité :} \textbf{Bitumage des axes structurants secondaires et installation d'un éclairage public solaire autonome sur les artères principales et places publiques (Partenariat Commune / MINHDU / ENEO / Partenaires privés).} \emph{Justification :} Une voirie praticable (actuellement 15 \% seulement) fluidifie l'économie locale, désenclave le quartier et l'éclairage public réduit drastiquement l'insécurité nocturne dénoncée par les habitants (Doc 4).
\item \textbf{Dimension sociale / Services de base :} \textbf{Construction de 4 écoles primaires publiques et d'un Centre de Santé Intégré (CSI) neuf sur foncier communal sécurisé, avec recrutement d'enseignants et de personnel soignant (État : MINEDUB, MINSANTE ; Appui : Commune, PTF).} \emph{Justification :} Le déficit de 8 écoles et le ratio 1 centre/60 000 hab. sont inacceptables pour une population à 55 \% jeune ; l'éducation et la santé sont des compétences régaliennes (État) et partagées (Commune), condition sine qua non de l'intégration par les services.
\end{enumerate}
\end{exoresolu}

\begin{attention}[Erreurs fréquentes à éviter au Bac et en TP]
\begin{itemize}
\item \textbf{Légende de croquis / carte :} Non numérotée, désordonnée (alphabétique au lieu de hiérarchique), placée \emph{à l'intérieur} du dessin (illisible), symboles non rapportés sur le dessin, absence d'échelle ou de nord.
\item \textbf{Description de photo :} Confondre \emph{description} (inventaire visuel ordonné : plan par plan) et \emph{explication/interprétation} (causes, conséquences). Ne pas citer les indices visuels précis (ex. : « maisons en planches » et pas « habitat précaire » seul).
\item \textbf{Classification des problèmes :} Mélanger les catégories (ex. mettre « chômage » en environnemental), oublier une dimension (souvent la sécuritaire ou l'économique), lister sans citer la source du document (Doc 1, Doc 3, Doc 4).
\item \textbf{Calculs :} Oublier l'unité (hab/km\textsuperscript{2}), inverser population/superficie, ne pas conclure (lien densité $\rightarrow$ pression sur services).
\item \textbf{Solutions :} Propositions irréalistes (ex. « démolir tout le quartier », « construire une autoroute »), sans acteur identifié (« l'État doit faire »), sans justification liée au diagnostic, ou hors sujet (ex. « créer une université » alors que le déficit est au primaire).
\end{itemize}
\end{attention}

\courssub{Bilan de l'activité}

Cette activité t'a permis de mettre en évidence trois idées essentielles de la leçon, directement mobilisables pour l'épreuve du Baccalauréat.

D'abord, la ville camerounaise est un véritable \cle{espace d'intégration nationale} : à New-Bell comme dans la plupart des quartiers populaires des métropoles, des populations de toutes les régions (Nord, Ouest, Sud, Est, Centre, Littoral) et de plusieurs confessions religieuses (musulmans, chrétiens, animistes) vivent, commercent, s'instruisent et se marient ensemble. La coexistence d'une mosquée et d'une église, la pratique du \cle{camfranglais} et la diversité des origines (Doc 3) en sont la preuve tangible.

Ensuite, la ville joue pleinement son rôle de \cle{centre d'impulsion économique} (proximité port, marché Mboppi, activités informelles denses) et de \cle{polarisation des services}, mais les calculs de densité (\num{20000} hab/km\textsuperscript{2}) et le déficit criant d'équipements (écoles, santé, eau, électricité, voirie) montrent que cette polarisation est souvent \textbf{débordée par la croissance urbaine rapide} et l'insuffisance des investissements publics. La ville « subit » son urbanisation plus qu'elle ne la « gère ».

Enfin, tu as constaté que les problèmes urbains — \cle{insalubrité}, \cle{anarchie} urbaine, \cle{incivisme}, \cle{insécurité}, déficits de logement, de transport, d'eau, d'électricité, \cle{voirie urbaine} dégradée — ne sont pas une fatalité. Ils peuvent être \textbf{diagnostiqués méthodiquement} (observation, classification, cartographie via le croquis) et faire l'objet de \textbf{solutions concrètes, hiérarchisées et actorialisées} (Commune, État, Habitants, Privé).

Tu as ainsi mobilisé l'ensemble des savoir-faire du programme officiel de Terminale :
\begin{itemize}
\item \cle{Observer et décrire} un paysage urbain (photo) ;
\item \cle{Localiser et légender} sur carte muette ;
\item \cle{Inventorier, classer et interpréter} des documents statistiques et textuels ;
\item \cle{Construire un croquis} géographique normé (titre, nord, légende hiérarchisée) ;
\item \cle{Argumenter} des propositions d'aménagement géographiquement fondées.
\end{itemize}
Ces compétences sont le cœur de l'évaluation au Baccalauréat (commentaire de document, croquis, étude de cas). Entraîne-toi à les reproduire sur d'autres quartiers (Nkolbisson, Bépanda, Mvog-Ada, Haussaré) et d'autres villes (Yaoundé, Bafoussam, Maroua).

\newpage

\coursec{Méthodes et exercices résolus}

\courssub{Méthode 1 : Observer, localiser et décrire un espace urbain (photographie, plan, carte)}

\begin{methode}[Décrire un paysage urbain en quatre étapes]
La description d'un document urbain (photographie de Douala, plan de Yaoundé, carte d'une région) est la première étape de tout commentaire au Bac. Elle obéit à une démarche rigoureuse.
\begin{enumerate}
\item \textbf{Identifier et localiser le document} : nature du document (photographie aérienne, carte au 1/200 000, plan de ville), lieu représenté, date, échelle, orientation (flèche du nord). Toujours commencer par une phrase de présentation complète.
\item \textbf{Décrire du général au particulier} : d'abord l'organisation globale (centre, périphérie, axes structurants), puis les détails (types d'habitat, voirie, équipements, espaces verts, quartiers spontanés).
\item \textbf{Nommer avec précision} : employer le vocabulaire géographique exact (\cle{centre-ville}, \cle{quartier résidentiel}, \cle{quartier sous-équipé}, \cle{zone industrielle}, \cle{voirie urbaine}, \cle{habitat spontané}), et localiser dans le document (au premier plan, au nord, le long de l'axe...).
\item \textbf{Hiérarchiser et conclure} : dégager le ou les traits dominants (ex. : contraste entre quartiers planifiés et quartiers spontanés) et annoncer l'interprétation qui suivra.
\end{enumerate}
\textbf{Réflexes} : ne jamais interpréter pendant la description ; ne jamais écrire « on voit » ; utiliser les repères de la légende ; chiffrer dès que le document le permet (distances, superficies, densités).
\end{methode}

\begin{exoresolu}[Description d'une photographie de Douala]
\textbf{Énoncé (type Bac)} : À partir d'une photographie aérienne récente de Douala montrant au premier plan le port et les entrepôts de Bonabéri, au centre le quartier des affaires d'Akwa avec ses immeubles de bureaux, et à l'arrière-plan des quartiers denses d'habitat spontané (New Bell, Ndogpassi), décrivez cette photographie et dégagez-en les traits dominants de l'espace urbain doualais.
\tcblower
\textbf{Solution détaillée.}
\textbf{Étape 1 — Présentation du document.} Le document est une photographie aérienne oblique de Douala, capitale économique du Cameroun et plus grande ville du pays (environ 3,5 à 4 millions d'habitants, ordre de grandeur). La prise de vue, orientée vers l'est, montre la ville depuis la rive du Wouri.
\textbf{Étape 2 — Description du général au particulier.} L'espace photographié s'organise en trois ensembles. Au premier plan, le long du fleuve Wouri, s'étend la zone portuaire et industrielle : quais, grues, entrepôts, silos, vastes emprises planes. Au centre, le quartier d'Akwa forme un centre d'affaires dense : immeubles de bureaux de plusieurs étages, banques, hôtels, voirie large et bitumée tracée en damier. À l'arrière-plan s'étendent à perte de vue des quartiers d'habitat dense et bas : toitures de tôle serrées les unes contre les autres, ruelles étroites et non bitumées, absence d'équipements collectifs visibles ; il s'agit de quartiers spontanés comme New Bell ou Ndogpassi.
\textbf{Étape 3 — Localisation précise.} Le port occupe la rive (premier plan), le centre d'affaires lui est immédiatement juxtaposé, et l'habitat spontané occupe la périphérie orientale, ce qui montre une organisation concentrique dégradée : les fonctions nobles au centre, l'habitat précaire en périphérie.
\textbf{Étape 4 — Traits dominants et conclusion.} Deux traits dominent : d'une part la concentration des fonctions économiques supérieures (port, affaires) au centre, qui fait de Douala le centre d'impulsion économique du Cameroun ; d'autre part le contraste spatial entre quartiers planifiés et équipés et quartiers spontanés sous-équipés, signe d'une urbanisation rapide et mal maîtrisée. Cette photographie illustre ainsi à la fois le rôle moteur de la ville et les problèmes urbains qu'elle concentre.
\end{exoresolu}

\courssub{Méthode 2 : Construire et légender un croquis urbain}

\begin{methode}[Réaliser un croquis de ville ou un croquis de localisation]
Le croquis est un exercice fréquent au Bac (TP de cartographie). Il vaut des points sur la forme autant que sur le fond.
\begin{enumerate}
\item \textbf{Analyser le sujet et sélectionner l'information} : un croquis n'est pas une carte : on ne reporte que ce qui éclaire le sujet (ex. : pour « Douala, métropole portuaire », on reporte le port, le fleuve, les axes, les quartiers — pas tout le reste).
\item \textbf{Tracer le fond} : contour simplifié de la ville ou du pays, principales limites (fleuve Wouri, littoral), en traits fins.
\item \textbf{Placer les faits essentiels} : quartiers (zones colorées ou hachurées), villes et équipements (points et symboles), axes de transport (traits), flux (flèches d'épaisseur proportionnelle à l'importance).
\item \textbf{Orienter et situer} : flèche du nord, échelle indicative, titre explicite en haut.
\item \textbf{Rédiger la légende} : ordonnée (d'abord les limites et le fond, puis les faits localisés, puis les dynamiques/flux), avec des symboles identiques à ceux du croquis. La légende doit permettre de lire le croquis sans le texte.
\end{enumerate}
\textbf{Réflexes} : sobriété (4 à 6 entrées de légende maximum), propreté, cohérence symbole/légende, titre qui reprend les mots du sujet.
\end{methode}

\begin{exoresolu}[Croquis : Douala, métropole portuaire et industrielle]
\textbf{Énoncé (type Bac)} : À l'aide des documents fournis (carte de l'estuaire du Wouri, tableau du trafic portuaire), construisez un croquis intitulé « Douala, métropole portuaire et industrielle du Cameroun ».
\tcblower
\textbf{Solution détaillée.}
\textbf{Raisonnement préalable.} Le sujet impose de montrer : (a) la situation de Douala sur l'estuaire du Wouri ; (b) le port, cœur de la fonction métropolitaine ; (c) les quartiers et la zone industrielle ; (d) les axes qui connectent Douala à son arrière-pays (Yaoundé, l'intérieur du pays, le Tchad et la Centrafrique dont le port de Douala est le débouché).
\textbf{Construction pas à pas.} (1) Je trace la côte et l'estuaire du Wouri en traits fins, avec la flèche du nord. (2) Je place par des zones : le port de Douala (Bonanjo sur la rive est, Bonabéri sur la rive ouest, relié par le pont sur le Wouri) en bleu foncé, la zone industrielle de Bonabéri-Bassa en hachures, le centre-ville (Akwa-Bonanjo) en pointillé dense, les quartiers périphériques en gris clair. (3) Je trace les axes : route Douala--Yaoundé, chemin de fer Transcamerounais vers Ngaoundéré, route vers le Nord-Ouest. (4) Je figure les flux par flèches : importations (vers l'intérieur) et exportations (pétrole, bois, cacao, bananes vers la mer), ainsi que les flux à destination du Tchad et de la Centrafrique. (5) Je rédige la légende ordonnée : 1. limites et hydrographie ; 2. les quartiers et fonctions ; 3. les axes ; 4. les flux.
\textbf{Conclusion.} Le croquis obtenu montre que la puissance de Douala repose sur la jonction entre la mer (port) et l'arrière-pays (axes), ce qui en fait la porte d'entrée du commerce extérieur camerounais : l'essentiel du commerce extérieur du pays (plus de 90 \%, ordre de grandeur couramment cité) y transite.
\end{exoresolu}

\courssub{Méthode 3 : Inventorier, classer et interpréter des données statistiques urbaines}

\begin{methode}[Exploiter un tableau statistique sur les villes]
\begin{enumerate}
\item \textbf{Lire le tableau en entier} : titre, unités, source, dates, lignes et colonnes. Une donnée sans unité ni date est inutilisable.
\item \textbf{Classer} : ranger les villes par ordre décroissant de population, regrouper par seuils (plus de 1 million ; 500 000 à 1 million ; moins de 500 000) ou par catégories (métropoles, capitales régionales, villes moyennes, petites villes).
\item \textbf{Calculer les indicateurs utiles} :
\begin{itemize}
\item part en pourcentage : $\text{part} = \dfrac{\text{effectif de la catégorie}}{\text{effectif total}} \times 100$ ;
\item taux d'urbanisation : $\text{TU} = \dfrac{\text{population urbaine}}{\text{population totale}} \times 100$ ;
\item taux d'accroissement (en \%) : $\dfrac{V_f - V_i}{V_i} \times 100$ ;
\item taux d'accroissement annuel moyen (en \%/an) : $\left[\left(\dfrac{V_f}{V_i}\right)^{1/n} - 1\right] \times 100$, où $n$ est le nombre d'années.
\end{itemize}
\item \textbf{Interpréter} : comparer (écarts, rapports), dégager les tendances (concentration, croissance rapide), puis expliquer par les mécanismes géographiques (exode rural, migrations, croissance naturelle, polarisation des emplois et services).
\end{enumerate}
\textbf{Réflexes} : toujours donner le résultat avec son unité et son arrondi ; vérifier qu'une somme de parts fait 100 \% ; justifier chaque calcul par la formule écrite.
\end{methode}

\begin{exoresolu}[Taux d'urbanisation et poids des deux métropoles]
\textbf{Énoncé (type Bac)} : On donne les données suivantes (ordres de grandeur, projections récentes) : population totale du Cameroun : 28 millions d'habitants ; population urbaine : 16,5 millions ; Douala : 3,8 millions ; Yaoundé : 4,2 millions ; ensemble des dix capitales régionales hors Douala et Yaoundé : 3,5 millions. (1) Calculez le taux d'urbanisation du Cameroun. (2) Calculez la part de Douala et Yaoundé dans la population urbaine. (3) Interprétez ces résultats en deux idées.
\tcblower
\textbf{Solution détaillée.}
\textbf{(1) Taux d'urbanisation.} On applique la formule :
\[ \text{TU} = \dfrac{\text{population urbaine}}{\text{population totale}} \times 100 = \dfrac{16{,}5}{28} \times 100 \approx 58{,}9\,\%. \]
Le taux d'urbanisation du Cameroun est d'environ \textbf{59 \%} : près de 6 Camerounais sur 10 vivent en ville, ce qui traduit une urbanisation avancée et rapide (il était inférieur à 30 \% dans les années 1970).
\textbf{(2) Part de Douala et Yaoundé dans la population urbaine.} Population cumulée des deux métropoles : $3{,}8 + 4{,}2 = 8$ millions d'habitants.
\[ \text{Part} = \dfrac{8}{16{,}5} \times 100 \approx 48{,}5\,\%. \]
Douala et Yaoundé concentrent à elles seules près de la \textbf{moitié de la population urbaine} du pays.
\textbf{(3) Interprétation.} Première idée : l'armature urbaine camerounaise est déséquilibrée, dominée par deux métropoles qui cumulent les fonctions politique (Yaoundé, capitale) et économique (Douala, port et industries) ; les dix autres capitales régionales ne pèsent ensemble que $3{,}5/16{,}5 \times 100 \approx 21\,\%$ de la population urbaine, soit moins de la moitié du poids des deux métropoles. Deuxième idée : cette macrocéphalie s'explique par la polarisation des emplois, des services (universités, hôpitaux, administrations) et des infrastructures, qui entretient l'exode rural vers les deux grandes villes et pose des problèmes de logement, de transport et d'assainissement.
\textbf{Conclusion.} Avec un taux d'urbanisation voisin de 59 \% et une concentration de près de la moitié des urbains dans deux villes, le Cameroun illustre une urbanisation rapide et polarisée, source à la fois d'intégration et de déséquilibres.
\end{exoresolu}

\courssub{Méthode 4 : Construire un diagramme (barres, circulaire) à partir de données urbaines}

\begin{methode}[Choisir et construire le bon diagramme]
\begin{enumerate}
\item \textbf{Choisir le type adapté} : diagramme en barres pour comparer des effectifs (population des villes) ou suivre une évolution (population urbaine en 1976, 1987, 2005, projections) ; diagramme circulaire pour des parts d'un tout (répartition de la population urbaine entre catégories de villes).
\item \textbf{Préparer les calculs} : pour un diagramme circulaire, convertir chaque part en angle : $\text{angle} = \dfrac{\text{effectif}}{\text{total}} \times 360^{\circ}$ ; vérifier que la somme des angles vaut $360^{\circ}$.
\item \textbf{Tracer proprement} : axes gradués régulièrement (barres), secteurs tracés au rapporteur (circulaire), barres de même largeur et espacées régulièrement.
\item \textbf{Légender et titrer} : titre complet (quoi, où, quand), légende ou étiquettes sur chaque élément, unités sur les axes, source des données.
\item \textbf{Exploiter} : le diagramme sert l'analyse : citer les valeurs remarquables (maximum, minimum, évolution) dans le commentaire.
\end{enumerate}
\textbf{Réflexes} : ne jamais mélanger les types de données dans un même diagramme ; une évolution se lit en \% (taux d'accroissement), pas seulement en valeurs absolues.
\end{methode}

\begin{exoresolu}[Diagramme circulaire de la répartition de la population urbaine]
\textbf{Énoncé (type Bac)} : La population urbaine du Cameroun (16,5 millions, ordres de grandeur) se répartit ainsi : Yaoundé : 4,2 millions ; Douala : 3,8 millions ; autres capitales régionales : 3,5 millions ; villes moyennes et petites villes : 5 millions. (1) Calculez les angles du diagramme circulaire représentant cette répartition. (2) Tracez le diagramme. (3) Formulez une conclusion.
\tcblower
\textbf{Solution détaillée.}
\textbf{(1) Calcul des angles.} On applique $\text{angle} = \dfrac{\text{effectif}}{\text{total}} \times 360^{\circ}$ avec un total de $4{,}2 + 3{,}8 + 3{,}5 + 5 = 16{,}5$ millions.
\begin{align*}
\text{Yaoundé} &: \dfrac{4{,}2}{16{,}5} \times 360 \approx 91{,}6^{\circ} \\
\text{Douala} &: \dfrac{3{,}8}{16{,}5} \times 360 \approx 82{,}9^{\circ} \\
\text{Autres capitales régionales} &: \dfrac{3{,}5}{16{,}5} \times 360 \approx 76{,}4^{\circ} \\
\text{Villes moyennes et petites villes} &: \dfrac{5}{16{,}5} \times 360 \approx 109{,}1^{\circ}
\end{align*}
Vérification : $91{,}6 + 82{,}9 + 76{,}4 + 109{,}1 = 360^{\circ}$. Le calcul est exact.
\textbf{(2) Tracé.} On trace un cercle ; à partir du rayon vertical (midi), on reporte au rapporteur le secteur de Yaoundé ($91{,}6^{\circ}$), puis celui de Douala ($82{,}9^{\circ}$), puis les autres capitales régionales ($76{,}4^{\circ}$), le dernier secteur complétant le cercle. Chaque secteur reçoit une couleur et une étiquette avec son pourcentage (25,5 \% ; 23 \% ; 21,2 \% ; 30,3 \%).
\textbf{(3) Conclusion.} Le diagramme montre que les deux métropoles (Yaoundé + Douala $\approx 48,5$ \%) pèsent presque autant que l'ensemble des villes moyennes et petites villes, confirmant la concentration de la population urbaine camerounaise au sommet de l'armature urbaine.
\end{exoresolu}

\courssub{Méthode 5 : Rédiger un commentaire de document sur la ville, espace d'intégration}

\begin{methode}[Structurer un commentaire composé en géographie urbaine]
\begin{enumerate}
\item \textbf{Dégager la problématique} à partir du sujet et des documents : ex. « En quoi la ville camerounaise est-elle un espace d'intégration nationale ? » ou « Quels problèmes la croissance urbaine pose-t-elle et quelles solutions ? ».
\item \textbf{Rédiger une introduction en trois temps} : présentation du sujet et des documents ; définition des termes clés (\cle{ville}, \cle{urbanisation}, \cle{intégration nationale}) ; annonce du plan.
\item \textbf{Construire un plan en deux ou trois parties équilibrées}, chaque partie s'appuyant sur les documents ET sur tes connaissances. Exemples de plans types :
\begin{itemize}
\item Ville et intégration : I. Un cadre de brassage des populations, cultures et religions ; II. Un espace où persistent des limites à l'intégration (regroupements communautaires, tensions).
\item Problèmes urbains : I. Les manifestations des problèmes (logement, transport, eau, insalubrité, insécurité) ; II. Les causes (croissance rapide, faiblesse des moyens, incivisme) ; III. Les solutions.
\end{itemize}
\item \textbf{Mobiliser des exemples précis et datés} : quartiers cosmopolites (New Bell à Douala, Briqueterie à Yaoundé), marchés (marché central de Yaoundé, marché Mboppi à Douala), équipements (hôpitaux, universités), politiques (grands travaux, programmes d'assainissement).
\item \textbf{Conclure} : réponse nette à la problématique, puis ouverture (ex. : enjeu des villes secondaires pour rééquilibrer l'armature urbaine).
\end{enumerate}
\textbf{Réflexes} : une idée = un paragraphe = un exemple ; citer explicitement les documents (« le document 2 montre que... ») ; soigner les transitions.
\end{methode}

\begin{exoresolu}[Commentaire : la ville camerounaise, espace d'intégration nationale]
\textbf{Énoncé (type Bac)} : À partir des documents (carte des migrations vers Douala et Yaoundé, photographie du quartier de la Briqueterie à Yaoundé, tableau de la diversité linguistique et religieuse urbaine) et de vos connaissances, montrez que la ville camerounaise est un espace d'intégration nationale, en relevant toutefois ses limites. (Introduction, développement structuré, conclusion.)
\tcblower
\textbf{Solution détaillée (plan rédigé et arguments).}
\textbf{Introduction.} La ville est, selon le programme, un espace de concentration des hommes et des activités. Au Cameroun, pays de plus de 250 groupes ethniques et de plusieurs grandes religions, les villes attirent des populations de toutes les régions : la problématique est de savoir en quoi elles fabriquent l'unité nationale, et quelles en sont les limites.
\textbf{I. La ville, cadre de brassage des populations, des cultures et des religions.}
\begin{itemize}
\item Les migrations convergent vers les villes : la carte montre des flux venus de toutes les régions vers Douala et Yaoundé, attirées par l'emploi et les services.
\item Les quartiers cosmopolites concrétisent ce brassage : à la Briqueterie (Yaoundé) ou à New Bell (Douala), voisins, commerçants et élèves appartiennent à des groupes ethniques et religieux différents ; le document 3 indique la coexistence de chrétiens, de musulmans et d'adeptes des religions traditionnelles, et l'usage des langues nationales mêlé au français et à l'anglais.
\item Les institutions urbaines unifient : écoles et universités brassent les élites de toutes les régions ; les marchés (Mboppi, marché central) sont des lieux d'échanges économiques et culturels ; les mariages mixtes y sont plus fréquents qu'en milieu rural.
\end{itemize}
\textbf{II. Les limites de l'intégration urbaine.}
\begin{itemize}
\item Le regroupement communautaire persiste : certains quartiers ou rues se spécialisent par origine régionale ou religion (quartiers à dominante musulmane comme la Briqueterie), ce qui peut freiner le mélange.
\item Les inégalités socio-spatiales (quartiers planifiés contre quartiers spontanés sous-équipés) créent des fractures qui se superposent parfois aux appartenances communautaires.
\item En période de tensions (crises politiques, conflits), les solidarités communautaires peuvent l'emporter sur la cohabitation urbaine.
\end{itemize}
\textbf{Conclusion.} La ville camerounaise est bien le principal creuset de l'intégration nationale, par le brassage quotidien des populations, des cultures et des religions ; mais cette intégration reste inachevée face aux regroupements communautaires et aux inégalités spatiales, d'où l'enjeu des politiques urbaines d'équipement et de mixité.
\end{exoresolu}

\courssub{Méthode 6 : Construire une argumentation « problèmes et solutions » sur les villes}

\begin{methode}[Traiter un sujet sur les problèmes urbains et leurs solutions]
\begin{enumerate}
\item \textbf{Inventorier les problèmes par catégories} (méthode du classement) : habitat (habitat spontané, surpeuplement), déplacements (congestion de la voirie urbaine, insuffisance des transports collectifs), environnement et santé (insalubrité, déchets, inondations, pollution), approvisionnement (eau potable, électricité délestée), social (chômage, insécurité, incivisme).
\item \textbf{Relier chaque problème à ses causes} : croissance démographique urbaine rapide (environ 4 à 5 \% par an dans les grandes villes, ordre de grandeur), exode rural, faiblesse des ressources financières des communes, absence ou non-application des plans d'urbanisme, comportements d'incivisme (dépôts sauvages d'ordures, constructions anarchiques).
\item \textbf{Associer à chaque problème une solution réaliste et localisée} : viabilisation et élargissement de la voirie, programmes d'assainissement (HYSACAM pour les déchets), restructuration des quartiers spontanés, lotissements, renforcement de la desserte en eau (Camwater) et en électricité (Eneo), transports collectifs (projets de bus et de transport urbain), éducation civique.
\item \textbf{Nuancer} : évaluer les limites des solutions (coût, lenteur, déplacement des problèmes) pour montrer un raisonnement critique.
\item \textbf{Rédiger} avec le vocabulaire exact : \cle{voirie urbaine}, \cle{habitat spontané}, \cle{assainissement}, \cle{restructuration}, \cle{quartier sous-équipé}.
\end{enumerate}
\end{methode}

\begin{exoresolu}[Étude de cas : les problèmes de Yaoundé et les solutions]
\textbf{Énoncé (type Bac)} : « Yaoundé, ville de plus de 4 millions d'habitants, connaît des embouteillages permanents, des coupures d'eau et d'électricité, des quartiers spontanés étendus et des problèmes d'insalubrité et d'insécurité. » (1) Classez ces problèmes en catégories. (2) Expliquez leurs causes. (3) Proposez des solutions argumentées.
\tcblower
\textbf{Solution détaillée.}
\textbf{(1) Classement.} Déplacements : embouteillages liés à une voirie urbaine insuffisante et encombrée. Approvisionnement : déficits en eau potable et délestages électriques. Habitat : extension des quartiers spontanés (Nsam, Mvog-Ada, vallées marécageuses et versants occupés). Environnement et santé : insalubrité (déchets, eaux usées), inondations dans les bas-fonds. Social : insécurité et chômage des jeunes.
\textbf{(2) Causes.} La cause principale est la croissance urbaine très rapide (exode rural + forte croissance naturelle) : la ville croît plus vite que ses équipements. S'y ajoutent l'occupation anarchique du sol (constructions sans permis, sur des sites à risque comme les versants et les bas-fonds), la faiblesse des moyens financiers de la Communauté urbaine, et l'incivisme (dépôts sauvages d'ordures qui bouchent les caniveaux et aggravent les inondations, branchements frauduleux).
\textbf{(3) Solutions.} Pour les déplacements : élargir et entretenir la voirie, créer des voies de contournement et développer un transport collectif structurant. Pour l'eau et l'électricité : renforcer les capacités de production et les réseaux (actions de Camwater et d'Eneo), lutter contre les fraudes. Pour l'habitat : restructurer les quartiers spontanés (ouverture de voies, raccordement aux réseaux) plutôt que démolir, et produire des lotissements viabilisés abordables. Pour l'insalubrité : renforcer la collecte des déchets (HYSACAM), curer les caniveaux, interdire et surveiller la construction dans les zones inondables. Pour l'insécurité : éclairage public, police de proximité, insertion économique des jeunes. Toutes ces solutions supposent aussi l'éducation civique des habitants, car aucune politique ne réussit contre l'incivisme.
\textbf{Conclusion.} Les problèmes de Yaoundé sont ceux d'une ville qui grandit plus vite qu'elle ne s'équipe ; les solutions existent mais exigent des moyens financiers, une planification respectée et la participation citoyenne.
\end{exoresolu}

\begin{attention}[Les 5 erreurs qui coûtent des points au Bac]
\begin{enumerate}
\item \textbf{Confondre « urbanisation » et « croissance démographique ».} L'urbanisation est l'augmentation de la \emph{part} de la population urbaine (en \%), pas la simple augmentation de la population. Calculer un taux d'urbanisation avec la seule population des villes sans la rapporter à la population totale est une erreur de définition et de calcul.
\item \textbf{Décrire sans localiser, ou interpréter pendant la description.} Une description sans présentation du document (nature, lieu, date, échelle) et sans repères spatiaux (au nord, au premier plan, le long du Wouri) perd la moitié des points ; mélanger description et explication brouille la démonstration.
\item \textbf{Oublier la légende, le titre ou le nord sur un croquis.} Un croquis sans titre reprenant le sujet, sans flèche du nord et sans légende ordonnée est sanctionné, même si le fond est correct. Vérifier aussi la cohérence exacte entre les symboles du croquis et ceux de la légende.
\item \textbf{Donner des chiffres faux ou invérifiables.} Inventer une population précise (« Douala : 3 847 215 habitants ») discrédite la copie. Utiliser des ordres de grandeur annoncés comme tels (« environ 3,5 à 4 millions », « plus de 90 \% du commerce extérieur transite par Douala ») et toujours accompagner un pourcentage de sa formule.
\item \textbf{Traiter les problèmes urbains sans causes ni solutions liées, ou avec un vocabulaire vague.} Écrire « les villes sont sales et dangereuses » au lieu d'\cle{insalubrité}, \cle{habitat spontané}, \cle{congestion de la voirie urbaine} relève du langage courant, pas de la géographie. Chaque problème cité doit être relié à une cause et, si le sujet le demande, à une solution localisée (HYSACAM, Camwater, restructuration des quartiers).
\end{enumerate}
\end{attention}

\newpage

\coursec{Exercices}

\courssub{Je vérifie mes connaissances}

\begin{exercice}[QCM : Urbanisation et villes camerounaises]
Parmi les propositions suivantes, une seule est exacte. Recopie la lettre correspondante et justifie brièvement ton choix.
\begin{enumerate}
    \item Le taux d'urbanisation du Cameroun en 2020 est d'environ :
    \begin{itemize}
        \item[A.] 30 \%
        \item[B.] 45 \%
        \item[C.] 58 \%
        \item[D.] 75 \%
    \end{itemize}
    \item La ville qui assure la fonction de capitale politique et administrative du Cameroun est :
    \begin{itemize}
        \item[A.] Douala
        \item[B.] Yaoundé
        \item[C.] Garoua
        \item[D.] Bafoussam
    \end{itemize}
    \item Le phénomène d'étalement urbain spontané, souvent non viabilisé, se traduit par l'apparition de :
    \begin{itemize}
        \item[A.] Quartiers résidentiels de standing
        \item[B.] Zones industrielles planifiées
        \item[C.] Quartiers spontanés (bidonvilles)
        \item[D.] Parcs urbains protégés
    \end{itemize}
    \item La ville de Douala doit son rang de métropole économique principalement à :
    \begin{itemize}
        \item[A.] Sa fonction de capitale politique
        \item[B.] Son port autonome et sa zone industrielle
        \item[C.] Son aéroport international unique
        \item[D.] Sa position au centre du pays
    \end{itemize}
\end{enumerate}
\end{exercice}

\begin{exercice}[Vrai ou Faux justifié]
Indique si chaque affirmation est VRAIE ou FAUSSE, puis justifie ta réponse par une phrase précise en mobilisant tes connaissances sur le Cameroun.
\begin{enumerate}
    \item La ville de Maroua est la plus grande ville de la région de l'Extrême-Nord et un carrefour commercial vers le Tchad et le Nigeria.
    \item L'exode rural est la cause unique de la croissance urbaine au Cameroun ; l'accroissement naturel n'y joue aucun rôle.
    \item Le « quartier cosmopolite » se caractérise par la cohabitation de populations d'origines ethniques, religieuses et culturelles diverses.
    \item Le Plan d'Aménagement et de Développement Urbain (PADU) est un document d'urbanisme obligatoire pour toute commune de plus de 50 000 habitants.
\end{enumerate}
\end{exercice}

\begin{exercice}[Définitions et concepts clés]
Donne une définition géographique précise (2 à 3 lignes) pour chacun des termes suivants en les situant dans le contexte camerounais :
\begin{enumerate}
    \item \cle{Urbanisation}
    \item \cle{Ville primatiale} (ou \cle{macrocéphalie urbaine})
    \item \cle{Quartier spontané} (ou \cle{habitat précaire})
    \item \cle{Polarisation des services}
    \item \cle{Voirie urbaine}
\end{enumerate}
\end{exercice}

\begin{exercice}[Calculs directs : Évolution de la population urbaine]
Le tableau ci-dessous donne l'évolution de la population totale et urbaine du Cameroun (chiffres arrondis, sources : BUCREP/INS, projections).

\begin{center}
\begin{tabularx}{\linewidth}{|l|X|X|X|X|}\hline
\textbf{Année} & \textbf{1976} & \textbf{1987} & \textbf{2005} & \textbf{2020 (proj.)} \\ \hline
Population totale (millions) & 7,7 & 10,5 & 17,5 & 26,5 \\ \hline
Population urbaine (millions) & 1,5 & 3,8 & 8,8 & 15,4 \\ \hline
\end{tabularx}
\end{center}

\begin{enumerate}
    \item Calcule le taux d'urbanisation (part de la population urbaine dans la population totale) pour chacune des quatre dates. Arrondis au dixième de pourcent.
    \item Calcule le taux de croissance annuel moyen de la population urbaine entre 1987 et 2005 (formule : $t = \left( \frac{P_f}{P_i} \right)^{\frac{1}{n}} - 1$). Exprime le résultat en \%.
    \item En te basant sur tes résultats, complète la phrase : « Entre 1976 et 2020, la population urbaine a été multipliée par environ \ldots , tandis que la population totale n'a été multipliée que par \ldots . »
\end{enumerate}
\end{exercice}

\courssub{Je m'entraîne}

\begin{exercice}[Analyse d'un croquis : La répartition des grandes villes du Cameroun]
On te propose un croquis simplifié de la répartition des principales villes camerounaises (seuil : 100 000 hab. en 2020). \emph{(Le croquis est à réaliser par l'élève ou fourni par le professeur : contour du pays, 10 régions, points proportionnels pour Yaoundé, Douala, Bafoussam, Garoua, Maroua, Ngaoundéré, Bertoua, Ebolowa, Bamenda, Kumba, Limbe, Kribi, Foumban, Dschang).}

\begin{enumerate}
    \item Légende ce croquis en respectant la sémiologie graphique (figé ponctuel proportionnel pour la population, flèches pour les axes de communication majeurs, hachures pour les zones de forte densité urbaine).
    \item Décris la répartition des grandes villes (concentration, alignements, vides).
    \item Explique cette répartition en reliant : le réseau hydrographique (Sanaga, Bénoué, Logone), le réseau de transport (Transcamerounais, route de l'Est, route du Nord), le relief (massif de l'Adamaoua, plaine du Bénoué) et l'héritage colonial/historique.
\end{enumerate}
\end{exercice}

\begin{exercice}[Étude de cas : Le quartier cosmopolite de « New Bell » à Douala]
Lis le texte suivant, extrait d'une enquête de terrain : 
\emph{« New Bell, arrondissement de Douala II, est l'un des plus anciens quartiers de la ville. Né de l'extension spontanée vers l'est du centre colonial, il accueille aujourd'hui plus de 200 000 habitants. On y croise des Bamileke (commerçants), des Douala (autochtones), des Bassa, des Haoussas (bouviers, vendeurs de bétail), des Nigérians (pêcheurs, commerçants), des Maliens (tailleurs, bijoutiers) et des populations venues des régions anglophones (Nord-Ouest, Sud-Ouest). Les langues véhiculées sont le français, l'anglais, le pidgin, le douala, le bamiléké, l'haoussa. On y trouve des églises (catholiques, protestantes, de réveil), des mosquées, des temples traditionnels. L'activité économique est dominée par le commerce informel (marché de New Bell, étals le long de la voirie), la petite industrie (menuiserie, soudure, mécanique) et les services (transport, restauration). La voirie est majoritairement non bitumée, l'évacuation des eaux usées défaillante, l'accès à l'eau potable et à l'électricité inégal. »}

\begin{enumerate}
    \item Identifie dans le texte trois éléments qui prouvent que New Bell est un \cle{quartier cosmopolite}.
    \item Montre que ce quartier joue un rôle d'\cle{intégration nationale} (brassage des populations, des cultures, des religions).
    \item Relevez deux indices de \cle{problèmes urbains} (insalubrité, voirie, habitat) mentionnés dans le texte.
    \item Propose deux actions concrètes que la Commune d'arrondissement pourrait engager pour améliorer le cadre de vie (voirie, assainissement, équipements collectifs).
\end{enumerate}
\end{exercice}

\begin{exercice}[Construction d'un diagramme : L'explosion urbaine de Yaoundé et Douala]
Le tableau suivant donne l'évolution de la population des deux métropoles (en millions d'habitants, estimations).

\begin{center}
\begin{tabularx}{\linewidth}{|l|X|X|X|X|X|}\hline
\textbf{Ville} & \textbf{1976} & \textbf{1987} & \textbf{2005} & \textbf{2010} & \textbf{2020} \\ \hline
Yaoundé & 0,30 & 0,70 & 1,70 & 2,10 & 3,90 \\ \hline
Douala & 0,45 & 1,00 & 2,00 & 2,40 & 3,80 \\ \hline
\end{tabularx}
\end{center}

\begin{enumerate}
    \item Construis, sur un même graphique (axes orthonormés), deux histogrammes groupés (ou un diagramme en barres groupées) représentant l'évolution de Yaoundé et de Douala. Respecte les règles : titre, axes gradués et légendés (unité), légende des séries, source.
    \item Compare la dynamique des deux villes sur la période 1976-2020 (rythme, écarts).
    \item Quelles conséquences cette croissance rapide a-t-elle sur la \cle{voirie urbaine}, le \cle{logement} et l'\cle{approvisionnement en eau/électricité} ? Donne un exemple précis pour chaque problème.
\end{enumerate}
\end{exercice}

\begin{exercice}[La ville, centre d'impulsion économique : Le cas de la filière bois à Douala]
La ville de Douala concentre l'essentiel de la transformation industrielle du bois camerounais (sciage, placage, contreplaqué, meubles). Les grumes arrivent par route (depuis l'Est et le Sud) ou par fleuve (Wouri, Sanaga) vers les zones industrielles de Bassa, Bonabéri et Yassa.

\begin{enumerate}
    \item En t'appuyant sur cet exemple, explique comment la ville agit comme un \cle{centre d'impulsion de l'activité économique} (transformation, valeur ajoutée, emploi).
    \item Identifie les \cle{services polarisés} par cette activité industrielle (transport, banque, assurance, portuaire, formation, maintenance).
    \item Pourquoi la localisation de ces usines à Douala (ville portuaire) est-elle un avantage comparatif majeur pour l'exportation ? Cite le principal port d'exportation.
    \item Cite deux problèmes environnementaux ou urbains générés par cette concentration industrielle dans la ville.
\end{enumerate}
\end{exercice}

\courssub{Je me prépare au Bac}

\begin{exercice}[Sujet type Bac : Exploitation de documents — « L'urbanisation du Cameroun : dynamiques et déséquilibres »]
\emph{Durée conseillée : 1 h 30. Les documents 1, 2 et 3 sont fournis.}

\textbf{Document 1 : Carte de l'urbanisation du Cameroun (2020)}
\emph{(Carte schématique à analyser : 10 régions, villes > 100 000 hab. en points proportionnels, axes routiers principaux, zones de forte densité rurale, frontière internationale).}

\textbf{Document 2 : Évolution de la population urbaine et du taux d'urbanisation (1960-2020)}
\begin{center}
\begin{tabularx}{\linewidth}{|l|X|X|X|X|X|X|X|}\hline
\textbf{Année} & \textbf{1960} & \textbf{1976} & \textbf{1987} & \textbf{2005} & \textbf{2010} & \textbf{2015} & \textbf{2020} \\ \hline
Pop. urbaine (millions) & 0,7 & 1,5 & 3,8 & 8,8 & 11,2 & 13,5 & 15,4 \\ \hline
Taux d'urbanisation (\%) & 12,5 & 19,5 & 36,2 & 50,3 & 53,5 & 55,8 & 58,1 \\ \hline
\end{tabularx}
\end{center}
Source : BUCREP, INS, World Urbanization Prospects (ONU).

\textbf{Document 3 : Extrait du discours du Ministre de l'Habitat et du Développement Urbain (MINHDU), 2022}
\emph{« Le déficit en logements décents au Cameroun est estimé à plus de 2,5 millions d'unités. 60 \% de la population urbaine vit dans des quartiers spontanés non viabilisés. L'insalubrité, l'insécurité foncière, la saturation de la voirie et les inondations récurrentes (Douala, Yaoundé, Limbe, Buea) sont les visages de cette crise urbaine. Le Plan National de Développement (PND 2020-2030) et la Stratégie Nationale de Développement (SND30) font de la ville durable une priorité : production de 50 000 logements/an, programmes d'assainissement (PADY, PADDouala), réforme foncière, mobilité urbaine (bus BRT, tramway projeté). »}

\textbf{Travail demandé :}
\begin{enumerate}
    \item \textbf{Analyse du Document 1 (4 pts) :} Présente la carte (titre, échelle, source, légende). Décris la répartition des grandes villes. Explique les facteurs de cette répartition (histoire, relief, hydrographie, axes de transport).
    \item \textbf{Exploitation du Document 2 (4 pts) :} Construis un graphique adapté (courbe ou barres) montrant l'évolution du taux d'urbanisation de 1960 à 2020. Dégage deux enseignements majeurs sur la dynamique urbaine camerounaise.
    \item \textbf{Synthèse Documents 1, 2 et 3 (7 pts) :} Montre que l'urbanisation rapide du Cameroun s'accompagne d'une « crise urbaine » majeure. Structure ta réponse en trois parties : 1. Une urbanisation rapide et macrocéphale (Douala/Yaoundé) ; 2. Les manifestations de la crise (logement, voirie, insalubrité, services) ; 3. Les réponses publiques (PND, SND30, programmes sectoriels).
    \item \textbf{Mise en perspective (5 pts) :} « La ville camerounaise, moteur de l'intégration nationale, peut-elle devenir une ville durable ? » Discute en mobilisant tes connaissances (décentralisation, gouvernance urbaine, financement, changement climatique).
\end{enumerate}
\end{exercice}

\begin{exercice}[Sujet type Bac : Dissertation — « La ville camerounaise : espace d'intégration nationale, mais espace en crise »]
\emph{Durée conseillée : 2 h. Sujet à traiter en dissertation structurée (Introduction, Développement en 2 ou 3 parties équilibrées, Conclusion).}

\textbf{Consignes :}
\begin{enumerate}
    \item Analyse les termes du sujet : « ville camerounaise », « espace d'intégration nationale », « espace en crise ». Définit-les.
    \item Problématise : Comment la ville, lieu par excellence du brassage et de l'activité économique, produit-elle simultanément des dysfonctionnements majeurs qui menacent la cohésion sociale et le développement durable ?
    \item Construis un plan détaillé (annonce des grandes parties et sous-parties).
    \item Rédige le développement en t'appuyant sur des \textbf{exemples précis et chiffrés} (Douala, Yaoundé, Maroua, Bamenda, Kribi, etc.) : brassage ethnique/religieux, marchés intercommunautaires, rôle des chefferies urbaines, mais aussi quartiers spontanés (Nkolbisson, New Bell, Mvog-Ada), embouteillages chroniques, pénuries d'eau (Eau du Cameroun/SCE), inondations, insécurité (coupeurs de route, bandes urbaines), gestion des déchets (HYSACAM), exclusion sociale.
    \item Conclus en ouvrant sur les politiques d'aménagement du territoire (SND30, villes intermédiaires, décentralisation, ville intelligente/durable).
\end{enumerate}
\end{exercice}

\begin{exercice}[Sujet type Bac : Travaux Pratiques (TP) — Construction d'un croquis de quartier cosmopolite]
\emph{Durée : 1 h. Matériel : Plan cadastral simplifié d'un quartier type (quadrillage orthogonal, voirie principale/secondaire, parcelles, équipements : marché, école, mosquée, église, centre de santé, terrain vague, zone marécageuse), papier calque, crayons de couleur, règle, équerre.}

\textbf{Consignes :}
\begin{enumerate}
    \item \textbf{Analyse du document source (2 pts) :} Identifie les éléments du plan fourni (échelle, orientation, types d'occupation du sol, équipements collectifs, voirie).
    \item \textbf{Construction du croquis (10 pts) :} Sur ton calque, réalise le croquis du « Quartier cosmopolite de \emph{Mvog-Mbi} (Yaoundé) » en respectant \textbf{rigoureusement} les étapes suivantes :
    \begin{itemize}
        \item Cadre, titre centré, flèche du Nord, échelle graphique.
        \item Légende organisée (hiérarchisée) \textbf{avant} le tracé : 
        \begin{itemize}
            \item \textbf{Habitat} : Habitat planifié/viabilisé (couleur claire), Habitat spontané/précaire (hachures rouges), Habitat rural résiduel.
            \item \textbf{Activités} : Marché principal (symbole), Commerce de rue/étals (points), Petites industries/artisanat (carreaux), Zone de stationnement/parc auto.
            \item \textbf{Équipements/Services} : Enseignement (école, lycée), Santé (CMA, hôpital), Culte (mosquée, église, temple), Administration (chefferie, mairie d'arrondissement), Espaces verts/loisirs.
            \item \textbf{Voirie et Réseaux} : Route bitumée (trait continu épais), Piste/voie non bitumée (trait pointillé), Réseau d'assainissement visible (caniveaux), Point d'eau/borne fontaine.
            \item \textbf{Contraintes/Risques} : Zone inondable/marécageuse (hachures bleues), Décharge sauvage, Érosion/gully.
        \end{itemize}
        \item Tracé soigné au crayon de couleur, respect de la sémiologie (surface, linéaire, ponctuel).
    \end{itemize}
    \item \textbf{Commentaire géographique (8 pts) :} Rédige un commentaire structuré (300 mots env.) de ton croquis pour montrer : 
    \begin{itemize}
        \item L'organisation morphologique du quartier (centre/peripherie, densité).
        \item La mixité sociale et fonctionnelle (cosmopolitisme : lieux de culte, marché, écoles).
        \item Les dysfonctionnements urbains visibles (voirie, assainissement, habitat précaire, risques).
        \item Les potentialités et pistes d'aménagement (restructuration, viabilisation, équipements manquants).
    \end{itemize}
\end{enumerate}
\end{exercice}

\newpage

\coursec{Corrigés des exercices}

\coursec{Leçon 3 : Les villes du Cameroun — Exercices corrigés et approfondis}

\courssub{Entraînement et préparation au Baccalauréat}

\begin{corrige}[Exercice 1 — QCM : Urbanisation et villes camerounaises]
\textbf{1. Réponse C — \num{58} \%.} Selon les projections du BUCREP/INS et de l'ONU, le taux d'urbanisation du Cameroun atteint environ \num{58} \% en 2020 (contre environ \num{12} \% en 1960). Le Cameroun est devenu un pays majoritairement urbain au milieu des années 2000.

\textbf{2. Réponse B — Yaoundé.} Yaoundé est la capitale politique et administrative du Cameroun (siège de la présidence, des ministères, des ambassades). Attention à ne pas confondre avec Douala, qui est la capitale \emph{économique}.

\textbf{3. Réponse C — Quartiers spontanés (bidonvilles).} L'étalement urbain rapide et non maîtrisé produit des quartiers spontanés, sans voirie viabilisée ni réseaux d'eau et d'assainissement (ex. : Nkolbisson à Yaoundé, New Bell à Douala).

\textbf{4. Réponse B — Son port autonome et sa zone industrielle.} Le Port Autonome de Douala (PAD) traite l'essentiel du commerce extérieur du pays (et une partie de celui du Tchad et de la Centrafrique), et la zone industrielle de Douala-Bassa/Bonabéri concentre la majorité des unités industrielles nationales.

\textbf{Point de vigilance :} au Bac, la justification est exigée même pour un QCM ; une lettre seule ne vaut pas tous les points.
\end{corrige}

\begin{corrige}[Exercice 2 — Vrai ou Faux justifié]
\textbf{1. VRAI.} Maroua (environ \num{500 000} à \num{600 000} habitants, ordre de grandeur) est le chef-lieu de la région de l'Extrême-Nord et un carrefour commercial majeur sur l'axe Ngaoundéré–Maroua–Kousseri, vers le Tchad (N'Djaména) et le Nigeria (via Banki).

\textbf{2. FAUX.} La croissance urbaine résulte de \emph{deux} mécanismes combinés : l'exode rural \textbf{et} l'accroissement naturel des populations déjà urbaines (natalité élevée, mortalité en baisse). Depuis les années 1990, l'accroissement naturel est même devenu la première composante de la croissance de nombreuses villes camerounaises.

\textbf{3. VRAI.} Le quartier cosmopolite (ex. : New Bell à Douala, Mvog-Mbi à Yaoundé) accueille des populations d'origines ethniques, régionales, nationales (Nigérians, Maliens, Tchadiens) et religieuses diverses, qui y cohabitent et y échangent : c'est un espace de brassage.

\textbf{4. FAUX.} Le PDAU (Projet de Développement Urbain, ex-PADY) est un document de planification urbaine élaboré pour les grandes agglomérations (Douala, Yaoundé, puis d'autres villes), mais il n'existe pas de seuil légal de \num{50 000} habitants le rendant « obligatoire pour toute commune » ; beaucoup de communes camerounaises ne disposent d'aucun document d'urbanisme à jour, ce qui aggrave l'anarchie foncière.

\textbf{Point de vigilance :} une justification vague (« c'est vrai car c'est dans le cours ») est sanctionnée ; il faut mobiliser un fait précis (ville, chiffre, mécanisme).
\end{corrige}

\begin{corrige}[Exercice 3 — Définitions et concepts clés]
\textbf{1. Urbanisation :} processus de concentration croissante de la population dans les villes et d'extension des modes de vie urbains. Au Cameroun, elle se mesure par le taux d'urbanisation, passé d'environ \num{12} \% (1960) à environ \num{58} \% (2020).

\textbf{2. Ville primatiale (macrocéphalie urbaine) :} ville dont le poids démographique et économique écrase celui des autres villes du réseau urbain. Au Cameroun, Douala et Yaoundé (environ \num{3,8} et \num{3,9} millions d'habitants en 2020) concentrent à elles deux près de la moitié de la population urbaine : on parle de \cle{macrocéphalie} à deux têtes.

\textbf{3. Quartier spontané (habitat précaire) :} espace urbain construit sans plan d'urbanisme ni titre foncier sécurisé, sur des terrains souvent inadaptés (marais, pentes), dépourvu de voirie bitumée, d'eau potable et d'assainissement. Exemples : New Bell (Douala), Mvog-Ada (Yaoundé).

\textbf{4. Polarisation des services :} concentration dans une ville des services supérieurs (banques, hôpitaux de référence, universités, administrations, sièges d'entreprises) qui attirent et desservent une vaste aire d'influence (hinterland). Douala polarise ainsi tout le littoral et une partie de la CEMAC.

\textbf{5. Voirie urbaine :} ensemble des voies de circulation d'une ville (avenues bitumées, rues, pistes, ronds-points, ponts) et de leurs dépendances (caniveaux, trottoirs, éclairage). Au Cameroun, la voirie est souvent insuffisante et dégradée face à la croissance du parc automobile (embouteillages de Douala et Yaoundé).
\end{corrige}

\begin{corrige}[Exercice 4 — Calculs directs : Évolution de la population urbaine]
\textbf{1. Taux d'urbanisation} : $T = \dfrac{P_{\text{urb}}}{P_{\text{tot}}} \times 100$.
\begin{align*}
1976 &: \frac{\num{1,5}}{\num{7,7}}\times 100 = \num{19,5}\,\% \\
1987 &: \frac{\num{3,8}}{\num{10,5}}\times 100 = \num{36,2}\,\% \\
2005 &: \frac{\num{8,8}}{\num{17,5}}\times 100 = \num{50,3}\,\% \\
2020 &: \frac{\num{15,4}}{\num{26,5}}\times 100 = \num{58,1}\,\%
\end{align*}

\textbf{2. Taux de croissance annuel moyen 1987--2005} : $n = 2005 - 1987 = 18$ ans.
\[
t = \left(\frac{\num{8,8}}{\num{3,8}}\right)^{\frac{1}{18}} - 1 = (\num{2,316})^{0,0556} - 1 \approx \num{1,0478} - 1 = \num{0,0478}
\]
Soit \textbf{$t \approx \num{4,8}\,\%$ par an}, un rythme très élevé (la population urbaine double en environ 15 ans à ce rythme).

\textbf{3. Phrase complétée :} « Entre 1976 et 2020, la population urbaine a été multipliée par environ \textbf{10} ($\num{15,4}/\num{1,5} \approx \num{10,3}$), tandis que la population totale n'a été multipliée que par \textbf{3,4} ($\num{26,5}/\num{7,7} \approx \num{3,4}$). »

\textbf{Point de vigilance :} ne pas confondre taux de variation global et taux annuel moyen ; penser à convertir le résultat en pourcentage et à préciser l'arrondi.
\end{corrige}

\begin{corrige}[Exercice 5 — Analyse d'un croquis : La répartition des grandes villes du Cameroun]
\textbf{1. Légendage attendu :}
\begin{itemize}
    \item Figé ponctuel proportionnel (cercles de surface proportionnelle à la population) : Douala et Yaoundé (grands cercles, $\approx$ \num{3,8}--\num{3,9} M), puis Bafoussam, Garoua, Maroua, Bamenda, Ngaoundéré (cercles moyens), Bertoua, Ebolowa, Kumba, Limbe, Kribi, Foumban, Dschang (petits cercles).
    \item Flèches pleines : axe Douala–Yaoundé ; route du Nord (Yaoundé–Ngaoundéré–Garoua–Maroua) ; route de l'Est (Yaoundé–Bertoua) ; Transcamerounais (Ngaoundéré–Yaoundé).
    \item Hachures : zones de forte densité (Ouest, Littoral, Centre, extrême nord soudanien).
\end{itemize}

\textbf{2. Description :} les grandes villes se concentrent dans le \textbf{Sud} (Littoral, Centre, Ouest, Sud-Ouest) et le long des \textbf{axes de communication} : alignement Douala–Yaoundé–Bafoussam, alignement nordique Ngaoundéré–Garoua–Maroua. De grands \textbf{vides} urbains existent : plateau de l'Adamaoua (hors Ngaoundéré), Est forestier (hors Bertoua), Nord anglophone montagneux peu urbanisé.

\textbf{3. Explication :}
\begin{itemize}
    \item \emph{Hydrographie} : Douala sur l'estuaire du Wouri (port), Garoua et Maroua dans les plaines de la Bénoué et du Logone (agriculture, élevage, commerce).
    \item \emph{Transports} : le chemin de fer Transcamerounais a fait naître ou grossir Ngaoundéré et Yaoundé ; les routes structurent les alignements.
    \item \emph{Relief} : les hautes terres de l'Ouest (climat frais, sols fertiles) portent de fortes densités rurales qui alimentent Bafoussam, Dschang, Bamenda ; l'Adamaoua, barrière de relief, est un vide.
    \item \emph{Histoire} : héritage colonial (Douala, port de Douala créé par les Allemands ; Yaoundé, station coloniale devenue capitale) et anciens centres précoloniaux (Foumban, Maroua).
\end{itemize}
\textbf{Point de vigilance :} un croquis sans titre, sans flèche du nord ou avec une légende non hiérarchisée perd des points, même si le fond est juste.
\end{corrige}

\begin{corrige}[Exercice 6 — Étude de cas : Le quartier cosmopolite de New Bell à Douala]
\textbf{1. Trois éléments prouvant le cosmopolitisme :}
\begin{itemize}
    \item Diversité ethnique et nationale : Bamiléké, Douala, Bassa, Haoussas, Nigérians, Maliens, anglophones du Nord-Ouest et du Sud-Ouest.
    \item Pluralité linguistique : français, anglais, pidgin, douala, bamiléké, haoussa.
    \item Pluralité religieuse : églises catholiques, protestantes et de réveil, mosquées, temples traditionnels.
\end{itemize}

\textbf{2. Rôle d'intégration nationale :} New Bell est un cadre de \cle{brassage} : des populations de toutes les régions du Cameroun et de l'étranger y vivent, y commercent (marché de New Bell), y parlent des langues communes (français, pidgin) et y pratiquent leurs cultes côte à côte. Les échanges économiques intercommunautaires (bétail haoussa, commerce bamiléké, pêche nigériane) tissent des solidarités qui dépassent les appartenances d'origine : le quartier fabrique une identité urbaine camerounaise commune.

\textbf{3. Deux indices de problèmes urbains :}
\begin{itemize}
    \item Voirie majoritairement non bitumée et évacuation des eaux usées défaillante (insalubrité, risques d'inondation et d'épidémies).
    \item Accès inégal à l'eau potable et à l'électricité ; habitat issu d'une extension spontanée (précarité, insécurité foncière).
\end{itemize}

\textbf{4. Deux actions concrètes :}
\begin{itemize}
    \item Bitumage des rues principales et curage/élargissement des caniveaux d'évacuation des eaux pluviales (programme type PDAU / PAD Douala).
    \item Raccordement prioritaire au réseau d'eau potable (bornes-fontaines, branchements sociaux) et création d'équipements collectifs (centre de santé, latrines publiques, aire de marché aménagée).
\end{itemize}
\end{corrige}

\begin{corrige}[Exercice 7 — Construction d'un diagramme : L'explosion urbaine de Yaoundé et Douala]
\textbf{1. Diagramme en barres groupées attendu :}

\begin{popfigure}
\begin{center}
\begin{tikzpicture}
\begin{axis}[
    ybar, bar width=7pt,
    width=0.95\linewidth, height=6.5cm,
    xlabel={Année}, ylabel={Population (millions d'habitants)},
    symbolic x coords={1976,1987,2005,2010,2020},
    xtick=data, ymin=0, ymax=4.5,
    legend style={at={(0.02,0.97)},anchor=north west,font=\small},
    ymajorgrids=true, grid style={popline!40},
]
\addplot[fill=popblue,draw=popdark] coordinates {(1976,0.30) (1987,0.70) (2005,1.70) (2010,2.10) (2020,3.90)};
\addplot[fill=poporange,draw=popdark] coordinates {(1976,0.45) (1987,1.00) (2005,2.00) (2010,2.40) (2020,3.80)};
\legend{Yaoundé, Douala}
\end{axis}
\end{tikzpicture}
\end{center}
\legende{Évolution comparée des populations de Yaoundé et Douala, 1976--2020 (millions d'habitants). Source : BUCREP/INS, estimations.}
\end{popfigure}

\textbf{2. Comparaison :} Douala part devant (\num{0,45} M contre \num{0,30} M en 1976) et garde l'avantage jusqu'en 2010 (\num{2,4} M contre \num{2,1} M). Mais Yaoundé croît plus vite (multipliée par 13 entre 1976 et 2020, contre 8,4 pour Douala) et dépasse Douala vers 2020 (\num{3,9} M contre \num{3,8} M). L'écart se réduit donc continûment : les deux métropoles forment une macrocéphalie à deux têtes.

\textbf{3. Conséquences :}
\begin{itemize}
    \item \emph{Voirie} : saturation des axes (embouteillages chroniques sur le pont sur le Wouri et le boulevard de l'Unité à Yaoundé), dégradation rapide des chaussées.
    \item \emph{Logement} : déficit estimé à plus de \num{2,5} millions d'unités au niveau national ; extension des quartiers spontanés (Nkolbisson, Mvog-Ada à Yaoundé).
    \item \emph{Eau/électricité} : pénuries et délestages ; à Douala, de nombreux quartiers dépendent des forages et des bornes-fontaines faute de raccordement au réseau de la Camerounaise des Eaux (CDE).
\end{itemize}
\textbf{Point de vigilance :} un graphique sans unité sur l'axe vertical ou sans source est incomplet au Bac.
\end{corrige}

\begin{corrige}[Exercice 8 — La ville, centre d'impulsion économique : la filière bois à Douala]
\textbf{1. Centre d'impulsion :} Douala transforme une matière première brute (grumes venues de l'Est et du Sud) en produits à plus forte \cle{valeur ajoutée} (sciages, placages, contreplaqués, meubles). Cette transformation crée des emplois industriels et tertiaires, des revenus fiscaux, et stimule les campagnes d'approvisionnement (demande de grumes, transport) : la ville impulse l'activité de tout son arrière-pays.

\textbf{2. Services polarisés :} transport et logistique (camionnage, manutention), banques et assurances (financement, couverture des risques), services portuaires (consignation, transit, douanes), formation technique (maintenance, mécanique du bois), services d'ingénierie et de réparation.

\textbf{3. Avantage comparatif :} la localisation portuaire supprime la rupture de charge coûteuse entre usine et navire : les produits finis sont exportés directement par le \textbf{Port Autonome de Douala}, principal port d'exportation du Cameroun (bois, cacao, coton, bananes). Cela réduit les coûts logistiques et rend la production compétitive à l'international.

\textbf{4. Deux problèmes :}
\begin{itemize}
    \item Pollution de l'air et des eaux (rejets industriels, sciures, solvants) dans les zones de Bassa et Bonabéri, dégradation de l'estuaire du Wouri.
    \item Congestion urbaine : flux de camions lourds dégradant la voirie et aggravant les embouteillages ; pression foncière industrielle concurrençant l'habitat.
\end{itemize}
\end{corrige}

\begin{corrige}[Exercice 9 — Sujet type Bac : « L'urbanisation du Cameroun : dynamiques et déséquilibres »]
\textbf{Barème indicatif :} Q1 : 4 pts ; Q2 : 4 pts ; Q3 : 7 pts ; Q4 : 5 pts.

\textbf{1. Analyse du Document 1 (4 pts).} \emph{Présentation (1 pt) :} carte de l'urbanisation du Cameroun en 2020, échelle nationale, sources BUCREP/INS ; légende : points proportionnels (villes > \num{100 000} hab.), lignes (axes routiers), hachures (fortes densités). \emph{Description (1,5 pt) :} concentration des grandes villes au Sud (Douala, Yaoundé, Bafoussam, Bamenda) et le long des axes ; alignement nordique Ngaoundéré–Garoua–Maroua ; vides sur l'Adamaoua et l'Est. \emph{Explication (1,5 pt) :} héritage colonial (port de Douala, capitale Yaoundé), relief (hautes terres de l'Ouest attractives, Adamaoua répulsif), hydrographie (Wouri, Bénoué, Logone), réseau de transport (Transcamerounais, route du Nord).

\textbf{2. Exploitation du Document 2 (4 pts).} Graphique attendu : courbe (ou barres) du taux d'urbanisation, axes légendés (années ; taux en \%), titre, source (2 pts). \emph{Enseignements (2 pts) :} (a) le taux a été multiplié par plus de 4,5 entre 1960 (\num{12,5} \%) et 2020 (\num{58,1} \%) : urbanisation très rapide ; (b) le seuil des \num{50} \% est franchi vers 2005 : le Cameroun est devenu un pays majoritairement urbain, et la population urbaine (\num{0,7} M $\rightarrow$ \num{15,4} M) croît plus vite que le taux, signe d'une explosion démographique urbaine absolue.

\textbf{3. Synthèse (7 pts) — plan attendu :}
\begin{enumerate}
    \item \emph{Une urbanisation rapide et macrocéphale (2 pts) :} taux passé de \num{12,5} \% à \num{58,1} \% ; Douala et Yaoundé ($\approx$ \num{3,8}--\num{3,9} M chacune) concentrent près de la moitié des urbains ; exode rural + accroissement naturel.
    \item \emph{Les manifestations de la crise (3 pts) :} déficit de \num{2,5} millions de logements ; \num{60} \% des urbains en quartiers spontanés non viabilisés ; insalubrité, insécurité foncière, voirie saturée, inondations récurrentes (Douala, Yaoundé, Limbe, Buea).
    \item \emph{Les réponses publiques (2 pts) :} PND 2020-2030 et SND30 (ville durable, \num{50 000} logements/an), programmes d'assainissement (PDAU pour Yaoundé et Douala), réforme foncière, mobilité urbaine (BRT, projet de tramway).
\end{enumerate}

\textbf{4. Mise en perspective (5 pts).} La ville camerounaise est bien un moteur d'intégration (brassage ethnique et religieux, marchés intercommunautaires, langues communes). Mais sa durabilité est conditionnée par : la \emph{décentralisation} effective (transfert de ressources aux communes), la \emph{gouvernance urbaine} (planification, sécurisation foncière), le \emph{financement} (fiscalité locale, partenariats public-privé) et l'adaptation au \emph{changement climatique} (inondations, zones à risque). Conclusion nuancée : la ville durable est possible si l'investissement suit la croissance démographique, ce qui n'est pas encore le cas.

\textbf{Points de vigilance :} citer les chiffres des documents (le correcteur les attend) ; ne pas recopier le Document 3 sans le relier aux Documents 1 et 2 ; annoncer le plan de la synthèse.
\end{corrige}

\begin{corrige}[Exercice 10 — Sujet type Bac : Dissertation « La ville camerounaise : espace d'intégration nationale, mais espace en crise »]
\textbf{1. Analyse des termes.} \emph{Ville camerounaise} : agglomération polarisant une aire d'influence (de Douala, $\approx$ \num{3,8} M hab., aux chefs-lieux de département). \emph{Espace d'intégration nationale} : lieu de brassage des populations, cultures et religions, où se construit une identité commune. \emph{Espace en crise} : accumulation de dysfonctionnements (habitat, voirie, eau, sécurité) qui menacent la cohésion.

\textbf{2. Problématique :} comment le même espace urbain peut-il être à la fois le principal creuset de l'unité nationale et le lieu concentré des fractures sociales et des défaillances de développement ?

\textbf{3. Plan détaillé attendu (dialectique) :}
\begin{itemize}
    \item \textbf{I. La ville, creuset de l'intégration nationale.} (a) Brassage des populations : quartiers cosmopolites (New Bell à Douala, Mvog-Mbi à Yaoundé), migrations interrégionales, langues véhiculaires communes (français, anglais, pidgin). (b) Brassage économique et culturel : marchés intercommunautaires (marché de New Bell, marché de Mokolo à Yaoundé, carrefour commercial de Maroua vers le Tchad et le Nigeria), rôle fédérateur des chefferies urbaines et des associations de ressortissants, cohabitation religieuse (églises, mosquées, temples). (c) Impulsion économique : concentration des industries (Douala), des services supérieurs et des administrations (Yaoundé).
    \item \textbf{II. Mais une ville en crise profonde.} (a) Habitat et foncier : \num{60} \% d'urbains en quartiers spontanés, déficit de \num{2,5} M de logements (Nkolbisson, Mvog-Ada, New Bell). (b) Services défaillants : voirie saturée (embouteillages du pont sur le Wouri), pénuries d'eau (CDE) et délestages électriques, gestion des déchets (HYSACAM débordé), inondations (Douala, Limbe, Buea). (c) Crise sociale : chômage et informalité, insécurité (bandes urbaines), exclusion et inégalités spatiales criantes entre quartiers viabilisés (Bastos, Bonanjo) et quartiers précaires.
    \item \textbf{III. (Synthèse) Des pistes pour réconcilier intégration et durabilité :} SND30 et PND 2020-2030, programmes PDAU, développement des villes intermédiaires (Kribi, Ngaoundéré, Bafoussam) pour décongestionner les métropoles, décentralisation et gouvernance locale, ville intelligente.
\end{itemize}

\textbf{4. Conclusion :} la ville camerounaise reste le moteur de l'unité nationale, mais sa crise menace ce rôle ; l'enjeu des politiques d'aménagement (SND30) est de transformer l'urbanisation subie en urbanisation maîtrisée.

\textbf{Points de vigilance :} exemples chiffrés obligatoires ; éviter le catalogue de problèmes sans lien avec l'intégration ; équilibrer les parties.
\end{corrige}

\begin{corrige}[Exercice 11 — TP : Construction d'un croquis de quartier cosmopolite (Mvog-Mbi, Yaoundé)]
\textbf{Barème indicatif :} analyse du document 2 pts ; croquis 10 pts ; commentaire 8 pts.

\textbf{1. Analyse du document source (2 pts).} L'élève doit relever : l'échelle (ex. 1/5 000), l'orientation (flèche du nord), le quadrillage orthogonal de la voirie (voirie principale bitumée, voies secondaires en piste), les types d'occupation du sol (parcelles bâties, terrain vague, zone marécageuse en bas de pente) et les équipements (marché, école, mosquée, église, centre de santé).

\textbf{2. Construction du croquis (10 pts) — grille de notation :}
\begin{itemize}
    \item Cadre, titre centré (« Le quartier cosmopolite de Mvog-Mbi (Yaoundé) »), flèche du nord, échelle graphique : 2 pts.
    \item Légende hiérarchisée tracée \emph{avant} le croquis, avec les 5 rubriques exigées (habitat, activités, équipements, voirie et réseaux, contraintes/risques) : 3 pts.
    \item Sémiologie correcte : surfaces pour l'habitat (couleur claire / hachures rouges), linéaires pour la voirie (trait continu épais = bitumé, pointillés = piste), ponctuels pour les équipements (symboles école, mosquée, église, CMA, borne-fontaine) : 3 pts.
    \item Cohérence spatiale : habitat spontané près de la zone marécageuse et en périphérie, marché au carrefour principal, équipements accessibles par la voirie bitumée : 2 pts.
\end{itemize}

\textbf{3. Commentaire géographique (8 pts) — éléments attendus :}
\begin{itemize}
    \item \emph{Morphologie (2 pts) :} centre dense et commerçant autour du marché, habitat planifié le long de la voirie bitumée, habitat spontané en périphérie et sur les bas-fonds ; densité décroissante du centre vers la lisière.
    \item \emph{Cosmopolitisme (2 pts) :} juxtaposition d'une mosquée, d'églises et d'un temple ; marché intercommunautaire ; écoles accueillant des enfants de toutes origines : le quartier est un espace de brassage culturel et religieux.
    \item \emph{Dysfonctionnements (2 pts) :} voirie secondaire non bitumée, caniveaux absents dans les bas-fonds, habitat précaire en zone inondable (hachures bleues), décharge sauvage : insalubrité et risque d'inondation.
    \item \emph{Pistes d'aménagement (2 pts) :} restructuration foncière et viabilisation des quartiers spontanés, curage des caniveaux et recul de l'habitat hors de la zone inondable, création des équipements manquants (centre de santé secondaire, espaces verts).
\end{itemize}

\textbf{Points de vigilance :} la légende doit être \emph{organisée avant} le tracé et chaque élément du croquis doit figurer dans la légende (et réciproquement) ; interdiction d'utiliser des couleurs non conventionnelles sans les définir ; le commentaire doit s'appuyer sur des éléments \emph{visibles} du croquis, pas sur des généralités hors sujet.
\end{corrige}

\newpage

\coursec{Fiche bilan}

\begin{popfigure}
\begin{tikzpicture}[mindmap, grow cyclic, every node/.style={concept}, root concept/.append style={concept color=popblue, font=\bfseries\small, minimum size=2.6cm}, level 1/.append style={level distance=3.6cm, sibling angle=60, font=\scriptsize}, level 2/.append style={level distance=2.2cm, font=\tiny, minimum size=1.1cm}]
\node[root concept]{Les villes du Cameroun}
  child[concept color=poporange]{ node{Généralités}
    child{ node{taux d'urbanisation $\approx$ 58\,\%}}
    child{ node{réseau urbain dominé par Douala et Yaoundé}}
  }
  child[concept color=popgreen]{ node{Espace d'intégration}
    child{ node{brassage des populations}}
    child{ node{quartiers cosmopolites (New Bell, Briqueterie)}}
  }
  child[concept color=poppurple]{ node{Impulsion économique}
    child{ node{industries, banques, port de Douala}}
    child{ node{polarisation des services}}
  }
  child[concept color=poppink]{ node{Problèmes urbains}
    child{ node{insalubrité, anarchie, incivisme}}
    child{ node{logement, transport, eau, électricité}}
  }
  child[concept color=popgold]{ node{Solutions}
    child{ node{planification, voirie, assainissement}}
    child{ node{décentralisation, villes secondaires}}
  }
  child[concept color=popmuted]{ node{Méthodes}
    child{ node{croquis, diagrammes}}
    child{ node{lecture de cartes et tableaux}}
  };
\end{tikzpicture}
\legende{Carte mentale de la leçon : les six grandes entrées à mobiliser le jour du Bac.}
\end{popfigure}

\begin{bilan}
\textbf{Définitions clés}
\begin{itemize}
  \item \cle{Ville} : espace concentré de population (au Cameroun, unité d'au moins \num{5000} habitants exerçant des activités majoritairement non agricoles), pôle de fonctions administratives, économiques et culturelles.
  \item \cle{Urbanisation} : processus de croissance des villes (démographique et spatiale) ; le taux d'urbanisation du Cameroun dépasse 55\,\% (ordre de grandeur : environ 58\,\% vers 2020).
  \item \cle{Quartier cosmopolite} : quartier où cohabitent des populations de régions, d'ethnies, de religions et parfois de nationalités différentes (ex. New Bell à Douala, Briqueterie à Yaoundé).
  \item \cle{Voirie urbaine} : ensemble des voies de circulation d'une ville (routes, avenues, rues, pistes) et leurs aménagements.
  \item \cle{Polarisation} : capacité d'une ville à concentrer les services (santé, éducation, banques, administration) et à attirer les populations de son aire d'influence.
\end{itemize}

\textbf{Repères chiffrés à connaître}
\begin{center}
\begin{tabularx}{\linewidth}{|l|X|}\hline
\rowcolor{popblueL} \textbf{Indicateur} & \textbf{Ordre de grandeur} \\ \hline
Taux d'urbanisation du Cameroun & environ 58\,\% (2020), contre moins de 30\,\% en 1976 \\ \hline
Douala (capitale économique) & plus de 3 millions d'habitants \\ \hline
Yaoundé (capitale politique) & plus de 3 millions d'habitants \\ \hline
Villes secondaires majeures & Garoua, Bamenda, Bafoussam, Maroua, Ngaoundéré, Buea \\ \hline
\end{tabularx}
\end{center}

\textbf{Idées-forces}
\begin{itemize}
  \item La ville camerounaise est un \cle{espace d'intégration nationale} : brassage des populations, des cultures et des religions ; langues véhiculaires (français, anglais, pidgin, fulfuldé, ewondo) ; mariages mixtes ; écoles, stades et marchés comme lieux de rencontre.
  \item La ville est un \cle{centre d'impulsion économique} : industries (Douala-Bonabéri, zone portuaire), sièges d'entreprises et de banques, marchés (marché Mboppi, marché Mokolo), et polarisation des services supérieurs (universités, hôpitaux de référence, ministères).
  \item Les \cle{problèmes} : incivisme, occupation anarchique du sol, insalubrité (ordures, eaux usées), insécurité, déficit de logements (quartiers spontanés), embouteillages, coupures d'eau et d'électricité.
  \item Les \cle{solutions} : plans d'urbanisme, lotissements, extension et entretien de la voirie, assainissement (HYSACAM), décentralisation et promotion des villes secondaires pour rééquilibrer l'armature urbaine.
\end{itemize}

\textbf{Méthodes express}
\begin{itemize}
  \item \emph{Croquis urbain} : titre, légende numérotée, orientation (nord), échelle indicative ; placer d'abord les grandes masses (centre, quartiers, axes, industries) avant les détails.
  \item \emph{Commentaire de document} : observer (nature, source, date) $\rightarrow$ décrire les faits $\rightarrow$ expliquer (causes, conséquences) $\rightarrow$ conclure par une ouverture.
  \item \emph{Diagramme} : vérifier que les pourcentages totalisent 100\,\% ; toujours indiquer l'unité et la source.
\end{itemize}
\end{bilan}

\begin{aretenir}[Checklist avant le Bac]
\begin{itemize}
  \item[$\square$] Je sais définir ville, urbanisation, quartier cosmopolite, voirie urbaine et polarisation.
  \item[$\square$] Je connais le taux d'urbanisation approximatif du Cameroun et son évolution depuis 1976.
  \item[$\square$] Je sais citer les deux métropoles (Douala, Yaoundé) et au moins quatre villes secondaires avec leur région.
  \item[$\square$] Je peux expliquer en quoi la ville est un espace d'intégration nationale avec deux exemples de quartiers cosmopolites.
  \item[$\square$] Je sais montrer le rôle d'impulsion économique des villes (port de Douala, industries, marchés, banques).
  \item[$\square$] Je connais au moins cinq problèmes des villes camerounaises et je sais les illustrer localement.
  \item[$\square$] Je peux proposer des solutions concrètes : planification, voirie, assainissement, décentralisation.
  \item[$\square$] Je sais construire un croquis complet : titre, légende, nord, échelle.
  \item[$\square$] Je sais lire un tableau statistique et calculer un taux de variation ou une part en pourcentage.
  \item[$\square$] Je rédige un commentaire de document en trois temps : décrire, expliquer, conclure.
\end{itemize}
\end{aretenir}

\findecours

\end{document}
